% =============================================================================
% APPENDIX C: The Complete FEPSDE Welfare Matrix
% =============================================================================
% Module: BCM2_04_INU_6040 / KNU_6041 / IDN_6042 (BEATRIX Architecture)
% Version: 2.2 (January 2026)
%
% Complete reference for:
% - Meta-Axiom U-M0 (universelle Nutzenstruktur)
% - 144-component utility structure
% - INU/KNU/IDN axiom systems
% - Context-welfare interaction (Γ-matrix)
%
% =============================================================================
% CRITICAL CROSS-REFERENCE: Appendix C ↔ Appendix AAA
% =============================================================================
%
% Appendix C defines the HORIZONTAL structure: utility DIMENSIONS (FEPSDE)
% Appendix AAA defines the VERTICAL structure: aggregation LEVELS (L0...L∞)
%
% Together they form the complete utility space: U^L_d
%   - Rows = Levels (from AAA)
%   - Columns = Dimensions (from C)
%
% ALL utilities in this appendix implicitly carry a level index:
%   U_d should be read as U^L_d (utility in dimension d at level L)
%
% See Section C.2a for the formal integration.
%
% =============================================================================
% CROSS-REFERENCE STRUCTURE: Appendix C ↔ Chapter 10
% =============================================================================
%
% ┌────────────────────────────────────────────────────────────────────────────┐
% │ App C Section │ Content                         │ Chapter 10 Section       │
% ├────────────────────────────────────────────────────────────────────────────┤
% │ C.1           │ 144-Component Structure         │ 10.3                     │
% │ C.2           │ Meta-Axiom U-M0                 │ 10.1                     │
% │ C.2a          │ Integration with AAA (Levels)   │ NEW                      │
% │ C.3           │ INU (59 Axiome)                 │ 10.1, 10.3               │
% │ C.4           │ KNU (Axiome)                    │ 10.1, 10.3               │
% │ C.5           │ IDN (Axiome)                    │ 10.1, 10.3               │
% │ C.6           │ Loss Aversion by Dimension      │ 10.4                     │
% │ C.7           │ Γ-Matrix (Context-Welfare)      │ 10.7                     │
% │ C.8           │ Context-Dependent Parameters    │ 10.7                     │
% │ C.9           │ Inter-Category Complementarity  │ 10.6                     │
% │ C.10          │ Counting Verification           │ 10.3                     │
% │ C.11          │ Worked Example                  │ 10.9                     │
% │ C.12          │ Validation                      │ 10.7                     │
% └────────────────────────────────────────────────────────────────────────────┘
%
% =============================================================================

\chapter{The Complete FEPSDE Welfare Matrix}
\label{app:fepsde-matrix}

\begin{abstract}
This appendix provides the complete reference documentation for the EBF 
utility framework: the Meta-Axiom U-M0 defining the universal structure, 
the 144-component decomposition, the axiom systems for INU/KNU/IDN, and 
the context-welfare interaction via the $\Gamma$-matrix.
\end{abstract}

%==============================================================================
% QUICK REFERENCE BOX (Required)
%==============================================================================

\begin{tcolorbox}[colback=yellow!10!white, colframe=yellow!75!black, 
    title=\textbf{Quick Reference}]

\textbf{CORE Question:} WHAT constitutes utility?

\textbf{Answer:} Utility decomposes into $\mathbf{144}$ distinct components 
across 4 orthogonal dimensions.

\textbf{The 144-Component Formula:}
\[
\text{Total} = \underbrace{3}_{\text{Category}} \times \underbrace{2}_{\text{Valence}} \times \underbrace{6}_{\text{FEPSDE}} \times \underbrace{4}_{\text{Time}} = 144
\]

\textbf{Key Structures:}
\begin{center}
\small
\begin{tabular}{@{}llc@{}}
\toprule
\textbf{Structure} & \textbf{Content} & \textbf{Axioms} \\
\midrule
Meta-Axiom U-M0 & Universal utility structure & 1 \\
INU & Individual utility ($U^0_d$) & 59 \\
KNU & Collective utility ($U^1_d$) & 49 \\
IDN & Identity utility ($U^{ID}_d$) & 35 \\
\midrule
& \textbf{Total} & \textbf{144} \\
\bottomrule
\end{tabular}
\end{center}

\textbf{FEPSDE Dimensions:}
\begin{itemize}[nosep]
    \item \textbf{F}inancial --- Income, wealth, economic security
    \item \textbf{E}motional --- Affect, mood, psychological wellbeing
    \item \textbf{P}hysical --- Health, fitness, bodily function
    \item \textbf{S}ocial --- Relationships, belonging, status
    \item \textbf{D}igital --- Online presence, data, connectivity
    \item \textbf{Eco}logical --- Environment, sustainability, nature
\end{itemize}

\textbf{Jump to:}
\begin{itemize}[nosep]
    \item 144-Component Structure $\rightarrow$ \S\ref{sec:app-c-structure}
    \item Meta-Axiom U-M0 $\rightarrow$ \S\ref{sec:app-c-meta-axiom}
    \item INU Axioms (59) $\rightarrow$ \S\ref{sec:app-c-inu}
    \item KNU Axioms (49) $\rightarrow$ \S\ref{sec:app-c-knu}
    \item IDN Axioms (35) $\rightarrow$ \S\ref{sec:app-c-idn}
    \item $\Gamma$-Matrix $\rightarrow$ \S\ref{sec:app-c-gamma}
    \item Worked Example $\rightarrow$ \S\ref{sec:app-c-example}
\end{itemize}

\end{tcolorbox}

%==============================================================================
% FORMAL COMPLIANCE BLOCK
%==============================================================================

% --- Terminology Reference ---
\begin{tcolorbox}[colback=gray!5!white, colframe=gray!50!black]
\textbf{Terminology:} See Appendix G (Glossary) for standard definitions.
Key terms in this appendix: FEPSDE dimensions ($d \in \{F,E,P,S,D,Eco\}$),
utility category (INU/KNU/IDN), valence (Gain/Pain), time horizon ($t \in \{0,1,2,3\}$),
loss aversion ($\lambda_d$), $\Gamma$-matrix (dimension interactions).
\end{tcolorbox}

% --- Chapter Linkage Box ---
\begin{tcolorbox}[colback=green!5!white, colframe=green!75!black,
    title=\textbf{Chapter References}]

\textbf{Primary Chapter:} Ch. 10 (Die Nutzenarchitektur)
\begin{itemize}[nosep]
    \item Ch. 10.1: Three utility categories (INU/KNU/IDN)
    \item Ch. 10.2: The six FEPSDE dimensions
    \item Ch. 10.3: The 144-component structure
\end{itemize}

\textbf{Secondary Chapters:}
\begin{itemize}[nosep]
    \item Ch. 6 (Reference Structure): Reference points by dimension
    \item Ch. 10.4 (Loss Aversion): $\lambda_d$ by dimension
    \item Ch. 10.6 (Inter-Category): INU$\times$KNU$\times$IDN complementarity
    \item Ch. 10.7 (Context Modulation): $\Gamma(\Psi)$ context effects
\end{itemize}

\textbf{Cross-References to CORE:}
\begin{itemize}[nosep]
    \item Appendix AAA (CORE-WHO): Levels $L$ on which dimensions operate
    \item Appendix B (CORE-HOW): $\gamma_{dd'}$ dimension complementarity
    \item Appendix V (CORE-WHEN): $\Psi$ modulates dimension weights
    \item Appendix BBB (CORE-WHERE): Dimension-specific parameters
    \item \textbf{XVIII (LIT-FEHR-METHOD):} Exclusion Principle --- $\gamma = 0$ as null hypothesis
\end{itemize}

\end{tcolorbox}

% --- Epistemic Status Tags ---
\begin{tcolorbox}[colback=blue!5!white, colframe=blue!75!black,
    title=\textbf{Epistemic Status: Key Parameters}]

\renewcommand{\arraystretch}{1.3}
\begin{tabular}{@{}llcc@{}}
\toprule
\textbf{Parameter} & \textbf{Description} & \textbf{Value/Range} & \textbf{Tag} \\
\midrule
$d$ & Number of FEPSDE dimensions & 6 & [THR] \\
144 & Total components & $3 \times 2 \times 6 \times 4$ & [THR] \\
$\lambda_F$ & Financial loss aversion & 2.0--2.5 & [EMP] \\
$\lambda_E$ & Emotional loss aversion & 1.5--2.0 & [EMP] \\
$\lambda_P$ & Physical loss aversion & 2.5--3.0 & [EMP] \\
$\lambda_S$ & Social loss aversion & 1.8--2.2 & [EMP] \\
$\lambda_D$ & Digital loss aversion & 1.2--1.5 & [LLM] \\
$\lambda_{Eco}$ & Ecological loss aversion & 1.0--1.3 & [LLM] \\
\addlinespace
\multicolumn{4}{@{}l}{\textit{Axiom counts:}} \\
INU axioms & Individual utility & 59 & [THR] \\
KNU axioms & Collective utility & 49 & [THR] \\
IDN axioms & Identity utility & 35 & [THR] \\
\bottomrule
\end{tabular}

\vspace{0.5em}
\textbf{Tag Legend:} [EMP] = Empirically validated; [THR] = Theoretically derived;
[LLM] = LLM Monte Carlo estimated

\end{tcolorbox}

% --- Bibliography Reference ---
\noindent\textbf{Full citations:} See \texttt{bibliography/bcm\_master.bib} \\
\textbf{Key sources:} Kahneman \& Tversky (1979); Fehr \& Schmidt (1999); Akerlof \& Kranton (2000)

\vspace{1em}

%%%%%%%%%%%%%%%%%%%%%%%%%%%%%%%%%%%%%%%%%%%%%%%%%%%%%%%%%%%%%%%%%%%%%%%%%%%%%%%
%%%%%%%%%%%%%%%%%%%%%%%%%%%%%%%%%%%%%%%%%%%%%%%%%%%%%%%%%%%%%%%%%%%%%%%%%%%%%%%
\section{The 144-Component Structure}
\label{sec:app-c-structure}
%%%%%%%%%%%%%%%%%%%%%%%%%%%%%%%%%%%%%%%%%%%%%%%%%%%%%%%%%%%%%%%%%%%%%%%%%%%%%%%
%%%%%%%%%%%%%%%%%%%%%%%%%%%%%%%%%%%%%%%%%%%%%%%%%%%%%%%%%%%%%%%%%%%%%%%%%%%%%%%

\begin{cross-reference}
\textbf{Chapter Reference:} See Chapter~\ref{sec:welfare-fepsde}, 
Section~\ref{subsec:10-3-structure} for conceptual introduction.
\end{cross-reference}

\subsection{Dimensional Decomposition}

The complete utility space consists of 144 distinct components, arising from 
four orthogonal dimensions:

\begin{table}[htbp]
\centering
\caption{The Four Dimensions of Utility Space}
\label{tab:app-c-dimensions}
\renewcommand{\arraystretch}{1.4}
\begin{tabular}{@{}llll@{}}
\toprule
\textbf{Dimension} & \textbf{Symbol} & \textbf{Values} & \textbf{Count} \\
\midrule
Utility Category & $i$ & INU, KNU, IDN & 3 \\
Valence & $v$ & Gain ($G$), Pain ($P$) & 2 \\
FEPSDE Dimension & $d$ & F, E, P, S, D, Eco & 6 \\
Time Horizon & $t$ & Instant, Short, Medium, Long & 4 \\
\midrule
\multicolumn{3}{l}{\textbf{Total Components}} & $\mathbf{3 \times 2 \times 6 \times 4 = 144}$ \\
\bottomrule
\end{tabular}
\end{table}

\subsection{Component Notation}

Each of the 144 components is uniquely identified:

\begin{equation}
X^i_{d,t} \quad \text{where} \quad 
\begin{cases}
X \in \{G, P\} & \text{(Gain or Pain)} \\
i \in \{\text{INU}, \text{KNU}, \text{IDN}\} & \text{(Utility Category)} \\
d \in \{F, E, P, S, D, \text{Eco}\} & \text{(FEPSDE Dimension)} \\
t \in \{0, 1, 2, 3\} & \text{(Time Horizon)}
\end{cases}
\end{equation}

%%%%%%%%%%%%%%%%%%%%%%%%%%%%%%%%%%%%%%%%%%%%%%%%%%%%%%%%%%%%%%%%%%%%%%%%%%%%%%%
%%%%%%%%%%%%%%%%%%%%%%%%%%%%%%%%%%%%%%%%%%%%%%%%%%%%%%%%%%%%%%%%%%%%%%%%%%%%%%%
\section{Core Theory: Meta-Axiom U-M0 (Universal Utility Structure)}
\label{sec:app-c-meta-axiom}
%%%%%%%%%%%%%%%%%%%%%%%%%%%%%%%%%%%%%%%%%%%%%%%%%%%%%%%%%%%%%%%%%%%%%%%%%%%%%%%
%%%%%%%%%%%%%%%%%%%%%%%%%%%%%%%%%%%%%%%%%%%%%%%%%%%%%%%%%%%%%%%%%%%%%%%%%%%%%%%

\begin{cross-reference}
\textbf{Chapter Reference:} See Chapter~\ref{sec:welfare-fepsde}, 
Section~\ref{subsec:10-1-categories} for conceptual foundations.
\end{cross-reference}

\subsection{The Meta-Axiom}

\begin{tcolorbox}[colback=red!5!white, colframe=red!75!black, 
    title=\textbf{META-AXIOM U-M0: Universal Utility Structure}]

\textbf{Statement:} All three utility categories (INU, KNU, IDN) follow identical 
mathematical structure. The difference lies not in the formula, but exclusively 
in the calculation of parameters.

\begin{equation}
\boxed{
U^{i} = \sum_{d \in \mathcal{D}} \sum_{t=0}^{3} \delta^{i}_{d}(t) \cdot 
\left[ G^{i}_{d,t} - \lambda^{i}_{d} \cdot P^{i}_{d,t} \right]
}
\label{eq:meta-axiom}
\end{equation}

\textbf{where:}
\begin{itemize}
\item $i \in \{\text{INU}, \text{KNU}, \text{IDN}\}$: utility category
\item $d \in \mathcal{D} = \{F, E, P, S, D, Eco\}$: FEPSDE dimension (6)
\item $t \in \{0, 1, 2, 3\}$: time horizon (4)
\item $G^{i}_{d,t}$: gain in category $i$, dimension $d$, time $t$
\item $P^{i}_{d,t}$: pain in category $i$, dimension $d$, time $t$
\item $\lambda^{i}_{d}$: loss aversion in category $i$, dimension $d$
\item $\delta^{i}_{d}(t)$: discount factor in category $i$, dimension $d$, time $t$
\end{itemize}

\end{tcolorbox}

\subsection{Structural Identity}

\begin{equation}
\text{Components per category} = 6 \times 4 \times 2 = 48
\end{equation}
\begin{equation}
\text{Total components} = 3 \times 48 = 144
\end{equation}

\subsection{Axiom Hierarchy}

\begin{table}[htbp]
\centering
\caption{Hierarchical Structure of Utility Axioms}
\label{tab:axiom-hierarchy}
\renewcommand{\arraystretch}{1.4}
\begin{tabular}{@{}clcl@{}}
\toprule
\textbf{Level} & \textbf{Designation} & \textbf{Count} & \textbf{Function} \\
\midrule
\textbf{0} & \textbf{Meta-Axiom (U-M0)} & \textbf{1} & \textbf{Master structure for all $i$} \\
\midrule
1 & INU-Axioms (INU-A0...A57) & 59 & Parameter calculation for $U^{INU}$ \\
1 & KNU-Axioms (KNU-A1...A65) & 65 & Parameter calculation for $U^{KNU}$ \\
1 & IDN-Axioms (IDN-A1...A94) & 95 & Parameter calculation for $U^{IDN}$ \\
\midrule
& \textbf{Total Axioms} & \textbf{220} & \\
\bottomrule
\end{tabular}
\end{table}

\subsection{The Three Utility Categories: Overview}

\begin{table}[htbp]
\centering
\caption{The Three Utility Categories}
\label{tab:app-c-three-categories}
\renewcommand{\arraystretch}{1.5}
\begin{tabular}{@{}clll@{}}
\toprule
\textbf{Symbol} & \textbf{Name} & \textbf{Core Question} & \textbf{Reference Point Source} \\
\midrule
$U^{INU}$ & Individual Utility & ``What do I want for MYSELF?'' & Status quo, expectations, peers \\
$U^{KNU}$ & Collective Utility & ``What do I want for US?'' & Group norm, fairness \\
$U^{IDN}$ & Identity Utility & ``Who AM I?'' & Self-concept, identity standard \\
\bottomrule
\end{tabular}
\end{table}


%%%%%%%%%%%%%%%%%%%%%%%%%%%%%%%%%%%%%%%%%%%%%%%%%%%%%%%%%%%%%%%%%%%%%%%%%%%%%%%
%%%%%%%%%%%%%%%%%%%%%%%%%%%%%%%%%%%%%%%%%%%%%%%%%%%%%%%%%%%%%%%%%%%%%%%%%%%%%%%
\section{Integration with Aggregation Levels (Appendix AAA)}
\label{sec:app-c-aaa-integration}
%%%%%%%%%%%%%%%%%%%%%%%%%%%%%%%%%%%%%%%%%%%%%%%%%%%%%%%%%%%%%%%%%%%%%%%%%%%%%%%
%%%%%%%%%%%%%%%%%%%%%%%%%%%%%%%%%%%%%%%%%%%%%%%%%%%%%%%%%%%%%%%%%%%%%%%%%%%%%%%

\begin{tcolorbox}[colback=blue!5!white, colframe=blue!75!black, 
    title=Critical Integration: Appendix C $\times$ Appendix AAA]
\textbf{Appendix C and Appendix AAA are orthogonal but interdependent:}
\begin{itemize}[nosep]
\item \textbf{Appendix C} defines the \textit{horizontal} structure: utility \textbf{dimensions} (FEPSDE)
\item \textbf{Appendix AAA} defines the \textit{vertical} structure: aggregation \textbf{levels} ($L = 0, 1, 2, ..., \infty$)
\end{itemize}

Together they form the complete utility space: $U^L_d$ --- utility at level $L$ in dimension $d$.
\end{tcolorbox}

\subsection{Why Level Structure Matters for FEPSDE}

Every utility component defined in this appendix implicitly carries a \textbf{level index}. 
When we write $U_F$ (Financial utility), we actually mean $U^L_F$ --- Financial utility 
\textit{at some level $L$}.

\subsubsection{Theoretical Justification for Level-Dimension Integration}

\begin{tcolorbox}[colback=gray!5!white, colframe=gray!75!black, 
    title=Why AAA $\times$ C Is Theoretically Necessary]

The integration of Appendix AAA (Levels) with Appendix C (FEPSDE Dimensions) is not 
merely a notational convenience---it is \textbf{theoretically required} for five reasons:

\paragraph{Reason 1: Avoiding the Aggregation Fallacy.}

Without level specification, statements like ``Financial utility increased'' are 
ambiguous and potentially fallacious:
\begin{itemize}[nosep]
\item GDP (national $U^6_F$) can rise while median individual income ($U^0_F$) falls
\item Family welfare ($U^1_E$) can improve while individual members suffer ($U^0_E$)
\item Organizational profit ($U^3_F$) does not imply employee financial welfare ($U^0_F$)
\end{itemize}

The ecological fallacy (Robinson, 1950) and Simpson's paradox demonstrate that 
aggregate patterns need not reflect individual patterns. Level indices make this explicit.

\paragraph{Reason 2: Explaining Cross-Level Trade-offs.}

Many real-world dilemmas involve trade-offs \textit{between levels within the same dimension}:
\begin{itemize}[nosep]
\item \textbf{Financial}: Individual tax burden ($U^0_F \downarrow$) vs. national investment ($U^6_F \uparrow$)
\item \textbf{Ecological}: Individual consumption ($U^0_{Eco} \uparrow$) vs. global sustainability ($U^{\infty}_{Eco} \downarrow$)
\item \textbf{Social}: Family loyalty ($U^1_S$) vs. organizational commitment ($U^3_S$)
\end{itemize}

Without $U^L_d$ notation, these trade-offs are invisible.

\paragraph{Reason 3: Level-Specific Parameter Variation.}

Behavioral parameters vary systematically by level:
\begin{itemize}[nosep]
\item \textbf{Loss aversion} ($\lambda$): Individuals show $\lambda \approx 2.25$ (Kahneman \& Tversky, 1979); 
      organizations show lower loss aversion; nations show extreme loss aversion for territory
\item \textbf{Discount rates} ($\delta$): Individuals discount hyperbolically; organizations use 
      market rates; nations should (but often don't) use near-zero rates for long-term policy
\item \textbf{Dimension weights} ($\omega_d$): Social dimension may dominate at tribal level ($L=2$) 
      but financial dimension dominates at organizational level ($L=3$)
\end{itemize}

These are not anomalies---they are \textbf{systematic} level effects requiring explicit modeling.

\paragraph{Reason 4: The Endogeneity of Levels.}

Appendix AAA (Axiom AL-2) establishes that levels are \textit{endogenous}---they exist 
because of shared identity (IDN) or collective utility (KNU). This means:
\begin{itemize}[nosep]
\item The \textit{same} FEPSDE dimensions apply at each level
\item But the \textit{existence} of the level depends on IDN/KNU
\item And the \textit{aggregation} from lower levels follows specific rules
\end{itemize}

FEPSDE (Appendix C) provides the dimensions; AAA provides the levels and aggregation rules.

\paragraph{Reason 5: Super-Additive Cooperation Across Levels.}

Recent research (Efferson, Bernhard, Fischbacher \& Fehr, 2024) demonstrates that 
cooperation mechanisms are \textbf{super-additive} across levels: neither individual-level 
repeated interactions nor group-level competition alone sustains cooperation, but their 
\textit{combination} does.

This directly implies that welfare analysis must consider multiple levels simultaneously---exactly 
what the $U^L_d$ framework provides.

\end{tcolorbox}

\paragraph{The Problem with Level-Free Notation.}

Consider ``Financial utility'' ($U_F$). Without specifying the level, this is ambiguous:

\begin{itemize}[nosep]
\item $U^0_F$: Individual income, wealth, consumption
\item $U^1_F$: Family/household financial resources
\item $U^2_F$: Tribal/community economic resources
\item $U^3_F$: Organizational revenue, profit
\item $U^6_F$: National GDP, sovereign wealth
\item $U^{\infty}_F$: Global economic welfare
\end{itemize}

These are \textit{different quantities} with different measurement scales, different 
determinants, and different policy implications.

\paragraph{The Complete Notation.}

The fully specified utility component is:
\begin{equation}
U^L_{i,d,v,t}
\end{equation}
where:
\begin{itemize}[nosep]
\item $L$: Level (from Appendix AAA)
\item $i$: Category (INU, KNU, IDN)
\item $d$: Dimension (F, E, P, S, D, Eco)
\item $v$: Valence (Gain, Pain)
\item $t$: Time horizon (Instant, Short, Medium, Long)
\end{itemize}

\subsection{The Level-Dimension Matrix}

Combining levels (AAA) with dimensions (C) yields a matrix structure:

\begin{equation}
\mathbf{U} = \begin{pmatrix}
U^0_F & U^0_E & U^0_P & U^0_S & U^0_D & U^0_{Eco} \\
U^1_F & U^1_E & U^1_P & U^1_S & U^1_D & U^1_{Eco} \\
U^2_F & U^2_E & U^2_P & U^2_S & U^2_D & U^2_{Eco} \\
\vdots & \vdots & \vdots & \vdots & \vdots & \vdots \\
U^{\infty}_F & U^{\infty}_E & U^{\infty}_P & U^{\infty}_S & U^{\infty}_D & U^{\infty}_{Eco}
\end{pmatrix}
\end{equation}

\begin{table}[htbp]
\centering
\caption{FEPSDE Dimensions at Different Levels}
\label{tab:app-c-level-dimension}
\renewcommand{\arraystretch}{1.3}
\small
\begin{tabular}{@{}clll@{}}
\toprule
\textbf{Level} & \textbf{Example Entity} & \textbf{Financial ($U^L_F$)} & \textbf{Social ($U^L_S$)} \\
\midrule
$L = 0$ & Individual & Personal income & Personal relationships \\
$L = 1$ & Family & Household income & Family bonds \\
$L = 2$ & Tribe/Community & Community resources & Social capital \\
$L = 3$ & Organization & Revenue, profit & Organizational culture \\
$L = 4$ & Network & Network value & Network effects \\
$L = 6$ & Nation & GDP & National cohesion \\
$L = \infty$ & All sentient beings & Global wealth & Universal solidarity \\
\bottomrule
\end{tabular}
\end{table}

\subsection{Aggregation Rules: From AAA to C}

Appendix AAA specifies how utilities aggregate \textit{across levels}. This applies 
to \textit{each dimension} defined in Appendix C:

\begin{tcolorbox}[colback=green!5!white, colframe=green!75!black, 
    title=Dimension-Specific Aggregation (from AAA Axiom AL-3)]
For each FEPSDE dimension $d \in \{F, E, P, S, D, Eco\}$:

\begin{equation}
U^L_d = \sum_{i \in M^L} \omega^L_{i,d} \cdot U^{L-1}_{i,d} + \sum_{i < j} \gamma^L_{ij,d} \cdot U^{L-1}_{i,d} \cdot U^{L-1}_{j,d}
\end{equation}

where:
\begin{itemize}[nosep]
\item $M^L$: Membership set of level $L$ (defined in AAA)
\item $\omega^L_{i,d}$: Weight of member $i$ in dimension $d$ at level $L$
\item $\gamma^L_{ij,d}$: Complementarity between members $i$ and $j$ in dimension $d$
\end{itemize}
\end{tcolorbox}

\paragraph{Key Insight: Dimension-Specific Complementarities.}

The complementarity parameter $\gamma^L_{ij,d}$ can vary by dimension:

\begin{itemize}[nosep]
\item \textbf{Financial}: Family members may have $\gamma^1_{ij,F} > 0$ (shared resources benefit all)
\item \textbf{Social}: Tribal members have strong $\gamma^2_{ij,S} > 0$ (social capital is superadditive)
\item \textbf{Ecological}: Nations may have $\gamma^6_{ij,Eco} < 0$ (environmental competition)
\end{itemize}

\subsection{Level-Specific Parameters}

Many parameters defined in Appendix C should be understood as level-specific:

\begin{table}[htbp]
\centering
\caption{Level-Specific Parameters}
\label{tab:app-c-level-params}
\renewcommand{\arraystretch}{1.3}
\begin{tabular}{@{}lll@{}}
\toprule
\textbf{Parameter} & \textbf{Level-Free (C)} & \textbf{Level-Specific (AAA $\times$ C)} \\
\midrule
Loss aversion & $\lambda_d$ & $\lambda^L_d$ \\
Discount factor & $\delta_d$ & $\delta^L_d$ \\
Dimension weight & $\omega_d$ & $\omega^L_d$ \\
Complementarity & $\gamma_{dd'}$ & $\gamma^L_{dd'}$ \\
\bottomrule
\end{tabular}
\end{table}

\paragraph{Example: Level-Specific Loss Aversion.}

Loss aversion in the Financial dimension may differ by level:

\begin{itemize}[nosep]
\item $\lambda^0_F \approx 2.25$: Individual loss aversion (Kahneman \& Tversky)
\item $\lambda^3_F \approx 1.5$: Organizational loss aversion (less emotional)
\item $\lambda^6_F \approx 3.0$: National loss aversion (territory, sovereignty)
\end{itemize}

\subsection{KNU and IDN as Level-Defining}

\begin{tcolorbox}[colback=yellow!5!white, colframe=yellow!75!black, 
    title=Critical Connection: KNU/IDN Define Levels (AAA Axiom AL-2)]
Appendix C defines KNU (Collective Utility) and IDN (Identity Utility) as 
utility categories. But Appendix AAA reveals a deeper role:

\begin{equation}
\boxed{\text{Level } L \text{ exists} \iff \exists \, IDN^L \lor \exists \, KNU^L}
\end{equation}

\textbf{Implication}: KNU and IDN do not merely \textit{contribute to} utility at 
pre-existing levels. They \textit{define which levels exist}.

\begin{itemize}[nosep]
\item A ``family'' (Level 1) exists because family members share identity ($IDN^1$) 
      or collective goals ($KNU^1$)
\item A ``tribe'' (Level 2) exists because tribe members share identity ($IDN^2$) 
      or collective welfare ($KNU^2$)
\item If $IDN^L = 0$ and $KNU^L = 0$, then Level $L$ does not exist for that group
\end{itemize}
\end{tcolorbox}

\subsection{Complete Structure: The $U^L_d$ Framework}

\begin{tcolorbox}[colback=blue!5!white, colframe=blue!75!black, 
    title=The Complete EBF Utility Structure]

Combining Appendix C (dimensions) with Appendix AAA (levels):

\begin{equation}
\boxed{
U^{total} = \sum_{L=0}^{\infty} \alpha^L(\Psi) \cdot \left[ \sum_{d \in FEPSDE} \omega^L_d \cdot U^L_d + \sum_{d < d'} \gamma^L_{dd'} \cdot U^L_d \cdot U^L_{d'} \right]
}
\end{equation}

where:
\begin{itemize}[nosep]
\item $\alpha^L(\Psi)$: Level salience given context $\Psi$ (from AAA, Axiom AL-6)
\item $\omega^L_d$: Weight of dimension $d$ at level $L$
\item $\gamma^L_{dd'}$: Complementarity between dimensions $d$ and $d'$ at level $L$
\item $U^L_d$: Utility in dimension $d$ at level $L$ (computed recursively per AAA)
\end{itemize}

The base case (individual level):
\begin{equation}
U^0_d = \sum_t \delta^t \cdot [G_{d,t} - \lambda_d \cdot P_{d,t}]
\end{equation}
\end{tcolorbox}

\subsection{Cross-Reference Summary}

\begin{table}[htbp]
\centering
\caption{Appendix C $\leftrightarrow$ Appendix AAA Interface}
\label{tab:app-c-aaa-interface}
\renewcommand{\arraystretch}{1.3}
\begin{tabular}{@{}p{6cm}p{6cm}@{}}
\toprule
\textbf{What C Provides to AAA} & \textbf{What AAA Provides to C} \\
\midrule
FEPSDE dimension definitions & Level structure ($L = 0, 1, ..., \infty$) \\
Dimension-specific parameters ($\lambda_d$, $\delta_d$) & Aggregation rules across levels \\
Within-dimension complementarity ($\gamma_{dd'}$) & Membership sets ($M^L$) \\
INU/KNU/IDN category structure & Level salience function $\alpha^L(\Psi)$ \\
144-component decomposition & Level-specific weights ($\omega^L_i$) \\
\bottomrule
\end{tabular}
\end{table}

\textbf{Key Principle}: When reading Appendix C, implicitly add a level index. 
When reading Appendix AAA, implicitly add dimension indices. The complete 
specification requires both: $U^L_d$.

%%%%%%%%%%%%%%%%%%%%%%%%%%%%%%%%%%%%%%%%%%%%%%%%%%%%%%%%%%%%%%%%%%%%%%%%%%%%%%%
%%%%%%%%%%%%%%%%%%%%%%%%%%%%%%%%%%%%%%%%%%%%%%%%%%%%%%%%%%%%%%%%%%%%%%%%%%%%%%%
\section{INU: Individual Utility (59 Axioms)}
\label{sec:app-c-inu}
%%%%%%%%%%%%%%%%%%%%%%%%%%%%%%%%%%%%%%%%%%%%%%%%%%%%%%%%%%%%%%%%%%%%%%%%%%%%%%%
%%%%%%%%%%%%%%%%%%%%%%%%%%%%%%%%%%%%%%%%%%%%%%%%%%%%%%%%%%%%%%%%%%%%%%%%%%%%%%%

\begin{cross-reference}
\textbf{Chapter Reference:} See Chapter~\ref{sec:welfare-fepsde}, 
Section~\ref{subsec:10-1-categories} for INU definition.

\textbf{Level Reference:} INU operates primarily at Level 0 (Individual). 
See Appendix AAA, Section~\ref{sec:aaa-bcm-connection} for level integration.
\end{cross-reference}

\subsection{INU Definition}

\begin{tcolorbox}[colback=blue!5!white, colframe=blue!75!black, 
    title=Individual Utility (INU)]

\textbf{Core Question:} ``What do I want for MYSELF?''

\textbf{Formula (from Meta-Axiom U-M0):}
\begin{equation}
U^{INU} = \sum_{d \in \mathcal{D}} \sum_{t=0}^{3} \delta^{INU}_{d}(t) \cdot 
\left[ G^{INU}_{d,t} - \lambda^{INU}_{d} \cdot P^{INU}_{d,t} \right]
\end{equation}

\textbf{Reference Point Formation:}
\begin{equation}
R^{INU}_{d,t} = \alpha_{sq} \cdot x^{status\_quo}_{d,t} + 
               \alpha_{exp} \cdot x^{expectation}_{d,t} + 
               \alpha_{soc} \cdot x^{social\_comparison}_{d,t}
\end{equation}

\textbf{Key Properties:}
\begin{itemize}
\item Self-regarding: Only considers consequences for the individual
\item Reference-dependent: Evaluated relative to personal reference points
\item Loss-averse: $\lambda^{INU}_d > 1$ for all dimensions
\end{itemize}

\end{tcolorbox}

%-----------------------------------------------------------------------
\subsection{Complete INU Axiom Table (59 Axioms)}
%-----------------------------------------------------------------------

\begin{longtable}{@{}llp{4.5cm}p{4.5cm}@{}}
\caption{Complete INU Axiom System (59 Axioms)} 
\label{tab:inu-axioms-complete} \\
\toprule
\textbf{Axiom} & \textbf{Title} & \textbf{Core Formula} & \textbf{Reference} \\
\midrule
\endfirsthead
\multicolumn{4}{c}{\tablename\ \thetable{} -- continued} \\
\toprule
\textbf{Axiom} & \textbf{Title} & \textbf{Core Formula} & \textbf{Reference} \\
\midrule
\endhead
\midrule
\multicolumn{4}{r}{Continued on next page...} \\
\endfoot
\bottomrule
\endlastfoot

\multicolumn{4}{@{}l}{\textit{\textbf{Block 1: Basis (A0--A7)}}} \\
\addlinespace
INU-A0 & Master Function & $V_{total} = V_{base} + \sum \omega [G - \lambda P] + C$ & Fehr \& Schmidt (1999); Kahneman \& Tversky (1979) \\
INU-A1 & FEPSDE Structure & $\sum_{d \in \{F,E,P,S,D,Eco\}} \omega_d U_d$ & Lancaster (1966); Thaler (1985) \\
INU-A2 & Loss Aversion & $U = G - \lambda P, \quad \lambda > 1$ & Kahneman \& Tversky (1979); Tversky \& Kahneman (1991) \\
INU-A3 & Baseline Value & $V_{base}(i)$ (value of ``doing'') & Deci \& Ryan (1985); Scitovsky (1976) \\
INU-A4 & Intrinsic Value & $V_{intr}(Task)$ & Bénabou \& Tirole (2003); Frey (1997) \\
INU-A5 & Extrinsic Value & $V_{extr}(Reward)$ & Gneezy \& Rustichini (2000); Lazear (2000) \\
INU-A6 & Complementarity & $\kappa(V_{intr}, V_{extr}) \neq 0$ & Milgrom \& Roberts (1990/1995); Bowles \& Polania-Reyes (2012) \\
INU-A7 & Reference Points & $U = f(x - R)$ & Kőszegi \& Rabin (2006); Kahneman (1992) \\
\addlinespace

\multicolumn{4}{@{}l}{\textit{\textbf{Block 2: Temporal (A8, A12, A13, A16)}}} \\
\addlinespace
INU-A8 & Time Valuation & $\sum \beta \delta^t U_t$ & Laibson (1997); O'Donoghue \& Rabin (1999) \\
INU-A12 & Dynamic Adaptation & $\dot{N} = \alpha(\text{Feedback} - N)$ & Sutton \& Barto (1998); Rescorla-Wagner (1972) \\
INU-A13 & Memory (Decay) & $M_t = \sum \rho^k U_{t-k}$ & Mullainathan (2002); Kahneman et al. (1993) \\
INU-A16 & Life Phases & $U(t) = f(\text{LifeCyclePhase})$ & Modigliani (1954); Heckman (2006) \\
\addlinespace

\multicolumn{4}{@{}l}{\textit{\textbf{Block 3: Uncertainty (A9--A11)}}} \\
\addlinespace
INU-A9 & Risk \& Uncertainty & $U \cdot (1 - \phi_{risk} - \phi_{amb})$ & Ellsberg (1961); Von Neumann-Morgenstern (1944) \\
INU-A10 & Motivated Beliefs & $B^M = B^R (1 - \text{PainAvoid})$ & Bénabou \& Tirole (2016); Epley \& Gilovich (2016) \\
INU-A11 & Blind Spots & $U_{eff} = U (1 - \phi_{Blind})$ & Gabaix (2014); Bazerman \& Chugh (2006) \\
\addlinespace

\multicolumn{4}{@{}l}{\textit{\textbf{Block 4: Interface (A14)}}} \\
\addlinespace
INU-A14 & Willingness Interface & $W(t) = f(N_{ind}(t), \theta(t), C_{context})$ & McFadden (1974); Ajzen (1991) \\
\addlinespace

\multicolumn{4}{@{}l}{\textit{\textbf{Block 5: Cognition (A15, A32, A36, A42)}}} \\
\addlinespace
INU-A15 & Salience & $Salience = f(Awareness)$ & Bordalo, Gennaioli, Shleifer (2012/2013) \\
INU-A32 & Bounded Attention & $\omega_{eff} = \omega \cdot \text{AttShare}$ & Simon (1955); Gabaix (2019) \\
INU-A36 & Bias Corrections & $U' = U (1 + \theta_{bias})$ & Tversky \& Kahneman (1974); DellaVigna (2009) \\
INU-A42 & Meta-Cognition & $U_{meta} \propto \text{Consistency}(\text{Self}, \text{Action})$ & Festinger (1957); Bénabou \& Tirole (2002) \\
\addlinespace

\multicolumn{4}{@{}l}{\textit{\textbf{Block 6: Context (A17, A18, A21)}}} \\
\addlinespace
INU-A17 & Gender Differentiation & $U_{d,gender} = U_d \cdot \gamma_{gender}$ & Croson \& Gneezy (2009); Niederle \& Vesterlund (2007) \\
INU-A18 & Cultural Context & $U_{d,cult} = f(C_{culture}, I_{inst})$ & Henrich et al. (2010); Hofstede (1980/2001) \\
INU-A21 & Intersectionality & $U_{eff} = f(G \times A \times SES \times E)$ & Crenshaw (1989); Davis (2008) \\
\addlinespace

\multicolumn{4}{@{}l}{\textit{\textbf{Block 7: Digital (A19, A24, A40, A43)}}} \\
\addlinespace
INU-A19 & Digital Utility Axis & $U_{dig} \in \{Priv, Conn, Rep\}$ & Goldfarb \& Tucker (2019); Acemoglu (2021) \\
INU-A24 & Tech Disruption & $U_{t+1} = U_t + \text{Shock}$ & Christensen (1997); Schumpeter (1942) \\
INU-A40 & Digital Amplification & $U_{dig} (1 + \alpha \cdot \text{Feedback})$ & Zuboff (2019); Sunstein (2001) \\
INU-A43 & Digital Addiction & $U + \alpha \text{Reward} - \beta \text{Cost}$ & Allcott, Gentzkow et al. (2022); Bernheim \& Rangel (2004) \\
\addlinespace

\multicolumn{4}{@{}l}{\textit{\textbf{Block 8: Behavior (A20, A33, A41)}}} \\
\addlinespace
INU-A20 & Habit Dampening & $U_{eff} = U \cdot (1-H)$ & Becker \& Murphy (1988); Pollak (1970) \\
INU-A33 & Self-Control & $\beta U_{long} + (1-\beta)U_{short}$ & Fudenberg \& Levine (2006); Thaler \& Shefrin (1981) \\
INU-A41 & Path Dependence & $U_{t+1} = U_t + \text{PastChoices}$ & David (1985); Arthur (1994) \\
\addlinespace

\multicolumn{4}{@{}l}{\textit{\textbf{Block 9: Stress (A22)}}} \\
\addlinespace
INU-A22 & Stress \& Regulation & $P_{eff} = P \cdot (1-Cope)(1-Res)$ & Lazarus \& Folkman (1984); Resilience literature \\
\addlinespace

\multicolumn{4}{@{}l}{\textit{\textbf{Block 10: Social/Moral (A23, A37, A38)}}} \\
\addlinespace
INU-A23 & Social Comparison & $U = f(x_{own} - x_{peer})$ & Veblen (1899); Fehr \& Schmidt (1999) \\
INU-A37 & Moral Preferences & $U_{moral} = \epsilon_{alt} U_{other} + \epsilon_{rec} \text{Fair}$ & Fehr \& Gächter (2000/2002); Andreoni (1990) \\
INU-A38 & Identity Investment & $U_{ID} \propto \text{FutureSelf}$ & Akerlof \& Kranton (2000/2010); Bénabou \& Tirole (2011) \\
\addlinespace

\multicolumn{4}{@{}l}{\textit{\textbf{Block 11: Side Effects (A25, A26)}}} \\
\addlinespace
INU-A25 & Collateral Damage & $P_{total} = P + \eta^- \cdot SE^-$ & Externalities literature \\
INU-A26 & Collateral Beauty & $G_{total} = G + \eta^+ \cdot SE^+$ & Serendipity literature \\
\addlinespace

\multicolumn{4}{@{}l}{\textit{\textbf{Block 12: Neuro/RT (A27--A29, A39)}}} \\
\addlinespace
INU-A27 & Reaction Time & $RT \propto 1/|\Delta U|$ & Krajbich et al. (2010); Fehr \& Rangel (2011) \\
INU-A28 & Preference Reconstruction & $\Delta U \propto -E[RT]$ & Konovalov \& Krajbich (2019) \\
INU-A29 & Relative RT & $RelRT = RT / (\text{Base} \cdot \text{Diff})$ & Chabris et al. (2009); Alós-Ferrer (2018) \\
INU-A39 & Noisy Preferences & $U_{obs} = U_{true} + \epsilon$ & McFadden (1974); Hey \& Orme (1994) \\
\addlinespace

\multicolumn{4}{@{}l}{\textit{\textbf{Block 13: Emotions (A30)}}} \\
\addlinespace
INU-A30 & Emotional Valence & $\mathcal{E} = \pi_{pride} Pride - \pi_{shame} Shame - \pi_{guilt} Guilt$ & Akerlof \& Kranton; Bowles \& Gintis \\
\addlinespace

\multicolumn{4}{@{}l}{\textit{\textbf{Block 14: Systemic (A31, A34)}}} \\
\addlinespace
INU-A31 & Systemic Feedback & $C_{t+1} = h(\sum_i a_i(t), C_t)$ & Reflexivity literature; Soros (1987) \\
INU-A34 & Spillover Effects & $U_d^{total} = U_d + \sum_{j \neq d} \tau_{jd} U_j$ & Cross-domain transfer \\
\addlinespace

\multicolumn{4}{@{}l}{\textit{\textbf{Block 15: Foundation (A35)}}} \\
\addlinespace
INU-A35 & Physiological Foundation & $\Phi_{phys} = \sigma(P_{phys} - \theta)$ & Maslow (1943); Basic needs \\
\addlinespace

\multicolumn{4}{@{}l}{\textit{\textbf{Block 16: Complementarity (A44--A57)}}} \\
\addlinespace
A44--A57 & Inter-dimensional Compl. & $C_{d_1,d_2} = \sigma U_{d_1} U_{d_2}$ & Milgrom \& Roberts (1990); Topkis (1998) \\

\end{longtable}

%-----------------------------------------------------------------------
\subsection{INU Parameter Calculation}
%-----------------------------------------------------------------------

The 59 INU axioms define how the four parameters ($G$, $P$, $\lambda$, $\delta$) 
are calculated.

\subsubsection{Parameter 1: Gain Calculation ($G^{INU}_{d,t}$)}

\begin{tcolorbox}[colback=blue!5!white, colframe=blue!50!black]
\begin{equation}
\boxed{
G^{INU}_{d,t} = \left[ V_{base,d} + V_{intr,d} + V_{extr,d} + \kappa_d V_{intr} V_{extr} + (x_{d,t} - R^{INU}_{d,t}) \right]^{+} \cdot \Omega^G_d
}
\end{equation}
\end{tcolorbox}

\textbf{Gain Modulator:}
\begin{equation}
\Omega^G_d = \underbrace{(1 - \phi_{risk} - \phi_{amb})}_{\text{A9}} \cdot 
             \underbrace{(1 - \phi_{motiv})}_{\text{A10}} \cdot 
             \underbrace{(1 - \phi_{blind})}_{\text{A11}} \cdot 
             \underbrace{S_d}_{\text{A15}} \cdot 
             \underbrace{(1 - H_d)}_{\text{A20}} \cdot 
             \underbrace{\alpha_d}_{\text{A32}}
\end{equation}

\textbf{Contributing Axioms:} A3, A4, A5, A6, A7, A9, A10, A11, A15, A20, A23, A26, A32, A34

\subsubsection{Parameter 2: Pain Calculation ($P^{INU}_{d,t}$)}

\begin{tcolorbox}[colback=blue!5!white, colframe=blue!50!black]
\begin{equation}
\boxed{
P^{INU}_{d,t} = \left[ R^{INU}_{d,t} - x_{d,t} \right]^{+} \cdot \Omega^P_d + \mathcal{E}^{-}_d
}
\end{equation}
\end{tcolorbox}

\textbf{Pain Modulator:}
\begin{equation}
\Omega^P_d = \underbrace{(1 - Cope_d) \cdot (1 - Res_d)}_{\text{A22: Stress buffer}}
\end{equation}

\textbf{Contributing Axioms:} A7, A9, A22, A25, A30, A34

\subsubsection{Parameter 3: Loss Aversion ($\lambda^{INU}_{d}$)}

\begin{tcolorbox}[colback=blue!5!white, colframe=blue!50!black]
\begin{equation}
\boxed{
\lambda^{INU}_{d} = \bar{\lambda}_d \cdot \Lambda_d(\Psi, Context)
}
\end{equation}
\end{tcolorbox}

\textbf{Context Modulator:}
\begin{equation}
\Lambda_d = \underbrace{\gamma_{gender,d}}_{\text{A17}} \cdot 
            \underbrace{\gamma_{culture,d}}_{\text{A18}} \cdot 
            \underbrace{f(G \times A \times S \times E)_d}_{\text{A21}} \cdot
            \underbrace{\gamma_{lifecycle,d}}_{\text{A16}}
\end{equation}

\textbf{Contributing Axioms:} A2, A16, A17, A18, A21

\subsubsection{Parameter 4: Discounting ($\delta^{INU}_{d}(t)$)}

\begin{tcolorbox}[colback=blue!5!white, colframe=blue!50!black]
\begin{equation}
\boxed{
\delta^{INU}_{d}(t) = \beta \cdot \bar{\delta}_d^{\,t} \cdot \Delta_d(t, Context)
}
\end{equation}
\end{tcolorbox}

\textbf{Contributing Axioms:} A8, A12, A13, A16, A33, A41

%-----------------------------------------------------------------------
\subsection{INU-Gain Components (24 cells)}
%-----------------------------------------------------------------------

\begin{table}[htbp]
\centering
\caption{INU-Gain: Individual Utility Gains by Dimension and Time}
\label{tab:app-c-inu-gain}
\small
\renewcommand{\arraystretch}{1.2}
\begin{tabular}{@{}lp{2.8cm}p{2.8cm}p{2.8cm}p{2.8cm}@{}}
\toprule
\textbf{Dim} & \textbf{Instant ($t=0$)} & \textbf{Short ($t=1$)} & \textbf{Medium ($t=2$)} & \textbf{Long ($t=3$)} \\
\midrule
F & Windfall, tip & Bonus, refund & Promotion, raise & Wealth accumulation \\
E & Joy, delight & Satisfaction & Fulfillment & Meaning, purpose \\
P & Energy burst & Recovery & Fitness gains & Long-term health \\
S & Compliment & Recognition & Status advancement & Legacy, reputation \\
D & Quick access & Efficiency gain & Digital competence & Digital capital \\
Eco & Nature moment & Clean environment & Sustainable living & Planet stewardship \\
\bottomrule
\end{tabular}
\end{table}

%-----------------------------------------------------------------------
\subsection{INU-Pain Components (24 cells)}
%-----------------------------------------------------------------------

\begin{table}[htbp]
\centering
\caption{INU-Pain: Individual Utility Losses by Dimension and Time}
\label{tab:app-c-inu-pain}
\small
\renewcommand{\arraystretch}{1.2}
\begin{tabular}{@{}lp{2.8cm}p{2.8cm}p{2.8cm}p{2.8cm}@{}}
\toprule
\textbf{Dim} & \textbf{Instant ($t=0$)} & \textbf{Short ($t=1$)} & \textbf{Medium ($t=2$)} & \textbf{Long ($t=3$)} \\
\midrule
F & Loss, theft & Expense, bill & Debt burden & Poverty, destitution \\
E & Frustration & Stress, worry & Anxiety, burnout & Depression, despair \\
P & Pain, injury & Illness, fatigue & Chronic condition & Disability, mortality \\
S & Embarrassment & Rejection & Social isolation & Loneliness, ostracism \\
D & Glitch, outage & Digital overload & Tech dependency & Digital exclusion \\
Eco & Pollution exposure & Environmental damage & Resource depletion & Climate catastrophe \\
\bottomrule
\end{tabular}
\end{table}

%%%%%%%%%%%%%%%%%%%%%%%%%%%%%%%%%%%%%%%%%%%%%%%%%%%%%%%%%%%%%%%%%%%%%%%%%%%%%%%
%%%%%%%%%%%%%%%%%%%%%%%%%%%%%%%%%%%%%%%%%%%%%%%%%%%%%%%%%%%%%%%%%%%%%%%%%%%%%%%
\section{KNU: Collective Utility}
\label{sec:app-c-knu}
%%%%%%%%%%%%%%%%%%%%%%%%%%%%%%%%%%%%%%%%%%%%%%%%%%%%%%%%%%%%%%%%%%%%%%%%%%%%%%%
%%%%%%%%%%%%%%%%%%%%%%%%%%%%%%%%%%%%%%%%%%%%%%%%%%%%%%%%%%%%%%%%%%%%%%%%%%%%%%%

\begin{cross-reference}
\textbf{Chapter Reference:} See Chapter~\ref{sec:welfare-fepsde}, 
Section~\ref{subsec:10-1-categories} for KNU definition.

\textbf{Level Reference:} KNU operates at all levels $L \geq 1$. Critically, 
KNU \textit{defines} which levels exist (see Appendix AAA, Axiom AL-2: 
``Level $L$ exists $\iff \exists IDN^L \lor \exists KNU^L$'').
\end{cross-reference}

\subsection{KNU Definition}

\begin{tcolorbox}[colback=green!5!white, colframe=green!75!black, 
    title=Collective Utility (KNU)]

\textbf{Core Question:} ``What do I want for US?''

\textbf{Level-Specific Notation:} $U^{KNU,L}$ denotes collective utility at level $L$.
The ``US'' in ``What do I want for US?'' is defined by the level:
\begin{itemize}[nosep]
\item $L = 1$: US = my family
\item $L = 2$: US = my tribe/community
\item $L = 3$: US = my organization
\item $L = 6$: US = my nation
\item $L = \infty$: US = all sentient beings
\end{itemize}

\textbf{Formula (from Meta-Axiom U-M0):}
\begin{equation}
U^{KNU} = \sum_{d \in \mathcal{D}} \sum_{t=0}^{3} \delta^{KNU}_{d}(t) \cdot 
\left[ G^{KNU}_{d,t} - \lambda^{KNU}_{d} \cdot P^{KNU}_{d,t} \right]
\end{equation}

\textbf{Reference Point Formation:}
\begin{equation}
R^{KNU}_{d,t} = \alpha_{norm} \cdot x^{group\_norm}_{d,t} + 
               \alpha_{fair} \cdot x^{fair\_distribution}_{d,t} + 
               \alpha_{peer} \cdot x^{peer\_groups}_{d,t}
\end{equation}

\textbf{Key Properties:}
\begin{itemize}
\item Other-regarding: Considers consequences for the relevant group
\item Fairness-weighted: Outcomes discounted by perceived unfairness
\item Conditional: Strength depends on group salience and identification
\end{itemize}

\end{tcolorbox}

%-----------------------------------------------------------------------
\subsection{Complete KNU Axiom Table (65 Axioms)}
%-----------------------------------------------------------------------

\begin{longtable}{@{}llp{4.2cm}p{4.2cm}@{}}
\caption{Complete KNU Axiom System (65 Axioms)} 
\label{tab:knu-axioms-complete} \\
\toprule
\textbf{Axiom} & \textbf{Title} & \textbf{Core Formula} & \textbf{Reference} \\
\midrule
\endfirsthead
\multicolumn{4}{c}{\tablename\ \thetable{} -- continued} \\
\toprule
\textbf{Axiom} & \textbf{Title} & \textbf{Core Formula} & \textbf{Reference} \\
\midrule
\endhead
\midrule
\multicolumn{4}{r}{Continued on next page...} \\
\endfoot
\bottomrule
\endlastfoot

\multicolumn{4}{@{}l}{\textit{\textbf{Block 1: Basis (A1--A7)}}} \\
\addlinespace
KNU-A1 & FEPSDE Structure & $N_{kol} = \sum_{d} G_d^{kol}$ & Ostrom (1990); Olson (1965) \\
KNU-A2 & Context Weights & $\sum \omega_d G_d^{kol}, \sum \omega = 1$ & Sandel (2012); Walzer (1983) \\
KNU-A3 & Pain/Gain Separation & $W = \sum G - \sum P$ & Kahneman \& Tversky (1979) \\
KNU-A4 & Asymmetry & $\partial W/\partial P > \partial W/\partial G$ & Loss aversion for public goods \\
KNU-A5 & Complementarity & $\kappa_{PG} = 0$ (no mixing) & Thaler (1999) Mental Accounting \\
KNU-A6 & Temporal Conditionality & $W_{tot} = \sum \beta_\tau U_\tau$ & Nordhaus (1994); Stern (2007) \\
KNU-A7 & Uncertainty Modulation & $\beta_\tau = \beta^0(1 - U_{env})$ & Pindyck (2007); Weitzman (2009) \\
\addlinespace

\multicolumn{4}{@{}l}{\textit{\textbf{Block 2: Cooperation (A8--A9)}}} \\
\addlinespace
KNU-A8 & Belief-based Cooperation & $N_{kol} \propto (\rho_K + f(B^{eff}))$ & Fischbacher, Gächter, Fehr (2001) \\
KNU-A9 & Actor Types & $\rho_K \in \{0, f(B), 1\}$ & Fehr \& Schmidt (1999); Kurzban \& Houser (2005) \\
\addlinespace

\multicolumn{4}{@{}l}{\textit{\textbf{Block 3: Dynamics (A10)}}} \\
\addlinespace
KNU-A10 & Life Changing Events & $\Delta U = LCE \cdot \theta$ & Taleb (2007); Scheffer et al. (2009) \\
\addlinespace

\multicolumn{4}{@{}l}{\textit{\textbf{Block 4: Goods Character (A11--A13)}}} \\
\addlinespace
KNU-A11 & Non-Excludability & $U_{ind} \ge 0$ if $N_{kol} > 0$ & Samuelson (1954) \\
KNU-A12 & Non-Rivalry & $\partial U_i / \partial U_j = 0$ & Stiglitz (1999) \\
KNU-A13 & Commons Rivalry & $\partial U_i / \partial U_j < 0$ & Hardin (1968); Ostrom (1990) \\
\addlinespace

\multicolumn{4}{@{}l}{\textit{\textbf{Block 5: FEPSDE Balances (A14--A19)}}} \\
\addlinespace
KNU-A14 & Social & $\gamma \cdot Coop - \delta \cdot Sanction$ & Coleman (1990); Putnam (2000) \\
KNU-A15 & Financial & $\gamma \cdot Exist - \delta \cdot Tax$ & Pigou (1920); Mirrlees (1971) \\
KNU-A16 & Psychological & $\gamma \cdot Trust - \delta \cdot Stress$ & Edmondson (1999) \\
KNU-A17 & Physical & $\gamma \cdot Safety - \delta \cdot Reg$ & Hobbes (1651) \\
KNU-A18 & Digital & $\gamma \cdot Access - \delta \cdot Surv$ & Zuboff (2019); Benkler (2006) \\
KNU-A19 & Ecological & $\gamma \cdot Sust - \delta \cdot Restr$ & Daly (1990) \\
\addlinespace

\multicolumn{4}{@{}l}{\textit{\textbf{Block 6: Temporal/Systemic (A20--A25)}}} \\
\addlinespace
KNU-A20 & Cumulative Experience & $\int (G - P) e^{-\mu t}$ & Acemoglu et al. (2005) \\
KNU-A21 & Intergenerational & $\sum \delta^g N(gen+g)$ & Rawls (1971) \\
KNU-A22 & Identity Complementarity & $N \cdot (1 + \kappa \cdot IDN)$ & Akerlof \& Kranton (2000) \\
KNU-A23 & Institutions & $f(N^0, Inst)$ & North (1990) \\
KNU-A24 & Resilience & $N \cdot (1 - \varepsilon)$ & Holling (1973) \\
KNU-A25 & Tipping Points & $\dot{N} = f(\dots) + \xi$ & Schelling (1978) \\
\addlinespace

\multicolumn{4}{@{}l}{\textit{\textbf{Block 7: Ethics/Governance (A26--A30)}}} \\
\addlinespace
KNU-A26 & Fairness & $\sum U_i - \phi \cdot Var(U_i)$ & Fehr \& Schmidt (1999); Rawls (1971) \\
KNU-A27 & Ethics/Privacy & s.t. Fairness, Privacy & Sen (1999); Nissenbaum (2009) \\
KNU-A28 & Transparency & $T(N)$ & Stiglitz (2002); Sunstein \\
KNU-A29 & Bias Control & $Bias \le \epsilon$ & Banerjee \& Duflo (2011) \\
KNU-A30 & Abuse Protection & $Protect(N)$ & Acemoglu \& Robinson (2012) \\
\addlinespace

\multicolumn{4}{@{}l}{\textit{\textbf{Block 8: Sanction/Enforcement (A31--A42)}}} \\
\addlinespace
KNU-A31 & Decentralized Punishment & $U - c$ (cost for punisher \& punished) & Fehr \& Gächter (2002) \\
KNU-A32 & Centralized Punishment & $U - P_{inst} \cdot I_{acc}$ & Becker (1968); Polinsky \& Shavell (2000) \\
KNU-A33 & Whistleblowing & $E[U] > E[C] \wedge Protect \ge \theta$ & Dyck, Morse, Zingales (2010) \\
KNU-A34 & Preventive Conformity & $f(NormStrength, Legit)$ & Cialdini (1984); Bicchieri (2006) \\
KNU-A35 & Institutional Quality & $Eff = f(Q, Acc, Fair)$ & Tyler (1990) \\
KNU-A36 & Elite Role Model & $P(Conf) = f(Elite)$ & Acemoglu \& Robinson (2006) \\
KNU-A37 & Positive Incentives & $U + r - c_r$ & Gneezy et al. (2011) \\
KNU-A38 & Rehabilitation & $U_{post} = U_{norm}$ & Bénabou \& Tirole (2011) \\
KNU-A39 & Inequality & $Impact(i) \neq Impact(j)$ & Piketty (2014); Glaeser et al. (2003) \\
KNU-A40 & Wrongful Punishment & $Err > 0 \Rightarrow Trust \downarrow$ & Greif (1993); Fehr (Betrayal Aversion) \\
KNU-A41 & Culture & $Form = f(Culture)$ & Henrich et al. (2006) \\
KNU-A42 & Dynamic Adaptation & $Pun = f(DefectionRate)$ & Boyd \& Richerson (1992) \\
\addlinespace

\multicolumn{4}{@{}l}{\textit{\textbf{Block 9: Dark Triad (A43--A46)}}} \\
\addlinespace
KNU-A43 & Antisocial Punishment & $Pun_{anti} = Pun(1-Legit)(1-Coh)$ & Herrmann, Thöni, Gächter (2008) \\
KNU-A44 & Dark Triad (Disposition) & $\Theta_{DT} = \alpha M + \beta N + \gamma P$ & Paulhus \& Williams (2002) \\
KNU-A45 & Behavioral Indicator & $\Theta_{Beh} = \sum \beta_i v_i$ & Stern (2007); Simon (1996) \\
KNU-A46 & Destruction Function & $P_{DTX} = P_{anti}(1+\Theta_{DT})(1+\Theta_{Beh})$ & FehrAdvice (2025) Synthesis \\
\addlinespace

\multicolumn{4}{@{}l}{\textit{\textbf{Block 10: System Dynamics (A47--A50)}}} \\
\addlinespace
KNU-A47 & Cooperation Complementarity & $\partial U_i / \partial U_j > 0$ & Topkis (1998); Bulow et al. (1985) \\
KNU-A48 & Feedback (Ind-Col) & $N_{t+1} = N_t + \eta \sum \Delta U_{ind}$ & Ostrom (2005) \\
KNU-A49 & Norm Coherence & $Coh_{t+1} = Coh_t + \gamma(Legit - Pun_{anti})$ & Bicchieri (2006); Sunstein (1996) \\
KNU-A50 & Systemic Equilibrium & $Stab = 1 - (w_1|\Delta N| + w_2|\Delta Coh| + w_3|\Delta\lambda|)$ & Nash (1950); Maynard Smith (1982) \\
\addlinespace

\multicolumn{4}{@{}l}{\textit{\textbf{Block 11: Inter-dimensional Complementarity (A51--A65)}}} \\
\addlinespace
KNU-A51 & F $\times$ E & $\sigma^{KNU}_{FE} \cdot U_F \cdot U_E$ & Easterlin (1974); Deaton (2008) \\
KNU-A52 & F $\times$ P & $\sigma^{KNU}_{FP} \cdot U_F \cdot U_P$ & Marmot (2004); Wilkinson \& Pickett (2009) \\
KNU-A53 & F $\times$ S & $\sigma^{KNU}_{FS} \cdot U_F \cdot U_S$ & Putnam (2000); Uslaner (2002) \\
KNU-A54 & F $\times$ D & $\sigma^{KNU}_{FD} \cdot U_F \cdot U_D$ & Benkler (2006); OECD Digital Economy \\
KNU-A55 & F $\times$ Eco & $\sigma^{KNU}_{F,Eco} \cdot U_F \cdot U_{Eco}$ & Ostrom (1990); Stern (2007) \\
KNU-A56 & E $\times$ P & $\sigma^{KNU}_{EP} \cdot U_E \cdot U_P$ & Helliwell \& Putnam (2004) \\
KNU-A57 & E $\times$ S & $\sigma^{KNU}_{ES} \cdot U_E \cdot U_S$ & Durkheim (1893); Collins (2004) \\
KNU-A58 & E $\times$ D & $\sigma^{KNU}_{ED} \cdot U_E \cdot U_D$ & Rheingold (2000); Wellman (2001) \\
KNU-A59 & E $\times$ Eco & $\sigma^{KNU}_{E,Eco} \cdot U_E \cdot U_{Eco}$ & Wilson (1984); Kals et al. (1999) \\
KNU-A60 & P $\times$ S & $\sigma^{KNU}_{PS} \cdot U_P \cdot U_S$ & Kawachi et al. (1997); Berkman \& Glass (2000) \\
KNU-A61 & P $\times$ D & $\sigma^{KNU}_{PD} \cdot U_P \cdot U_D$ & WHO Digital Health; Topol (2019) \\
KNU-A62 & P $\times$ Eco & $\sigma^{KNU}_{P,Eco} \cdot U_P \cdot U_{Eco}$ & WHO Environmental Health; Haines et al. (2006) \\
KNU-A63 & S $\times$ D & $\sigma^{KNU}_{SD} \cdot U_S \cdot U_D$ & Van Dijk (2020); Hargittai (2002) \\
KNU-A64 & S $\times$ Eco & $\sigma^{KNU}_{S,Eco} \cdot U_S \cdot U_{Eco}$ & Stern et al. (1999); Dietz et al. (2005) \\
KNU-A65 & D $\times$ Eco & $\sigma^{KNU}_{D,Eco} \cdot U_D \cdot U_{Eco}$ & GeSI (2015); Hilty \& Aebischer (2015) \\

\end{longtable}

%-----------------------------------------------------------------------
\subsection{KNU Parameter Calculation}
%-----------------------------------------------------------------------

The 65 KNU axioms define how the four parameters ($G$, $P$, $\lambda$, $\delta$) 
are calculated for collective utility.

\subsubsection{Parameter 1: Collective Gain ($G^{KNU}_{d,t}$)}

\begin{tcolorbox}[colback=green!5!white, colframe=green!50!black]
\begin{equation}
\boxed{
G^{KNU}_{d,t} = \left[ \sum_d \omega_d \cdot \gamma_d \cdot Benefit_d - R^{KNU}_{d,t} \right]^{+} \cdot \Omega^{G,KNU}_d \cdot (1 + \kappa_{IDN} \cdot IDN)
}
\end{equation}
\end{tcolorbox}

\textbf{Gain Modulator:}
\begin{equation}
\Omega^{G,KNU}_d = \underbrace{(1 - U_{env})}_{\text{A7: Uncertainty}} \cdot 
                   \underbrace{(\rho_K + f(B^{eff}))}_{\text{A8: Belief-Coop}} \cdot 
                   \underbrace{(1 - \varepsilon)}_{\text{A24: Resilience}}
\end{equation}

\textbf{Contributing Axioms:} 
\begin{itemize}
\item A1 (FEPSDE): $\sum_d G_d^{kol}$
\item A2 (Weights): $\omega_d$
\item A7 (Uncertainty): $(1 - U_{env})$
\item A8 (Belief-Coop): $(\rho_K + f(B^{eff}))$
\item A11-A12 (Public Goods): Non-excludability, non-rivalry
\item A14-A19 (FEPSDE Balances): $\gamma_d \cdot Benefit_d$
\item A22 (Identity Compl.): $(1 + \kappa_{IDN} \cdot IDN)$
\item A24 (Resilience): $(1 - \varepsilon)$
\item A47 (Coop Compl.): $\partial U_i / \partial U_j > 0$
\end{itemize}

\subsubsection{Parameter 2: Collective Pain ($P^{KNU}_{d,t}$)}

\begin{tcolorbox}[colback=green!5!white, colframe=green!50!black]
\begin{equation}
\boxed{
P^{KNU}_{d,t} = \left[ \sum_d \omega_d \cdot \delta_d \cdot Cost_d + R^{KNU}_{d,t} \right]^{+} \cdot \Omega^{P,KNU}_d + P_{sanction} + P_{DTX}
}
\end{equation}
\end{tcolorbox}

\textbf{Pain Modulator:}
\begin{equation}
\Omega^{P,KNU}_d = \underbrace{(1 + Err \cdot \Delta Trust)}_{\text{A40: Wrongful Punishment}}
\end{equation}

\textbf{Sanction Terms:}
\begin{align}
P_{sanction} &= \underbrace{c}_{\text{A31: Decentralized}} + \underbrace{P_{inst} \cdot I_{acc}}_{\text{A32: Centralized}} \\
P_{DTX} &= \underbrace{P_{anti}(1 + \Theta_{DT})(1 + \Theta_{Beh})}_{\text{A46: Dark Triad Destruction}}
\end{align}

\textbf{Contributing Axioms:}
\begin{itemize}
\item A3-A4 (Asymmetry): $\partial W/\partial P > \partial W/\partial G$
\item A13 (Commons): $\partial U_i / \partial U_j < 0$
\item A14-A19 (FEPSDE Balances): $\delta_d \cdot Cost_d$
\item A25 (Tipping Points): Shock term $\xi$
\item A31 (Decentralized Pun.): $c$
\item A32 (Centralized Pun.): $P_{inst} \cdot I_{acc}$
\item A40 (Wrongful Pun.): $Err \rightarrow Trust \downarrow$
\item A43-A46 (Dark Triad): $P_{DTX}$
\end{itemize}

\subsubsection{Parameter 3: Collective Loss Aversion ($\lambda^{KNU}_{d}$)}

\begin{tcolorbox}[colback=green!5!white, colframe=green!50!black]
\begin{equation}
\boxed{
\lambda^{KNU}_{d} = \bar{\lambda}^{KNU}_d \cdot \Lambda^{KNU}_d
}
\end{equation}
\end{tcolorbox}

\textbf{Context Modulator:}
\begin{equation}
\Lambda^{KNU}_d = \underbrace{f(Fairness)}_{\text{A26}} \cdot 
                  \underbrace{f(Q, Acc, Fair)}_{\text{A35: Inst. Quality}} \cdot 
                  \underbrace{f(Culture)}_{\text{A41}} \cdot 
                  \underbrace{f(Inst)}_{\text{A23}}
\end{equation}

\textbf{Contributing Axioms:}
\begin{itemize}
\item A4 (Asymmetry): $\lambda > 1$ baseline
\item A23 (Institutions): $f(Inst)$
\item A26 (Fairness): $\sum U_i - \phi \cdot Var(U_i)$
\item A35 (Inst. Quality): $f(Q, Acc, Fair)$
\item A41 (Culture): $f(Culture)$
\end{itemize}

\subsubsection{Parameter 4: Collective Discounting ($\delta^{KNU}_{d}(t)$)}

\begin{tcolorbox}[colback=green!5!white, colframe=green!50!black]
\begin{equation}
\boxed{
\delta^{KNU}_{d}(t) = \beta^{KNU} \cdot \bar{\delta}^{KNU,t}_d \cdot \Delta^{KNU}_d(t)
}
\end{equation}
\end{tcolorbox}

\textbf{Time Modulator:}
\begin{equation}
\Delta^{KNU}_d(t) = \underbrace{(1 - U_{env})}_{\text{A7: Uncertainty}} \cdot 
                    \underbrace{e^{-\mu t}}_{\text{A20: Cumulative}} \cdot 
                    \underbrace{\delta^g}_{\text{A21: Intergenerational}}
\end{equation}

\textbf{Contributing Axioms:}
\begin{itemize}
\item A6 (Temporal): $\sum \beta_\tau U_\tau$
\item A7 (Uncertainty): $(1 - U_{env})$
\item A20 (Cumulative): $\int (G-P) e^{-\mu t}$
\item A21 (Intergenerational): $\sum \delta^g N(gen+g)$
\item A24-A25 (Resilience/Tipping): $(1 - \varepsilon)$, $\xi$
\end{itemize}

\subsubsection{Stability Metric (A50)}

\begin{tcolorbox}[colback=green!5!white, colframe=green!50!black]
\begin{equation}
\boxed{
Stab_{KNU} = 1 - (w_1 \cdot |\Delta N_{kol}| + w_2 \cdot |\Delta NormCoh| + w_3 \cdot |\Delta \lambda_{pos}|)
}
\end{equation}
\textbf{Threshold:} $Stab_{KNU} \ge 0.95$ $\Rightarrow$ CE\_OK\_FULL (System stable)
\end{tcolorbox}

\subsubsection{Inter-dimensional Complementarities}

\begin{equation}
U^{KNU}_{total} = \sum_d U^{KNU}_d + \sum_{d_1 < d_2} \sigma^{KNU}_{d_1,d_2} \cdot U^{KNU}_{d_1} \cdot U^{KNU}_{d_2}
\end{equation}

\begin{table}[htbp]
\centering
\caption{$\sigma^{KNU}_{d_1,d_2}$: Collective Complementarity Coefficients}
\label{tab:knu-sigma-matrix}
\small
\renewcommand{\arraystretch}{1.3}
\begin{tabular}{@{}lccccc@{}}
\toprule
$\sigma^{KNU}$ & E & P & S & D & Eco \\
\midrule
F & A51 & A52 & A53 & A54 & A55 \\
E & -- & A56 & A57 & A58 & A59 \\
P & -- & -- & A60 & A61 & A62 \\
S & -- & -- & -- & A63 & A64 \\
D & -- & -- & -- & -- & A65 \\
\bottomrule
\end{tabular}
\end{table}

%-----------------------------------------------------------------------
\subsection{KNU-Gain Components (24 cells)}
%-----------------------------------------------------------------------

\begin{table}[htbp]
\centering
\caption{KNU-Gain: Collective Utility Gains by Dimension and Time}
\label{tab:app-c-knu-gain}
\small
\renewcommand{\arraystretch}{1.2}
\begin{tabular}{@{}lp{2.8cm}p{2.8cm}p{2.8cm}p{2.8cm}@{}}
\toprule
\textbf{Dim} & \textbf{Instant ($t=0$)} & \textbf{Short ($t=1$)} & \textbf{Medium ($t=2$)} & \textbf{Long ($t=3$)} \\
\midrule
F & Group bonus & Team reward & Profit sharing & Generational wealth \\
E & Collective joy & Team spirit & Community pride & Shared cultural meaning \\
P & Emergency help & Group healthcare & Public health gains & Societal longevity \\
S & Group recognition & Team cohesion & Community trust & Social capital legacy \\
D & Shared access & Group connectivity & Digital commons & Collective infrastructure \\
Eco & Clean moment & Local improvement & Regional sustainability & Climate stability \\
\bottomrule
\end{tabular}
\end{table}

%-----------------------------------------------------------------------
\subsection{KNU-Pain Components (24 cells)}
%-----------------------------------------------------------------------

\begin{table}[htbp]
\centering
\caption{KNU-Pain: Collective Utility Losses by Dimension and Time}
\label{tab:app-c-knu-pain}
\small
\renewcommand{\arraystretch}{1.2}
\begin{tabular}{@{}lp{2.8cm}p{2.8cm}p{2.8cm}p{2.8cm}@{}}
\toprule
\textbf{Dim} & \textbf{Instant ($t=0$)} & \textbf{Short ($t=1$)} & \textbf{Medium ($t=2$)} & \textbf{Long ($t=3$)} \\
\midrule
F & Collective loss & Group expense & Organizational debt & Systemic poverty \\
E & Shared grief & Team stress & Community anxiety & Collective trauma \\
P & Group injury & Team illness & Public health crisis & Generational harm \\
S & Group shame & Team conflict & Community breakdown & Social collapse \\
D & System outage & Network failure & Digital divide & Technological exclusion \\
Eco & Pollution event & Environmental damage & Resource depletion & Climate catastrophe \\
\bottomrule
\end{tabular}
\end{table}

%%%%%%%%%%%%%%%%%%%%%%%%%%%%%%%%%%%%%%%%%%%%%%%%%%%%%%%%%%%%%%%%%%%%%%%%%%%%%%%
%%%%%%%%%%%%%%%%%%%%%%%%%%%%%%%%%%%%%%%%%%%%%%%%%%%%%%%%%%%%%%%%%%%%%%%%%%%%%%%
\section{IDN: Identity Utility}
\label{sec:app-c-idn}
%%%%%%%%%%%%%%%%%%%%%%%%%%%%%%%%%%%%%%%%%%%%%%%%%%%%%%%%%%%%%%%%%%%%%%%%%%%%%%%
%%%%%%%%%%%%%%%%%%%%%%%%%%%%%%%%%%%%%%%%%%%%%%%%%%%%%%%%%%%%%%%%%%%%%%%%%%%%%%%

\begin{cross-reference}
\textbf{Chapter Reference:} See Chapter~\ref{sec:welfare-fepsde}, 
Section~\ref{subsec:10-1-categories} for IDN definition.

\textbf{Level Reference:} IDN operates at all levels $L \geq 0$. Critically, 
IDN \textit{defines} which levels exist (see Appendix AAA, Axiom AL-2). 
Furthermore, IDN creates a \textit{hierarchy} of identities (see AAA, Axiom AL-7: 
Identity Hierarchy).
\end{cross-reference}

\subsection{IDN Definition}

\begin{tcolorbox}[colback=orange!5!white, colframe=orange!75!black, 
    title=Identity Utility (IDN)]

\textbf{Core Question:} ``Who AM I?''

\textbf{Level-Specific Notation:} $U^{IDN,L}$ denotes identity utility at level $L$.
The ``I'' in ``Who AM I?'' can be defined at multiple levels:
\begin{itemize}[nosep]
\item $L = 0$: I = individual self
\item $L = 1$: I = family member (``I am a Smith'')
\item $L = 2$: I = tribe/community member (``I am a New Yorker'')
\item $L = 3$: I = organizational member (``I am a Googler'')
\item $L = 6$: I = national citizen (``I am American'')
\item $L = \infty$: I = sentient being (``I am human'')
\end{itemize}

\textbf{Formula (from Meta-Axiom U-M0):}
\begin{equation}
U^{IDN} = \sum_{d \in \mathcal{D}} \sum_{t=0}^{3} \delta^{IDN}_{d}(t) \cdot 
\left[ G^{IDN}_{d,t} - \lambda^{IDN}_{d} \cdot P^{IDN}_{d,t} \right]
\end{equation}

\textbf{Reference Point Formation (Identity-Prescribed State):}
\begin{equation}
R^{IDN}_{d,t} = x^{*}_{d,t} = f(\text{roles}, \text{values}, \text{group\_membership}, \text{narrative})
\end{equation}

\textbf{Key Properties:}
\begin{itemize}
\item Self-definitional: About who one IS, not what one HAS
\item Can be instant: Identity threats activate immediately
\item Non-fungible: Cannot be compensated by INU gains
\item Asymmetric: Falling short often worse than exceeding
\item \textbf{Level-defining}: Shared identity creates social levels (see AAA, Axiom AL-2)
\end{itemize}

\end{tcolorbox}

%-----------------------------------------------------------------------
\subsection{Complete IDN Axiom Table (95 Axioms)}
%-----------------------------------------------------------------------

\begin{longtable}{@{}llp{4.2cm}p{4.2cm}@{}}
\caption{Complete IDN Axiom System (95 Axioms)} 
\label{tab:idn-axioms-complete} \\
\toprule
\textbf{Axiom} & \textbf{Title} & \textbf{Core Formula} & \textbf{Reference} \\
\midrule
\endfirsthead
\multicolumn{4}{c}{\tablename\ \thetable{} -- continued} \\
\toprule
\textbf{Axiom} & \textbf{Title} & \textbf{Core Formula} & \textbf{Reference} \\
\midrule
\endhead
\midrule
\multicolumn{4}{r}{Continued on next page...} \\
\endfoot
\bottomrule
\endlastfoot

\multicolumn{4}{@{}l}{\textit{\textbf{Block 1: Basis (A1--A8)}}} \\
\addlinespace
IDN-A1 & Definition & $ID = f(\text{Belonging}, \text{Difference})$ & Tajfel \& Turner (1979); Akerlof \& Kranton (2000) \\
IDN-A2 & Multiplicity & $U_{tot} = \sum \pi_j U^j$ & Sen (2006); Stryker (1980) \\
IDN-A3 & Independence & $U_{IDN} \neq U_{INU} + U_{KNU}$ & Akerlof \& Kranton (2000) \\
IDN-A4 & Net Utility & $(\alpha \text{Recog} + \dots) - (\theta \text{Exclus} + \dots)$ & Honneth (1992); Goffman (1963) \\
IDN-A5 & Temporality & $\sum \delta \omega (G - C)$ & Frederick et al. (2002) \\
IDN-A6 & Aggregation & $\sum \pi_j U^j$ (weighting focus) & Roccas \& Brewer (2002) \\
IDN-A7 & Sociality & $U^j = f(\text{Self}, \text{Others})$ & Mead (1934); Festinger (1954) \\
IDN-A8 & Context & $\pi_{j,context} = f(KON)$ & Turner et al. (1987) \\
\addlinespace

\multicolumn{4}{@{}l}{\textit{\textbf{Block 2: Awareness/Dynamics (A9--A12)}}} \\
\addlinespace
IDN-A9 & Motivated Awareness & $U_{eff} = \sum A_j \pi_j U^j$ & Bénabou \& Tirole (2016) \\
IDN-A10 & Complementarity & $+ \sum \kappa_{jk} \text{Interaction}(j,k)$ & Crenshaw (1989); Milgrom \& Roberts (1990) \\
IDN-A11 & Path Dependence & $U_t = f(U_{t-1}, \text{Events})$ & Sydow et al. (2009) \\
IDN-A12 & Identity Journey & $\{\text{Build}, \text{Crisis}, \text{Trans}, \text{Stable}\}$ & Erikson (1959); Marcia (1966) \\
\addlinespace

\multicolumn{4}{@{}l}{\textit{\textbf{Block 3: Norms/Institutions (A13--A15, A41, A44, A59)}}} \\
\addlinespace
IDN-A13 & Norms & $+ \eta \text{Conform} - \zeta \text{Conflict}$ & Bicchieri (2006); Akerlof (1980) \\
IDN-A14 & Institutionalization & $C = f(\text{Stab}_{inst})$ & North (1990) \\
IDN-A15 & Uncertainty & $Var(U) = f(U_{ep}, U_{ip}, \dots)$ & Hogg (2007) \\
IDN-A41 & Norm Integration & (see A13, reinforcement focus) & Parsons (1951) \\
IDN-A44 & Adaptive Uncertainty & $\Delta U = -[\phi_1 \text{Uncert} - \phi_2 A]$ & Heiner (1983) \\
IDN-A59 & Legitimacy & $U_{eff} = U + \zeta \text{Legit}$ & Tyler (2006) \\
\addlinespace

\multicolumn{4}{@{}l}{\textit{\textbf{Block 4: Temporal (A16--A20, A31--A35)}}} \\
\addlinespace
IDN-A16 & Temporal Integration & $\sum \delta \omega U_i$ (inst-long) & Loewenstein \& Prelec (1992) \\
IDN-A17 & Temporal Profiles & $Profil^j = \{\omega_{inst}, \omega_{short}, \omega_{mid}, \omega_{long}\}$ & Frederick et al. (2002) \\
IDN-A18 & Discounting & $\sum \delta (G - C)$ & Laibson (1997) \\
IDN-A19 & Temporal Dynamics & $\omega_t^j = f(\text{LifePhase}, \text{Events})$ & Carstensen (2006) \\
IDN-A20 & Temporal Feedback & $\Delta U_{long} = g(U_{short})$ & Becker \& Murphy (1988) \\
IDN-A31 & Temporal Gains & $G^j = \{G_{inst}, \dots, G_{long}\}$ & Samuelson (1937) \\
IDN-A32 & Temporal Pains & $C^j = \{C_{inst}, \dots, C_{long}\}$ & O'Donoghue \& Rabin (1999) \\
IDN-A33 & Time Definition Type & $TimeProfile = f(Type)$ & Braudel (1958) \\
IDN-A34 & Reversibility & $\Delta U \in \{+,-\}$ reversible & Lewin (1947) \\
IDN-A35 & Time Feedback & $\Delta U_{long} = g(U_{short})$ & Becker \& Murphy (1988) \\
\addlinespace

\multicolumn{4}{@{}l}{\textit{\textbf{Block 5: Gains/Pains (A21--A23, A36--A37)}}} \\
\addlinespace
IDN-A21 & Identity Gains & $\alpha \text{Recog} + \beta \text{Belong} + \dots$ & Maslow (1943); Baumeister \& Leary (1995) \\
IDN-A22 & Identity Pains & $\theta \text{Exclus} + \lambda \text{Stigma} + \dots$ & Williams (2007); Eisenberger (2003) \\
IDN-A23 & Asymmetry (Magnitude) & $|C_i^j| \gg |G_i^j|$ for $\Delta x_{gain} = \Delta x_{loss}$ & Kahneman \& Tversky (1979) \\
IDN-A36 & Independence (Costs) & $C_{IDN} \neq C_{INU}$ & Akerlof \& Kranton (2000) \\
IDN-A37 & Independence (Gains) & $G_{IDN} \neq G_{INU}$ & Akerlof \& Kranton (2000) \\
\addlinespace

\multicolumn{4}{@{}l}{\textit{\textbf{Block 6: Complementarity/Conflicts (A24--A30, A40)}}} \\
\addlinespace
IDN-A24 & Complementary Gains & $G_{tot} = \sum G + \sum \kappa^+ (G^j G^k)$ & Milgrom \& Roberts (1995) \\
IDN-A25 & Conflict Costs & $C_{tot} = \sum C + \sum \kappa^- \text{Conflict}(j,k)$ & Festinger (1957) \\
IDN-A26 & Dominance & $U^{primed} = \max(\pi \cdot U^j)$ & Stryker (1980) \\
IDN-A27 & Stability & $Stab = f(Inst, Persistence)$ & Marcia (1966) \\
IDN-A28 & Strength & $Strength = f(Salience, Emotion, Relevance)$ & Tajfel (1974) \\
IDN-A29 & Context Dependence & $\pi_{context} = f(KON)$ & Turner et al. (1987) \\
IDN-A30 & Relationality & $U = f(\text{Belonging}, \text{Distinction})$ & Barth (1969) \\
IDN-A40 & Contextual Conflicts & $C_{Conflict} = f(KON)$ & Goffman (1959) \\
\addlinespace

\multicolumn{4}{@{}l}{\textit{\textbf{Block 7: Segments/Interventions (A38--A39, A42--A43, A47)}}} \\
\addlinespace
IDN-A38 & Asymmetry (Persistence) & $\int_t^\infty C_{IDN}^j d\tau > \int_t^\infty G_{IDN}^j d\tau$ & Baumeister et al. (2001) \\
IDN-A39 & Segment Differences & $U_{seg} = f(Age, Gender, Culture)$ & Hofstede (2001) \\
IDN-A42 & Interventions & $\Delta U = f(\text{Symbols}, \text{Narratives})$ & Thaler \& Sunstein (2008) \\
IDN-A43 & Identity Journeys & $\{Build, Crisis, Trans, Stable\}$ & Erikson (1959) \\
IDN-A47 & Fragmentation & $C_{Frag} \propto \text{Count}(Contradictions)$ & Donahue et al. (1993) \\
\addlinespace

\multicolumn{4}{@{}l}{\textit{\textbf{Block 8: Macro/Resources (A49--A52, A54--A56)}}} \\
\addlinespace
IDN-A49 & Macro Level & $U = f(\text{Culture}, \text{Nation})$ & Anderson (1983) \\
IDN-A50 & Resource Binding & $C = \tau(\text{Time} + \text{Money} + \text{Energy})$ & Spence (1973) \\
IDN-A51 & Unconscious Identity & $U_{eff} = A \cdot U + (1-A)\epsilon U$ & Greenwald \& Banaji (1995) \\
IDN-A52 & Shocks & $\Delta U = f(\text{Shock}, \text{Resilience})$ & Taleb (2012) \\
IDN-A54 & Dynamic Weighting & $\pi_t = f(\pi_{t-1}, \text{Events})$ & Sydow (2009) \\
IDN-A55 & Multiplicative Effects & $\Pi (1+\kappa)$ & Crenshaw (1989) \\
IDN-A56 & Collective Identity & $U_{coll} = f(IDN, KNU)$ & Tajfel \& Turner (1979) \\
\addlinespace

\multicolumn{4}{@{}l}{\textit{\textbf{Block 9: Emotions/Complexity (A57, A71--A72, A74)}}} \\
\addlinespace
IDN-A57 & Emotional Intensity & $U_{emo} = \psi \text{Emotion} \cdot U_{IDN}$ & Elster (1996); Loewenstein (1996) \\
IDN-A71 & Complexity Anchor & $\text{Salience} = \alpha \cdot K(t)$ & Luhmann (1979) \\
IDN-A72 & Saturation Priority & $U_{IDN} > U_{INU} + U_{KNU}$ at $Sat > 1$ & Simon (1957) \\
IDN-A74 & Polarization & $+ \beta K$ (small) vs $-\gamma K$ (large) & Sunstein (2002) \\
\addlinespace

\multicolumn{4}{@{}l}{\textit{\textbf{Block 10: Coherence/Activation (A70, A70b, A75--A76)}}} \\
\addlinespace
IDN-A70 & Coherence Principle & $U_{tot} = \sum U^j - \Omega \cdot \text{Fragmentation}$ & Antonovsky (1979); Festinger (1957) \\
IDN-A70b & Self-Image Function & $\Delta I_t = f(|I_t - S_t|, A_j, E^j, K^j)$ & Cooley (1902); Festinger (1954) \\
IDN-A75 & Latent vs. Salient & $U_{eff} = U_{pot} \cdot A(t)$ & Freud (1915); Kahneman (2011) \\
IDN-A76 & Priming & $A(t) = f(\text{Prime}, \text{Context})$ & Bargh (2006) \\
\addlinespace

\multicolumn{4}{@{}l}{\textit{\textbf{Block 11: Bridges/Systemic (A45--A46, A48, A53, A58)}}} \\
\addlinespace
IDN-A45 & Bridge Function & $f(U_{INU}, U_{KNU})$ & Granovetter (1973) \\
IDN-A46 & Digital Identity & $U_{dig} \neq U_{phys}$ & Belk (2013) \\
IDN-A48 & Intergenerational & $\rho (\text{Continuity} - \text{Break})$ & Bourdieu (1986) \\
IDN-A53 & Intersectionality & $\sum \chi_{jk} \text{Intersection}$ & Crenshaw (1989); Cole (2009) \\
IDN-A58 & Identity Investment & $ROI = (U - C_{inv}) / C_{inv}$ & Bourdieu (1986) \\
\addlinespace

\multicolumn{4}{@{}l}{\textit{\textbf{Block 12: External/Resilience/Meta (A60--A69, A73, A77)}}} \\
\addlinespace
IDN-A60 & External Awareness & $U = \sum A_{other} \cdot \text{Effect}$ & Cooley (1902); Goffman \\
IDN-A61 & Parallel Journeys & $J_{tot} = \{J_{Work}, J_{Private}, \dots\}$ & Luhmann (1984) \\
IDN-A62 & Reversibility & $\Delta U \in \{+,-\}$ & Ebaugh (1988) \\
IDN-A63 & Neutrality & $U^{neutral} = 0$ & Freud (Latency) \\
IDN-A64 & Conflict Hierarchy & $C_{tot} = \max(C_{Conflict})$ & Maslow (1943) \\
IDN-A65 & Mobility & $Mob = f(Stab, Inst)$ & Bourdieu (1984) \\
IDN-A66 & Global vs. Local & $U = f(N_{glob}, N_{loc})$ & Robertson (1995) \\
IDN-A67 & Resilience & $Res = \Delta U / \text{Shock}$ & Bonanno (2004) \\
IDN-A68 & Identity ROI & $ROI = (Utility - Cost)/Inv$ & Becker (1964) \\
IDN-A69 & Meta-Relationality & $f(\text{Self}, \text{Other})$ & Mead (1934) \\
IDN-A73 & Segmented Response & $\Delta U^s = f(K, \kappa_s)$ & Kegan (1982) \\
IDN-A77 & Dynamic Multiplicity & $U_{tot} = \sum U_{pot} \cdot A \cdot w$ & Roccas \& Brewer (2002) \\
\addlinespace

\multicolumn{4}{@{}l}{\textit{\textbf{Block 13: Synthesis/Collapse (A78--A79)}}} \\
\addlinespace
IDN-A78 & Identity Synthesis & $Action(t) = \text{ArgMax}(\sum_j w_j U_{IDN}^j)$ & Ricoeur (1990); Korsgaard (2009) \\
IDN-A79 & Identity Collapse & $State = \begin{cases} Stable & \int(G-C)>\epsilon \\ Collapse & Coh < Min \end{cases}$ & Durkheim (1897); Erikson (1968) \\
\addlinespace

\multicolumn{4}{@{}l}{\textit{\textbf{Block 14: Inter-dimensional Complementarity (A80--A94)}}} \\
\addlinespace
IDN-A80 & F $\times$ E & $\sigma^{IDN}_{FE} \cdot U_F \cdot U_E$ & Lea \& Webley (2006); Belk (1988) \\
IDN-A81 & F $\times$ P & $\sigma^{IDN}_{FP} \cdot U_F \cdot U_P$ & Marmot (2004) \\
IDN-A82 & F $\times$ S & $\sigma^{IDN}_{FS} \cdot U_F \cdot U_S$ & Bourdieu (1984); Veblen (1899) \\
IDN-A83 & F $\times$ D & $\sigma^{IDN}_{FD} \cdot U_F \cdot U_D$ & Belk (2013) \\
IDN-A84 & F $\times$ Eco & $\sigma^{IDN}_{F,Eco} \cdot U_F \cdot U_{Eco}$ & Stern et al. (1999) \\
IDN-A85 & E $\times$ P & $\sigma^{IDN}_{EP} \cdot U_E \cdot U_P$ & Damasio (1994) \\
IDN-A86 & E $\times$ S & $\sigma^{IDN}_{ES} \cdot U_E \cdot U_S$ & Baumeister \& Leary (1995) \\
IDN-A87 & E $\times$ D & $\sigma^{IDN}_{ED} \cdot U_E \cdot U_D$ & Turkle (2011) \\
IDN-A88 & E $\times$ Eco & $\sigma^{IDN}_{E,Eco} \cdot U_E \cdot U_{Eco}$ & Clayton (2003) \\
IDN-A89 & P $\times$ S & $\sigma^{IDN}_{PS} \cdot U_P \cdot U_S$ & Goffman (1963); Butler (1990) \\
IDN-A90 & P $\times$ D & $\sigma^{IDN}_{PD} \cdot U_P \cdot U_D$ & Yee \& Bailenson (2007) \\
IDN-A91 & P $\times$ Eco & $\sigma^{IDN}_{P,Eco} \cdot U_P \cdot U_{Eco}$ & Wilson (1984); Schultz (2002) \\
IDN-A92 & S $\times$ D & $\sigma^{IDN}_{SD} \cdot U_S \cdot U_D$ & boyd (2014); Wellman (2001) \\
IDN-A93 & S $\times$ Eco & $\sigma^{IDN}_{S,Eco} \cdot U_S \cdot U_{Eco}$ & Fielding et al. (2008) \\
IDN-A94 & D $\times$ Eco & $\sigma^{IDN}_{D,Eco} \cdot U_D \cdot U_{Eco}$ & Longo et al. (2019) \\

\end{longtable}

%-----------------------------------------------------------------------
\subsection{IDN Parameter Calculation}
%-----------------------------------------------------------------------

The 95 IDN axioms define how the four parameters ($G$, $P$, $\lambda$, $\delta$) 
are calculated for identity utility.

\subsubsection{Parameter 1: Identity Gain ($G^{IDN}_{d,t}$)}

\begin{tcolorbox}[colback=orange!5!white, colframe=orange!50!black]
\begin{equation}
\boxed{
G^{IDN}_{d,t} = \sum_j \pi_j(t) \cdot A_j(t) \cdot \left[ \alpha Recog_j + \beta Belong_j + \gamma Meaning_j - R^{IDN}_{d,t} \right]^{+} \cdot \Omega^{G,IDN}_d
}
\end{equation}
\end{tcolorbox}

\textbf{Gain Modulator:}
\begin{equation}
\Omega^{G,IDN}_d = \underbrace{A_j(t)}_{\text{A9: Awareness}} \cdot 
                   \underbrace{(1 + \sum_k \kappa^+_{jk})}_{\text{A24: Compl. Gains}} \cdot 
                   \underbrace{\psi \cdot Emotion}_{\text{A57: Intensity}}
\end{equation}

\textbf{Contributing Axioms:}
\begin{itemize}
\item A1 (Definition): $f(Belonging, Difference)$
\item A2 (Multiplicity): $\sum \pi_j U^j$
\item A4 (Net Utility): $\alpha Recog + \beta Belong + ...$
\item A9 (Motivated Awareness): $A_j$
\item A21 (Identity Gains): $\alpha Recog + \beta Belong + \gamma Meaning$
\item A24 (Complementary Gains): $\sum \kappa^+ (G^j G^k)$
\item A26 (Dominance): $\max(\pi \cdot U^j)$
\item A28 (Strength): $f(Salience, Emotion, Relevance)$
\item A57 (Emotional Intensity): $\psi \cdot Emotion$
\item A75-A76 (Activation): $A(t) = f(Prime, Context)$
\end{itemize}

\subsubsection{Parameter 2: Identity Pain ($P^{IDN}_{d,t}$)}

\begin{tcolorbox}[colback=orange!5!white, colframe=orange!50!black]
\begin{equation}
\boxed{
P^{IDN}_{d,t} = \sum_j \pi_j(t) \cdot \left[ \theta Exclus_j + \lambda Stigma_j \right] \cdot \Omega^{P,IDN}_d + C_{conflict} + C_{frag}
}
\end{equation}
\end{tcolorbox}

\textbf{Pain Modulator:}
\begin{equation}
\Omega^{P,IDN}_d = \underbrace{\psi \cdot Emotion}_{\text{A57: Intensity}}
\end{equation}

\textbf{Conflict and Fragmentation Terms:}
\begin{align}
C_{conflict} &= \underbrace{\sum_{j \neq k} \kappa^-_{jk} \cdot |U^j - U^k|}_{\text{A25: Conflict Costs}} \\
C_{frag} &= \underbrace{\Omega \cdot Count(Contradictions)}_{\text{A47, A70: Fragmentation}}
\end{align}

\textbf{Contributing Axioms:}
\begin{itemize}
\item A4 (Net Utility): $\theta Exclus + \lambda Stigma + ...$
\item A22 (Identity Pains): $\theta Exclus + \lambda Stigma$
\item A23 (Asymmetry Magnitude): $|C| \gg |G|$
\item A25 (Conflict Costs): $\sum \kappa^- Conflict(j,k)$
\item A38 (Asymmetry Persistence): $\int C d\tau > \int G d\tau$
\item A40 (Contextual Conflicts): $f(KON)$
\item A47 (Fragmentation): $C_{Frag} \propto Count(Contradictions)$
\item A57 (Emotional Intensity): $\psi \cdot Emotion$
\item A64 (Conflict Hierarchy): $\max(C_{Conflict})$
\item A70 (Coherence): $\Omega \cdot Fragmentation$
\end{itemize}

\subsubsection{Parameter 3: Identity Loss Aversion ($\lambda^{IDN}_{d}$)}

\begin{tcolorbox}[colback=orange!5!white, colframe=orange!50!black]
\begin{equation}
\boxed{
\lambda^{IDN}_{d} = \bar{\lambda}^{IDN}_d \cdot \Lambda^{IDN}_d
}
\end{equation}
\end{tcolorbox}

\textbf{Context Modulator:}
\begin{equation}
\Lambda^{IDN}_d = \underbrace{f(Stability)}_{\text{A27}} \cdot 
                  \underbrace{f(Salience, Emotion)}_{\text{A28}} \cdot 
                  \underbrace{\int_t^\infty Persistence(\tau) d\tau}_{\text{A38: Persistence}}
\end{equation}

\textbf{Note:} Identity loss aversion is typically higher than INU/KNU:
\begin{equation}
\lambda^{IDN}_d > \lambda^{INU}_d \quad \text{and} \quad \lambda^{IDN}_d > \lambda^{KNU}_d
\end{equation}
due to persistence effects (A38): identity pains are ``stickier'' than other pains.

\textbf{Contributing Axioms:}
\begin{itemize}
\item A23 (Asymmetry Magnitude): $|C| \gg |G|$
\item A27 (Stability): $f(Inst, Persistence)$
\item A28 (Strength): $f(Salience, Emotion, Relevance)$
\item A38 (Asymmetry Persistence): $\int_t^\infty C(\tau) d\tau > \int_t^\infty G(\tau) d\tau$
\item A57 (Emotional Intensity): $\psi \cdot Emotion$
\end{itemize}

\subsubsection{Parameter 4: Identity Discounting ($\delta^{IDN}_{d}(t)$)}

\begin{tcolorbox}[colback=orange!5!white, colframe=orange!50!black]
\begin{equation}
\boxed{
\delta^{IDN}_{d}(t) = \beta^{IDN} \cdot \bar{\delta}^{IDN,t}_d \cdot \Delta^{IDN}_d(t)
}
\end{equation}
\end{tcolorbox}

\textbf{Time Modulator:}
\begin{equation}
\Delta^{IDN}_d(t) = \underbrace{f(LifePhase)}_{\text{A19}} \cdot 
                    \underbrace{f(IdentityType)}_{\text{A33}} \cdot 
                    \underbrace{\omega_{profile}^j}_{\text{A17: Temporal Profile}}
\end{equation}

\textbf{Temporal Profile (A17):}
\begin{equation}
Profile^j = \{\omega_{inst}^j, \omega_{short}^j, \omega_{mid}^j, \omega_{long}^j\}
\end{equation}
Each identity has a specific ``temporal fingerprint'' -- e.g., institutional identities 
discount slower than social identities (A33: Braudel's \textit{longue durée}).

\textbf{Contributing Axioms:}
\begin{itemize}
\item A5 (Temporality): $\sum \delta \omega (G - C)$
\item A16 (Temporal Integration): $\sum \delta \omega U_i$
\item A17 (Temporal Profiles): $\{\omega_{inst}, \omega_{short}, \omega_{mid}, \omega_{long}\}$
\item A18 (Discounting): $\sum \delta (G - C)$
\item A19 (Temporal Dynamics): $f(LifePhase, Events)$
\item A33 (Time Definition Type): $f(Type)$
\item A35 (Time Feedback): $\Delta U_{long} = g(U_{short})$
\end{itemize}

\subsubsection{Synthesis Function (A78)}

\begin{tcolorbox}[colback=orange!5!white, colframe=orange!50!black]
\begin{equation}
\boxed{
Action(t) = \text{ArgMax}\left( \sum_j w_j(t) \cdot U^j_{IDN} \right)
}
\end{equation}
\textbf{Interpretation:} Despite multiplicity, the individual synthesizes all active identity vectors into a single action direction at each moment.
\end{tcolorbox}

\subsubsection{Collapse Condition (A79)}

\begin{tcolorbox}[colback=red!5!white, colframe=red!50!black]
\begin{equation}
\boxed{
State(t) = \begin{cases} 
\text{Stable} & \text{if } \int (G - C) dt > \epsilon \text{ and } Coherence > Min \\
\text{Collapse} & \text{if } Coherence < Min_{critical}
\end{cases}
}
\end{equation}
\textbf{Warning:} Identity collapse leads to anomie, radicalization, or pathology. This is the tipping point of the identity module.
\end{tcolorbox}

\subsubsection{Inter-dimensional Complementarities}

\begin{equation}
U^{IDN}_{total} = \sum_d U^{IDN}_d + \sum_{d_1 < d_2} \sigma^{IDN}_{d_1,d_2} \cdot U^{IDN}_{d_1} \cdot U^{IDN}_{d_2}
\end{equation}

\begin{table}[htbp]
\centering
\caption{$\sigma^{IDN}_{d_1,d_2}$: Identity Complementarity Coefficients}
\label{tab:idn-sigma-matrix}
\small
\renewcommand{\arraystretch}{1.3}
\begin{tabular}{@{}lccccc@{}}
\toprule
$\sigma^{IDN}$ & E & P & S & D & Eco \\
\midrule
F & A80 & A81 & A82 & A83 & A84 \\
E & -- & A85 & A86 & A87 & A88 \\
P & -- & -- & A89 & A90 & A91 \\
S & -- & -- & -- & A92 & A93 \\
D & -- & -- & -- & -- & A94 \\
\bottomrule
\end{tabular}
\end{table}

\subsubsection{Identity Coherence Effect}

\begin{tcolorbox}[colback=orange!5!white, colframe=orange!50!black]
\textbf{Coherence Index:}
\begin{equation}
K^{IDN} = \frac{\sum_{d_1 < d_2} \mathbb{1}[\sigma_{d_1,d_2} > 0]}{15}
\end{equation}

Positive complementarities ($\sigma > 0$) strengthen identity coherence.
Negative complementarities ($\sigma < 0$) create identity conflicts.

\textbf{Example:} ``Career identity'' vs. ``Parent identity'' may have $\sigma_{FS} < 0$.
\end{tcolorbox}

%-----------------------------------------------------------------------
\subsection{IDN-Gain Components (24 cells)}
%-----------------------------------------------------------------------

\begin{table}[htbp]
\centering
\caption{IDN-Gain: Identity Utility Gains by Dimension and Time}
\label{tab:app-c-idn-gain}
\small
\renewcommand{\arraystretch}{1.2}
\begin{tabular}{@{}lp{2.8cm}p{2.8cm}p{2.8cm}p{2.8cm}@{}}
\toprule
\textbf{Dim} & \textbf{Instant ($t=0$)} & \textbf{Short ($t=1$)} & \textbf{Medium ($t=2$)} & \textbf{Long ($t=3$)} \\
\midrule
F & ``I earned this'' & ``I'm succeeding'' & ``I'm financially competent'' & ``I built wealth'' \\
E & ``I feel like myself'' & ``I'm authentic'' & ``I found my meaning'' & ``I lived with purpose'' \\
P & ``I'm strong'' & ``I'm healthy'' & ``I take care of myself'' & ``I aged with dignity'' \\
S & ``I belong here'' & ``I'm respected'' & ``I matter to others'' & ``I'll be remembered'' \\
D & ``I'm connected'' & ``I'm tech-savvy'' & ``I'm digitally fluent'' & ``I shaped digital world'' \\
Eco & ``I care about nature'' & ``I act sustainably'' & ``I'm environmentally responsible'' & ``I protected the planet'' \\
\bottomrule
\end{tabular}
\end{table}

%-----------------------------------------------------------------------
\subsection{IDN-Pain Components (24 cells)}
%-----------------------------------------------------------------------

\begin{table}[htbp]
\centering
\caption{IDN-Pain: Identity Utility Losses by Dimension and Time}
\label{tab:app-c-idn-pain}
\small
\renewcommand{\arraystretch}{1.2}
\begin{tabular}{@{}lp{2.8cm}p{2.8cm}p{2.8cm}p{2.8cm}@{}}
\toprule
\textbf{Dim} & \textbf{Instant ($t=0$)} & \textbf{Short ($t=1$)} & \textbf{Medium ($t=2$)} & \textbf{Long ($t=3$)} \\
\midrule
F & ``I'm a failure'' & ``I can't provide'' & ``I'll never be successful'' & ``I wasted my potential'' \\
E & ``I'm not myself'' & ``I've lost my way'' & ``Who am I anymore?'' & ``My life was meaningless'' \\
P & ``I'm weak'' & ``I'm unhealthy'' & ``I've let myself go'' & ``I aged badly'' \\
S & ``I don't belong'' & ``No one respects me'' & ``I don't matter'' & ``I'll be forgotten'' \\
D & ``I'm left behind'' & ``I can't keep up'' & ``Technology passed me by'' & ``I'm obsolete'' \\
Eco & ``I'm destroying nature'' & ``I'm part of the problem'' & ``I failed the planet'' & ``I left a damaged world'' \\
\bottomrule
\end{tabular}
\end{table}

%%%%%%%%%%%%%%%%%%%%%%%%%%%%%%%%%%%%%%%%%%%%%%%%%%%%%%%%%%%%%%%%%%%%%%%%%%%%%%%
%%%%%%%%%%%%%%%%%%%%%%%%%%%%%%%%%%%%%%%%%%%%%%%%%%%%%%%%%%%%%%%%%%%%%%%%%%%%%%%
\section{Loss Aversion by Dimension}
\label{sec:app-c-loss-aversion}
%%%%%%%%%%%%%%%%%%%%%%%%%%%%%%%%%%%%%%%%%%%%%%%%%%%%%%%%%%%%%%%%%%%%%%%%%%%%%%%
%%%%%%%%%%%%%%%%%%%%%%%%%%%%%%%%%%%%%%%%%%%%%%%%%%%%%%%%%%%%%%%%%%%%%%%%%%%%%%%

\begin{cross-reference}
\textbf{Chapter Reference:} See Chapter~\ref{sec:welfare-fepsde}, 
Section~\ref{subsec:10-4-calculation} for loss aversion treatment.
\end{cross-reference}

\begin{table}[htbp]
\centering
\caption{Loss Aversion Coefficients by FEPSDE Dimension}
\label{tab:app-c-lambda}
\renewcommand{\arraystretch}{1.4}
\begin{tabular}{@{}clcl@{}}
\toprule
\textbf{Dim} & \textbf{Dimension} & \textbf{$\lambda_d$} & \textbf{Interpretation} \\
\midrule
F & Financial & $2.0$ & Kahneman-Tversky baseline \\
E & Emotional & $1.5$--$2.5$ & Varies with affect intensity \\
P & Physical & $2.0$--$3.0$ & Health losses weighted heavily \\
S & Social & $2.0$--$4.0$ & Status losses especially painful \\
D & Digital & $1.5$--$2.0$ & Moderate, newer dimension \\
Eco & Ecological & $1.2$--$2.0$ & Often underweighted (temporal distance) \\
\bottomrule
\end{tabular}
\begin{flushleft}
\small\textit{Note:} Values from LLM-MC estimation (see Appendix~AN). 
Ranges indicate heterogeneity across populations.
\end{flushleft}
\end{table}

%%%%%%%%%%%%%%%%%%%%%%%%%%%%%%%%%%%%%%%%%%%%%%%%%%%%%%%%%%%%%%%%%%%%%%%%%%%%%%%
%%%%%%%%%%%%%%%%%%%%%%%%%%%%%%%%%%%%%%%%%%%%%%%%%%%%%%%%%%%%%%%%%%%%%%%%%%%%%%%
\section{The $\Gamma$-Matrix: Context-Welfare Interaction}
\label{sec:app-c-gamma}
%%%%%%%%%%%%%%%%%%%%%%%%%%%%%%%%%%%%%%%%%%%%%%%%%%%%%%%%%%%%%%%%%%%%%%%%%%%%%%%
%%%%%%%%%%%%%%%%%%%%%%%%%%%%%%%%%%%%%%%%%%%%%%%%%%%%%%%%%%%%%%%%%%%%%%%%%%%%%%%

\begin{cross-reference}
\textbf{Chapter Reference:} See Chapter~10, Section~10.7 for conceptual treatment.
See Chapter~9 for the 8$\Psi$-Dimensions.
\end{cross-reference}

%-----------------------------------------------------------------------
\subsection{The Context-Welfare Connection}
%-----------------------------------------------------------------------

\paragraph{The Core Claim.}
Welfare is not produced in a vacuum. The same action, the same outcome, 
the same person can experience vastly different utility depending on context:

\begin{itemize}
\item A \$1,000 bonus feels different in a recession vs. a boom
\item Social rejection hurts more in a low-trust society
\item Health improvements matter more where healthcare is scarce
\end{itemize}

\paragraph{The $\Psi$-Dimensions (from Chapter 9).}
Context is captured by eight dimensions:

\begin{table}[htbp]
\centering
\caption{The Eight $\Psi$-Dimensions}
\label{tab:app-c-psi-dimensions}
\renewcommand{\arraystretch}{1.3}
\begin{tabular}{@{}lll@{}}
\toprule
\textbf{Symbol} & \textbf{Dimension} & \textbf{Content} \\
\midrule
$\Psi_I$ & Institutional & Rule of law, property rights, governance quality \\
$\Psi_S$ & Social & Trust, social cohesion, cooperation norms \\
$\Psi_C$ & Cognitive & Education, skills, human capital \\
$\Psi_K$ & Information & Knowledge access, transparency, media quality \\
$\Psi_E$ & Economic & Development level, income, market maturity \\
$\Psi_T$ & Temporal & Stability, predictability, time horizons \\
$\Psi_M$ & Market & Market scope, competition, integration \\
$\Psi_F$ & Flexibility & Adaptability, mobility, openness to change \\
\bottomrule
\end{tabular}
\end{table}

%-----------------------------------------------------------------------
\subsection{Definition: The $\Gamma$-Matrix Structure}
%-----------------------------------------------------------------------

\begin{tcolorbox}[colback=blue!5!white, colframe=blue!50!black, title=Definition: $\Gamma$-Matrix]
\begin{equation}
\boxed{
U_d(\Psi) = \bar{U}_d + \sum_{j=1}^{8} \gamma_{dj} \cdot \Psi_j
}
\label{eq:app-c-gamma-definition}
\end{equation}
where:
\begin{itemize}
\item $d \in \{F, E, P, S, D, Eco\}$: welfare dimension (6 rows)
\item $j \in \{I, S, C, K, E, T, M, F\}$: context dimension (8 columns)
\item $\gamma_{dj}$: effect of context dimension $j$ on welfare dimension $d$
\item $\bar{U}_d$: baseline welfare in dimension $d$
\end{itemize}
\end{tcolorbox}

\paragraph{Matrix Dimensions.}
The $\Gamma$-matrix is $6 \times 8$:
\begin{itemize}
\item 6 rows: FEPSDE welfare dimensions
\item 8 columns: $\Psi$ context dimensions
\item 48 coefficients: $\gamma_{dj}$ for each combination
\end{itemize}

%-----------------------------------------------------------------------
\subsection{The Complete $\Gamma$-Matrix}
%-----------------------------------------------------------------------

\begin{table}[htbp]
\centering
\caption{The $\Gamma$-Matrix: Context-Welfare Interaction Coefficients}
\label{tab:app-c-gamma-full}
\small
\renewcommand{\arraystretch}{1.4}
\begin{tabular}{@{}lcccccccc@{}}
\toprule
& $\Psi_I$ & $\Psi_S$ & $\Psi_C$ & $\Psi_K$ & $\Psi_E$ & $\Psi_T$ & $\Psi_M$ & $\Psi_F$ \\
& \scriptsize{Inst.} & \scriptsize{Social} & \scriptsize{Cogn.} & \scriptsize{Info.} & \scriptsize{Econ.} & \scriptsize{Temp.} & \scriptsize{Market} & \scriptsize{Flex.} \\
\midrule
$U^F$ & 0.31 & 0.12 & 0.18 & 0.25 & \textbf{0.42} & 0.14 & 0.28 & 0.33 \\
\scriptsize{Financial} & \scriptsize{(.05)} & \scriptsize{(.04)} & \scriptsize{(.05)} & \scriptsize{(.04)} & \scriptsize{(.06)} & \scriptsize{(.04)} & \scriptsize{(.05)} & \scriptsize{(.05)} \\
\addlinespace
$U^E$ & 0.15 & \textbf{0.47} & 0.11 & 0.09 & 0.08 & 0.29 & 0.02 & 0.06 \\
\scriptsize{Emotional} & \scriptsize{(.04)} & \scriptsize{(.06)} & \scriptsize{(.04)} & \scriptsize{(.03)} & \scriptsize{(.03)} & \scriptsize{(.05)} & \scriptsize{(.02)} & \scriptsize{(.03)} \\
\addlinespace
$U^P$ & 0.28 & 0.14 & \textbf{0.31} & 0.16 & 0.24 & 0.26 & 0.03 & 0.09 \\
\scriptsize{Physical} & \scriptsize{(.05)} & \scriptsize{(.04)} & \scriptsize{(.05)} & \scriptsize{(.04)} & \scriptsize{(.05)} & \scriptsize{(.05)} & \scriptsize{(.02)} & \scriptsize{(.03)} \\
\addlinespace
$U^S$ & 0.11 & \textbf{0.51} & 0.04 & 0.08 & 0.01 & 0.12 & \textbf{--0.18} & 0.02 \\
\scriptsize{Social} & \scriptsize{(.04)} & \scriptsize{(.06)} & \scriptsize{(.03)} & \scriptsize{(.03)} & \scriptsize{(.02)} & \scriptsize{(.04)} & \scriptsize{(.04)} & \scriptsize{(.02)} \\
\addlinespace
$U^D$ & 0.13 & 0.03 & 0.27 & \textbf{0.39} & 0.22 & 0.08 & 0.24 & 0.11 \\
\scriptsize{Digital} & \scriptsize{(.04)} & \scriptsize{(.02)} & \scriptsize{(.05)} & \scriptsize{(.05)} & \scriptsize{(.04)} & \scriptsize{(.03)} & \scriptsize{(.05)} & \scriptsize{(.04)} \\
\addlinespace
$U^{Eco}$ & 0.21 & 0.09 & 0.12 & 0.14 & \textbf{--0.19} & 0.23 & \textbf{--0.21} & 0.04 \\
\scriptsize{Ecological} & \scriptsize{(.04)} & \scriptsize{(.03)} & \scriptsize{(.04)} & \scriptsize{(.04)} & \scriptsize{(.04)} & \scriptsize{(.05)} & \scriptsize{(.04)} & \scriptsize{(.02)} \\
\bottomrule
\end{tabular}
\begin{flushleft}
\small\textit{Note:} LLM-MC estimates with standard errors in parentheses. 
Bold indicates largest magnitude per row. Negative values indicate structural trade-offs.
\end{flushleft}
\end{table}

%-----------------------------------------------------------------------
\subsection{Row Reading: What Drives Each Welfare Dimension?}
%-----------------------------------------------------------------------

\paragraph{Financial Welfare ($U^F$).}
\begin{itemize}
\item Strongest driver: $\gamma_{F,E} = 0.42$ (Economic development)
\item Also important: $\gamma_{F,F} = 0.33$ (Flexibility), $\gamma_{F,I} = 0.31$ (Institutions)
\item Interpretation: Financial welfare requires development, flexibility, and institutions
\end{itemize}

\paragraph{Emotional Welfare ($U^E$).}
\begin{itemize}
\item Strongest driver: $\gamma_{E,S} = 0.47$ (Social trust)
\item Also important: $\gamma_{E,T} = 0.29$ (Temporal stability)
\item Interpretation: Happiness depends on trust and stability, not just money
\end{itemize}

\paragraph{Physical Welfare ($U^P$).}
\begin{itemize}
\item Strongest driver: $\gamma_{P,C} = 0.31$ (Cognitive capacity/education)
\item Also important: $\gamma_{P,I} = 0.28$ (Institutions), $\gamma_{P,T} = 0.26$ (Stability)
\item Interpretation: Health requires education and stable healthcare systems
\end{itemize}

\paragraph{Social Welfare ($U^S$).}
\begin{itemize}
\item Strongest driver: $\gamma_{S,S} = 0.51$ (Social trust---self-reinforcing!)
\item \textbf{Negative:} $\gamma_{S,M} = -0.18$ (Market integration harms social welfare)
\item Interpretation: Trust builds trust; markets can erode social fabric
\end{itemize}

\paragraph{Digital Welfare ($U^D$).}
\begin{itemize}
\item Strongest driver: $\gamma_{D,K} = 0.39$ (Information access)
\item Also important: $\gamma_{D,C} = 0.27$ (Cognitive capacity)
\item Interpretation: Digital welfare requires information infrastructure and skills
\end{itemize}

\paragraph{Ecological Welfare ($U^{Eco}$).}
\begin{itemize}
\item \textbf{Negative:} $\gamma_{Eco,M} = -0.21$ (Market integration harms ecology)
\item \textbf{Negative:} $\gamma_{Eco,E} = -0.19$ (Economic development harms ecology)
\item Positive: $\gamma_{Eco,T} = 0.23$ (Stability enables environmental investment)
\item Interpretation: Growth-environment trade-off is real
\end{itemize}

%-----------------------------------------------------------------------
\subsection{Column Reading: What Does Each Context Dimension Affect?}
%-----------------------------------------------------------------------

\paragraph{Institutional Quality ($\Psi_I$).}
\begin{itemize}
\item Strongly affects: $U^F$ (0.31), $U^P$ (0.28), $U^{Eco}$ (0.21)
\item Interpretation: Institutions matter for material and health outcomes
\end{itemize}

\paragraph{Social Trust ($\Psi_S$).}
\begin{itemize}
\item Strongly affects: $U^S$ (0.51), $U^E$ (0.47)
\item Weakly affects: $U^D$ (0.03), $U^{Eco}$ (0.09)
\item Interpretation: Trust matters most for social and emotional welfare
\end{itemize}

\paragraph{Cognitive Capacity ($\Psi_C$).}
\begin{itemize}
\item Strongly affects: $U^P$ (0.31), $U^D$ (0.27)
\item Interpretation: Education drives health and digital welfare
\end{itemize}

\paragraph{Information Access ($\Psi_K$).}
\begin{itemize}
\item Strongly affects: $U^D$ (0.39), $U^F$ (0.25)
\item Interpretation: Information infrastructure matters for digital and financial outcomes
\end{itemize}

\paragraph{Economic Development ($\Psi_E$).}
\begin{itemize}
\item Strongly affects: $U^F$ (0.42), $U^P$ (0.24), $U^D$ (0.22)
\item \textbf{Negative:} $U^{Eco}$ (--0.19)
\item Interpretation: Development helps material welfare but harms environment
\end{itemize}

\paragraph{Temporal Stability ($\Psi_T$).}
\begin{itemize}
\item Broad positive effects: $U^E$ (0.29), $U^P$ (0.26), $U^{Eco}$ (0.23)
\item Interpretation: Stability enables long-term investment across dimensions
\end{itemize}

\paragraph{Market Integration ($\Psi_M$).}
\begin{itemize}
\item Positive effects: $U^F$ (0.28), $U^D$ (0.24)
\item \textbf{Negative effects:} $U^S$ (--0.18), $U^{Eco}$ (--0.21)
\item Interpretation: Markets boost financial/digital but harm social/ecological
\end{itemize}

\paragraph{Flexibility ($\Psi_F$).}
\begin{itemize}
\item Strongest effect: $U^F$ (0.33)
\item Weak effects elsewhere
\item Interpretation: Flexibility primarily matters for financial adaptation
\end{itemize}

%-----------------------------------------------------------------------
\subsection{Structural Trade-offs Revealed}
%-----------------------------------------------------------------------

\paragraph{The Three Negative Coefficients.}
The $\Gamma$-matrix reveals three structural trade-offs:

\begin{table}[htbp]
\centering
\caption{Structural Trade-offs in the $\Gamma$-Matrix}
\label{tab:app-c-tradeoffs}
\renewcommand{\arraystretch}{1.4}
\begin{tabular}{@{}llll@{}}
\toprule
\textbf{Coefficient} & \textbf{Value} & \textbf{Trade-off} & \textbf{Literature} \\
\midrule
$\gamma_{S,M}$ & --0.18 & Markets $\uparrow$ $\Rightarrow$ Social welfare $\downarrow$ & Autor et al. (2013) \\
$\gamma_{Eco,M}$ & --0.21 & Markets $\uparrow$ $\Rightarrow$ Ecological welfare $\downarrow$ & Globalization studies \\
$\gamma_{Eco,E}$ & --0.19 & Development $\uparrow$ $\Rightarrow$ Ecological welfare $\downarrow$ & EKC literature \\
\bottomrule
\end{tabular}
\end{table}

\paragraph{These Are Not Estimation Errors.}
The negative coefficients reflect \emph{genuine tensions} documented in the literature:

\begin{itemize}
\item \textbf{$\gamma_{S,M} = -0.18$:} Autor et al. (2013) documented how trade 
shocks from China devastated American communities---job losses, family breakdown, 
opioid epidemics. Market integration increased aggregate GDP but destroyed 
local social fabric.

\item \textbf{$\gamma_{Eco,M} = -0.21$:} Globalization enables ``pollution havens''---
production moves to jurisdictions with weak environmental regulation. 
Global supply chains increase carbon footprints.

\item \textbf{$\gamma_{Eco,E} = -0.19$:} The Environmental Kuznets Curve literature 
documents how development initially increases environmental degradation. 
The inverted-U only emerges at high income levels, and even then, only for 
some pollutants.
\end{itemize}

\paragraph{Policy Implication.}
These trade-offs mean that policies maximizing one welfare dimension may 
harm others. Optimal policy requires explicit trade-off analysis, not 
single-dimension optimization.

%-----------------------------------------------------------------------
\subsection{Context Effects on Behavioral Parameters}
%-----------------------------------------------------------------------

\paragraph{Beyond Welfare Levels.}
The $\Gamma$-matrix captures context effects on welfare \emph{levels}.
But context also affects \emph{behavioral parameters}---the deep coefficients 
governing how outcomes translate to utility.

\paragraph{Three Context-Dependent Parameters.}

\begin{table}[htbp]
\centering
\caption{Context-Dependent Behavioral Parameters}
\label{tab:app-c-context-parameters}
\renewcommand{\arraystretch}{1.4}
\begin{tabular}{@{}llll@{}}
\toprule
\textbf{Parameter} & \textbf{Context Effect} & \textbf{Coefficient} & \textbf{Interpretation} \\
\midrule
Discount $\delta$ & $\delta = \bar{\delta} + \gamma_{\delta,T} \cdot \Psi_T$ & $\gamma_{\delta,T} = 0.08$ & Stability $\uparrow$ $\Rightarrow$ Patience $\uparrow$ \\
Loss aversion $\lambda$ & $\lambda = \bar{\lambda} - \gamma_{\lambda,I} \cdot \Psi_I$ & $\gamma_{\lambda,I} = 0.15$ & Institutions $\uparrow$ $\Rightarrow$ Loss aversion $\downarrow$ \\
Collective weight $\omega$ & $\omega = \bar{\omega} + \gamma_{\omega,S} \cdot \Psi_S$ & $\gamma_{\omega,S} = 0.12$ & Trust $\uparrow$ $\Rightarrow$ Collective concern $\uparrow$ \\
\bottomrule
\end{tabular}
\end{table}

\paragraph{Interpretation.}

\textbf{Discount Factor ($\delta$):}
People in stable environments ($\Psi_T$ high) are more patient:
\begin{itemize}
\item Stable institutions make future payoffs more certain
\item Long time horizons encourage investment over consumption
\item Empirically: Patience correlates with institutional quality across countries
\end{itemize}

\textbf{Loss Aversion ($\lambda$):}
People in well-functioning institutions ($\Psi_I$ high) are less loss-averse:
\begin{itemize}
\item Good institutions provide insurance against losses
\item Property rights make recovery from losses more feasible
\item Empirically: Risk tolerance correlates with institutional quality
\end{itemize}

\textbf{Collective Weight ($\omega_{KNU}$):}
People in high-trust societies ($\Psi_S$ high) weight collective welfare more:
\begin{itemize}
\item Trust makes collective action credible
\item Social capital enables group-level thinking
\item Empirically: Prosociality correlates with social trust
\end{itemize}

%-----------------------------------------------------------------------
\subsection{Country Comparison Example: Denmark vs. Brazil}
%-----------------------------------------------------------------------

\paragraph{Different $\Psi$-Profiles.}
Consider two countries with different context profiles:

\begin{table}[htbp]
\centering
\caption{$\Psi$-Profiles: Denmark vs. Brazil}
\label{tab:app-c-denmark-brazil-psi}
\renewcommand{\arraystretch}{1.3}
\begin{tabular}{@{}lcc@{}}
\toprule
\textbf{$\Psi$-Dimension} & \textbf{Denmark} & \textbf{Brazil} \\
\midrule
$\Psi_I$ (Institutional) & 0.85 & 0.45 \\
$\Psi_S$ (Social trust) & 0.90 & 0.35 \\
$\Psi_C$ (Cognitive) & 0.80 & 0.50 \\
$\Psi_K$ (Information) & 0.85 & 0.55 \\
$\Psi_E$ (Economic) & 0.75 & 0.55 \\
$\Psi_T$ (Temporal) & 0.85 & 0.40 \\
$\Psi_M$ (Market) & 0.80 & 0.60 \\
$\Psi_F$ (Flexibility) & 0.75 & 0.50 \\
\bottomrule
\end{tabular}
\end{table}

\paragraph{Predicted Welfare Differences.}
Using the $\Gamma$-matrix:

\begin{table}[htbp]
\centering
\caption{Predicted Welfare: Denmark vs. Brazil}
\label{tab:app-c-denmark-brazil-welfare}
\renewcommand{\arraystretch}{1.3}
\begin{tabular}{@{}lccc@{}}
\toprule
\textbf{Welfare Dimension} & \textbf{Denmark} & \textbf{Brazil} & \textbf{$\Delta$} \\
\midrule
$U^F$ (Financial) & 0.72 & 0.48 & +0.24 \\
$U^E$ (Emotional) & \textbf{0.78} & 0.42 & \textbf{+0.36} \\
$U^P$ (Physical) & 0.69 & 0.45 & +0.24 \\
$U^S$ (Social) & \textbf{0.75} & 0.38 & \textbf{+0.37} \\
$U^D$ (Digital) & 0.71 & 0.52 & +0.19 \\
$U^{Eco}$ (Ecological) & 0.58 & 0.42 & +0.16 \\
\bottomrule
\end{tabular}
\end{table}

\paragraph{Key Insight.}
The largest differences are in Emotional ($+0.36$) and Social ($+0.37$) welfare---driven 
primarily by Denmark's higher $\Psi_S$ (trust). Brazil's lower institutional quality 
and trust create welfare deficits that cannot be fully compensated by other dimensions.

%-----------------------------------------------------------------------
\subsection{The $\Gamma$-Matrix in EBF}
%-----------------------------------------------------------------------

\paragraph{Integration with Utility Categories.}
The $\Gamma$-matrix applies to all three utility categories:

\begin{equation}
U^i_d(\Psi) = \bar{U}^i_d + \sum_{j} \gamma_{dj} \cdot \Psi_j \quad \text{for } i \in \{INU, KNU, IDN\}
\end{equation}

The same context affects INU, KNU, and IDN welfare in each dimension.

\paragraph{Category-Specific Sensitivity.}
However, categories may have different \emph{sensitivity} to context:

\begin{tcolorbox}[colback=green!5!white, colframe=green!50!black]
\begin{equation}
\boxed{
U^i_d(\Psi) = \bar{U}^i_d + \kappa^i \cdot \sum_{j} \gamma_{dj} \cdot \Psi_j
}
\label{eq:app-c-kappa-scaling}
\end{equation}
\end{tcolorbox}

where $\kappa^i$ is a category-specific scaling factor:

\begin{table}[htbp]
\centering
\caption{Category-Specific Context Sensitivity}
\label{tab:app-c-kappa-factors}
\renewcommand{\arraystretch}{1.3}
\begin{tabular}{@{}lcl@{}}
\toprule
\textbf{Category} & $\kappa^i$ & \textbf{Interpretation} \\
\midrule
INU & $\approx 1.0$ & INU responds normally to context \\
KNU & $\approx 1.2$ & KNU is more context-sensitive (group welfare depends more on environment) \\
IDN & $\approx 0.8$ & IDN is less context-sensitive (identity is more internal) \\
\bottomrule
\end{tabular}
\end{table}

\paragraph{Full Integration.}
The complete context-modulated utility for each component is:

\begin{tcolorbox}[colback=blue!5!white, colframe=blue!50!black]
\begin{equation}
\boxed{
U^i_{d,t}(\Psi) = (\delta^i_d(\Psi))^t \cdot \left[ G^i_{d,t}(\Psi) - \lambda^i_d(\Psi) \cdot P^i_{d,t}(\Psi) \right]
}
\label{eq:app-c-full-context-integration}
\end{equation}
where all parameters ($\delta$, $\lambda$, $G$, $P$) are functions of context $\Psi$.
\end{tcolorbox}

%-----------------------------------------------------------------------
\subsection{Validation: $\Gamma$-Matrix vs. Empirical Literature}
%-----------------------------------------------------------------------

\paragraph{The Estimation Challenge.}
Directly estimating 48 $\gamma_{dj}$ coefficients would require:
\begin{itemize}
\item Cross-country panel data with consistent FEPSDE measures
\item Variation in all eight $\Psi$ dimensions
\item Instruments for causal identification
\end{itemize}

This is beyond most research budgets.

\paragraph{The LLM-MC Solution.}
Following the methodology in Appendix~N, we used LLM Monte Carlo to extract 
the $\Gamma$ coefficients from the implicit knowledge in scientific literature.

\paragraph{Validation Against Empirical Benchmarks.}

\begin{table}[htbp]
\centering
\caption{$\Gamma$-Matrix Validation}
\label{tab:app-c-gamma-validation}
\renewcommand{\arraystretch}{1.3}
\begin{tabular}{@{}llccc@{}}
\toprule
\textbf{Coefficient} & \textbf{Empirical Source} & \textbf{Empirical} & \textbf{LLM-MC} & \textbf{$|\Delta|$} \\
\midrule
$\gamma_{F,I}$ & Acemoglu et al. (2001) & 0.28 & 0.31 & 0.03 \\
$\gamma_{E,S}$ & Helliwell \& Putnam (2004) & 0.44 & 0.47 & 0.03 \\
$\gamma_{P,C}$ & Cutler \& Lleras-Muney (2010) & 0.35 & 0.31 & 0.04 \\
$\gamma_{S,M}$ & Autor et al. (2013) & --0.21 & --0.18 & 0.03 \\
$\gamma_{Eco,E}$ & EKC meta-analyses & --0.15 & --0.19 & 0.04 \\
\midrule
\multicolumn{4}{r}{\textbf{Mean Absolute Error:}} & \textbf{0.034} \\
\multicolumn{4}{r}{\textbf{Sign Agreement:}} & \textbf{5/5 (100\%)} \\
\bottomrule
\end{tabular}
\end{table}

The LLM-MC estimates are close to empirical benchmarks (MAE = 0.034) and 
correctly predict all signs---including the negative trade-off coefficients.

%-----------------------------------------------------------------------
\subsection{Summary: The $\Gamma$-Matrix}
%-----------------------------------------------------------------------

\begin{tcolorbox}[colback=gray!5!white, colframe=gray!75!black, title=Section C.7 Summary]
\textbf{The $\Gamma$-Matrix:}
$$U_d(\Psi) = \bar{U}_d + \sum_{j=1}^{8} \gamma_{dj} \cdot \Psi_j$$

\textbf{Structure:}
\begin{itemize}
\item 6 $\times$ 8 matrix (48 coefficients)
\item Rows: FEPSDE welfare dimensions
\item Columns: $\Psi$ context dimensions
\end{itemize}

\textbf{Key Findings:}
\begin{itemize}
\item Strongest positive: $\gamma_{S,S} = 0.51$ (trust $\rightarrow$ social welfare)
\item Strongest positive: $\gamma_{E,S} = 0.47$ (trust $\rightarrow$ emotional welfare)
\item Negative trade-offs: $\gamma_{S,M} = -0.18$, $\gamma_{Eco,M} = -0.21$, $\gamma_{Eco,E} = -0.19$
\end{itemize}

\textbf{Structural Trade-offs:}
\begin{itemize}
\item Market integration harms social and ecological welfare
\item Economic development harms ecological welfare
\item These are genuine tensions, not estimation errors
\end{itemize}

\textbf{Context-Dependent Parameters:}
\begin{itemize}
\item $\delta(\Psi_T)$: Stability increases patience ($\gamma_{\delta,T} = 0.08$)
\item $\lambda(\Psi_I)$: Institutions reduce loss aversion ($\gamma_{\lambda,I} = 0.15$)
\item $\omega_{KNU}(\Psi_S)$: Trust increases collective concern ($\gamma_{\omega,S} = 0.12$)
\end{itemize}

\textbf{Category-Specific Sensitivity:}
\begin{itemize}
\item $\kappa^{INU} \approx 1.0$, $\kappa^{KNU} \approx 1.2$, $\kappa^{IDN} \approx 0.8$
\end{itemize}

\textbf{Validation:}
\begin{itemize}
\item MAE = 0.034 against empirical benchmarks
\item 100\% sign agreement including negative coefficients
\end{itemize}
\end{tcolorbox}

%%%%%%%%%%%%%%%%%%%%%%%%%%%%%%%%%%%%%%%%%%%%%%%%%%%%%%%%%%%%%%%%%%%%%%%%%%%%%%%
%%%%%%%%%%%%%%%%%%%%%%%%%%%%%%%%%%%%%%%%%%%%%%%%%%%%%%%%%%%%%%%%%%%%%%%%%%%%%%%
\section{Context-Dependent Parameters}
\label{sec:app-c-context-params}
%%%%%%%%%%%%%%%%%%%%%%%%%%%%%%%%%%%%%%%%%%%%%%%%%%%%%%%%%%%%%%%%%%%%%%%%%%%%%%%
%%%%%%%%%%%%%%%%%%%%%%%%%%%%%%%%%%%%%%%%%%%%%%%%%%%%%%%%%%%%%%%%%%%%%%%%%%%%%%%

\begin{table}[htbp]
\centering
\caption{Context-Dependent Behavioral Parameters}
\label{tab:app-c-context-params}
\renewcommand{\arraystretch}{1.4}
\begin{tabular}{@{}llll@{}}
\toprule
\textbf{Parameter} & \textbf{Context Effect} & \textbf{Coefficient} & \textbf{Interpretation} \\
\midrule
Discount $\delta$ & $\delta = \bar{\delta} + \gamma_{\delta,T} \cdot \Psi_T$ & $\gamma_{\delta,T} = 0.08$ & Stability $\uparrow$ $\Rightarrow$ Patience $\uparrow$ \\
Loss aversion $\lambda$ & $\lambda = \bar{\lambda} - \gamma_{\lambda,I} \cdot \Psi_I$ & $\gamma_{\lambda,I} = 0.15$ & Institutions $\uparrow$ $\Rightarrow$ Loss aversion $\downarrow$ \\
Collective weight $\omega$ & $\omega = \bar{\omega} + \gamma_{\omega,S} \cdot \Psi_S$ & $\gamma_{\omega,S} = 0.12$ & Trust $\uparrow$ $\Rightarrow$ Collective concern $\uparrow$ \\
\bottomrule
\end{tabular}
\end{table}

%%%%%%%%%%%%%%%%%%%%%%%%%%%%%%%%%%%%%%%%%%%%%%%%%%%%%%%%%%%%%%%%%%%%%%%%%%%%%%%
%%%%%%%%%%%%%%%%%%%%%%%%%%%%%%%%%%%%%%%%%%%%%%%%%%%%%%%%%%%%%%%%%%%%%%%%%%%%%%%
\section{Inter-Category Complementarities}
\label{sec:app-c-inter-category}
%%%%%%%%%%%%%%%%%%%%%%%%%%%%%%%%%%%%%%%%%%%%%%%%%%%%%%%%%%%%%%%%%%%%%%%%%%%%%%%
%%%%%%%%%%%%%%%%%%%%%%%%%%%%%%%%%%%%%%%%%%%%%%%%%%%%%%%%%%%%%%%%%%%%%%%%%%%%%%%

\begin{cross-reference}
\textbf{Chapter Reference:} See Chapter~10, Section~10.6.
\end{cross-reference}

%-----------------------------------------------------------------------
\subsection{The Fundamental Insight: Utility Categories Interact}
%-----------------------------------------------------------------------

\paragraph{Beyond Additive Aggregation.}
A naive model would simply add the three utilities:
\begin{equation}
Q_{naive} = \omega_{INU} \cdot U^{INU} + \omega_{KNU} \cdot U^{KNU} + \omega_{IDN} \cdot U^{IDN}
\end{equation}

But this misses crucial interactions:
\begin{itemize}
\item Does individual success enhance or threaten collective welfare?
\item Does personal gain confirm or violate identity?
\item Does group success strengthen or weaken individual standing?
\end{itemize}

\paragraph{The Complementarity Extension.}
The full model includes interaction terms:

\begin{tcolorbox}[colback=blue!5!white, colframe=blue!50!black]
\begin{equation}
\boxed{
Q = \sum_{i} \omega_i \cdot U^i + \sum_{i < j} \gamma_{ij} \cdot U^i \cdot U^j
}
\label{eq:app-c-complementarity-model}
\end{equation}
where $\gamma_{ij}$ captures the complementarity between categories $i$ and $j$.
\end{tcolorbox}

%-----------------------------------------------------------------------
\subsection{The Three Complementarity Coefficients}
%-----------------------------------------------------------------------

\paragraph{The Complementarity Matrix.}
With three categories, there are three pairwise interactions:

\begin{table}[htbp]
\centering
\caption{The Inter-Category Complementarity Matrix}
\label{tab:app-c-complementarity-matrix}
\renewcommand{\arraystretch}{1.4}
\begin{tabular}{@{}l|ccc@{}}
\toprule
& INU & KNU & IDN \\
\midrule
INU & --- & $\gamma_{INU,KNU}$ & $\gamma_{INU,IDN}$ \\
KNU & $\gamma_{KNU,INU}$ & --- & $\gamma_{KNU,IDN}$ \\
IDN & $\gamma_{IDN,INU}$ & $\gamma_{IDN,KNU}$ & --- \\
\bottomrule
\end{tabular}
\end{table}

\textbf{Note:} The matrix is symmetric ($\gamma_{ij} = \gamma_{ji}$), so there are 
three unique coefficients.

\paragraph{Interpretation of $\gamma_{ij}$.}

\begin{table}[htbp]
\centering
\caption{Interpreting Complementarity Coefficients}
\label{tab:app-c-gamma-interpretation}
\renewcommand{\arraystretch}{1.3}
\begin{tabular}{@{}lll@{}}
\toprule
\textbf{Sign} & \textbf{Name} & \textbf{Meaning} \\
\midrule
$\gamma_{ij} > 0$ & Complement & High $U^i$ \emph{enhances} value of high $U^j$ \\
$\gamma_{ij} < 0$ & Substitute & High $U^i$ \emph{reduces} value of high $U^j$ \\
$\gamma_{ij} = 0$ & Independent & $U^i$ and $U^j$ do not interact \\
\bottomrule
\end{tabular}
\end{table}

%-----------------------------------------------------------------------
\subsection{$\gamma_{INU,KNU}$: Individual--Collective Complementarity}
%-----------------------------------------------------------------------

\paragraph{The Question.}
Does individual gain enhance or diminish collective welfare---and vice versa?

\paragraph{Positive Complementarity ($\gamma_{INU,KNU} > 0$).}
Individual and collective interests align:

\begin{table}[htbp]
\centering
\caption{Positive INU-KNU Complementarity Examples}
\label{tab:app-c-inu-knu-positive}
\renewcommand{\arraystretch}{1.3}
\begin{tabular}{@{}lll@{}}
\toprule
\textbf{Situation} & \textbf{Mechanism} & \textbf{Effect} \\
\midrule
Team bonus & My success = team success & Both rise together \\
Cooperative venture & Individual effort benefits group & Mutual reinforcement \\
Family prosperity & My income supports family & Aligned interests \\
Public goods (high trust) & My contribution valued & Individual and collective gain \\
\bottomrule
\end{tabular}
\end{table}

When $\gamma_{INU,KNU} > 0$:
\begin{equation}
\frac{\partial^2 Q}{\partial U^{INU} \partial U^{KNU}} > 0
\end{equation}
Gains in one category make gains in the other \emph{more valuable}.

\paragraph{Negative Complementarity ($\gamma_{INU,KNU} < 0$).}
Individual and collective interests conflict:

\begin{table}[htbp]
\centering
\caption{Negative INU-KNU Complementarity Examples}
\label{tab:app-c-inu-knu-negative}
\renewcommand{\arraystretch}{1.3}
\begin{tabular}{@{}lll@{}}
\toprule
\textbf{Situation} & \textbf{Mechanism} & \textbf{Effect} \\
\midrule
Zero-sum competition & My gain = your loss & INU$\uparrow$ implies KNU$\downarrow$ \\
Free-riding & I benefit, group pays & Individual gain, collective harm \\
Corruption & Personal enrichment & Individual gain undermines collective \\
Exploitation & Extract from group & Negative-sum interaction \\
\bottomrule
\end{tabular}
\end{table}

When $\gamma_{INU,KNU} < 0$:
\begin{equation}
\frac{\partial^2 Q}{\partial U^{INU} \partial U^{KNU}} < 0
\end{equation}
High individual utility \emph{reduces} the marginal value of collective utility.

\paragraph{Empirical Estimates by Context.}

\begin{table}[htbp]
\centering
\caption{$\gamma_{INU,KNU}$ by Context}
\label{tab:app-c-gamma-inu-knu-context}
\renewcommand{\arraystretch}{1.3}
\begin{tabular}{@{}llc@{}}
\toprule
\textbf{Context} & \textbf{Characteristic} & $\gamma_{INU,KNU}$ \\
\midrule
High-trust cooperative & Aligned incentives & $+0.3$ to $+0.5$ \\
Family/close group & Strong bonds & $+0.2$ to $+0.4$ \\
Competitive market & Zero-sum elements & $-0.1$ to $+0.1$ \\
Corrupt institution & Extraction dominates & $-0.2$ to $-0.4$ \\
Prisoner's dilemma & Structural conflict & $-0.3$ to $-0.5$ \\
\bottomrule
\end{tabular}
\end{table}

%-----------------------------------------------------------------------
\subsection{$\gamma_{INU,IDN}$: Individual--Identity Complementarity}
%-----------------------------------------------------------------------

\paragraph{The Question.}
Does personal success confirm or violate one's identity?

\paragraph{Positive Complementarity ($\gamma_{INU,IDN} > 0$).}
Personal success aligns with identity:

\begin{table}[htbp]
\centering
\caption{Positive INU-IDN Complementarity Examples}
\label{tab:app-c-inu-idn-positive}
\renewcommand{\arraystretch}{1.3}
\begin{tabular}{@{}lll@{}}
\toprule
\textbf{Situation} & \textbf{Identity} & \textbf{Effect} \\
\midrule
Entrepreneur succeeds & ``I'm a winner'' & Success confirms identity \\
Artist sells work & ``I'm a real artist'' & Market validates identity \\
Athlete wins & ``I'm a champion'' & Victory confirms self-image \\
Scholar publishes & ``I'm an intellectual'' & Achievement fits narrative \\
\bottomrule
\end{tabular}
\end{table}

When $\gamma_{INU,IDN} > 0$, success feels \emph{more} satisfying because it 
confirms who you are.

\paragraph{Negative Complementarity ($\gamma_{INU,IDN} < 0$).}
Personal success conflicts with identity:

\begin{table}[htbp]
\centering
\caption{Negative INU-IDN Complementarity Examples}
\label{tab:app-c-inu-idn-negative}
\renewcommand{\arraystretch}{1.3}
\begin{tabular}{@{}lll@{}}
\toprule
\textbf{Situation} & \textbf{Identity Conflict} & \textbf{Effect} \\
\midrule
``Selling out'' & Artist takes corporate job & INU$\uparrow$ but IDN$\downarrow\downarrow$ \\
Betraying class & Working-class person ``goes posh'' & Financial gain, identity loss \\
Coal miner $\to$ coder & Manual labor identity & Higher pay, identity violation \\
Inherited wealth & ``Self-made'' identity & Wealth doesn't fit narrative \\
\bottomrule
\end{tabular}
\end{table}

When $\gamma_{INU,IDN} < 0$, success feels \emph{hollow} or even \emph{shameful}
because it violates who you are.

\paragraph{The ``Selling Out'' Phenomenon.}
Negative $\gamma_{INU,IDN}$ explains why people reject lucrative opportunities:
\begin{equation}
\Delta Q = \omega_{INU} \cdot \Delta U^{INU} + \gamma_{INU,IDN} \cdot U^{INU} \cdot \Delta U^{IDN}
\end{equation}

If $\Delta U^{INU} > 0$ (financial gain) but $\Delta U^{IDN} < 0$ (identity violation),
and $|\gamma_{INU,IDN} \cdot \Delta U^{IDN}|$ is large, then $\Delta Q < 0$ despite 
the financial gain.

\paragraph{Empirical Estimates by Situation.}

\begin{table}[htbp]
\centering
\caption{$\gamma_{INU,IDN}$ by Situation}
\label{tab:app-c-gamma-inu-idn-context}
\renewcommand{\arraystretch}{1.3}
\begin{tabular}{@{}llc@{}}
\toprule
\textbf{Situation} & \textbf{Characteristic} & $\gamma_{INU,IDN}$ \\
\midrule
Success confirms identity & ``This is who I am'' & $+0.3$ to $+0.5$ \\
Neutral (identity not salient) & No identity implications & $0$ to $+0.1$ \\
Success questions identity & ``Is this really me?'' & $-0.1$ to $-0.3$ \\
Success violates identity & ``Selling out'' & $-0.3$ to $-0.6$ \\
\bottomrule
\end{tabular}
\end{table}

%-----------------------------------------------------------------------
\subsection{$\gamma_{KNU,IDN}$: Collective--Identity Complementarity}
%-----------------------------------------------------------------------

\paragraph{The Question.}
Does group success strengthen or threaten personal identity?

\paragraph{Positive Complementarity ($\gamma_{KNU,IDN} > 0$).}
Group and identity reinforce each other:

\begin{table}[htbp]
\centering
\caption{Positive KNU-IDN Complementarity Examples}
\label{tab:app-c-knu-idn-positive}
\renewcommand{\arraystretch}{1.3}
\begin{tabular}{@{}lll@{}}
\toprule
\textbf{Situation} & \textbf{Mechanism} & \textbf{Effect} \\
\midrule
``We are miners'' & Group identity & Collective success = identity affirmation \\
National pride & Patriotic identity & Nation's success enhances self-worth \\
Team identification & ``I am a [team] fan'' & Team victory feels personal \\
Religious community & Faith identity & Community flourishing = personal meaning \\
\bottomrule
\end{tabular}
\end{table}

When $\gamma_{KNU,IDN} > 0$, collective success \emph{validates} personal identity.
This is the foundation of group-based identity and nationalism.

\paragraph{Negative Complementarity ($\gamma_{KNU,IDN} < 0$).}
Group demands conflict with personal identity:

\begin{table}[htbp]
\centering
\caption{Negative KNU-IDN Complementarity Examples}
\label{tab:app-c-knu-idn-negative}
\renewcommand{\arraystretch}{1.3}
\begin{tabular}{@{}lll@{}}
\toprule
\textbf{Situation} & \textbf{Identity Conflict} & \textbf{Effect} \\
\midrule
Conformity pressure & ``I must fit in'' vs. authenticity & Group success, identity loss \\
Groupthink & Personal values overridden & Collective gain, self-betrayal \\
Cult dynamics & Identity subsumed by group & KNU$\uparrow$ but IDN$\downarrow$ \\
Whistleblower dilemma & Group loyalty vs. integrity & Identity demands betraying group \\
\bottomrule
\end{tabular}
\end{table}

When $\gamma_{KNU,IDN} < 0$, serving the group feels like \emph{losing oneself}.

\paragraph{Empirical Estimates by Context.}

\begin{table}[htbp]
\centering
\caption{$\gamma_{KNU,IDN}$ by Context}
\label{tab:app-c-gamma-knu-idn-context}
\renewcommand{\arraystretch}{1.3}
\begin{tabular}{@{}llc@{}}
\toprule
\textbf{Context} & \textbf{Characteristic} & $\gamma_{KNU,IDN}$ \\
\midrule
Strong group identity & ``We are one'' & $+0.4$ to $+0.6$ \\
Moderate identification & Member but individual & $+0.1$ to $+0.3$ \\
Individualistic culture & Self $>$ group & $0$ to $-0.1$ \\
Oppressive group & Group demands conformity & $-0.2$ to $-0.4$ \\
\bottomrule
\end{tabular}
\end{table}

%-----------------------------------------------------------------------
\subsection{The Complementarity Triad: Coherence and Conflict}
%-----------------------------------------------------------------------

\paragraph{Utility Coherence.}
The three complementarities together determine whether the utility system 
is coherent or conflicted:

\begin{tcolorbox}[colback=blue!5!white, colframe=blue!50!black, title=Definition: Utility Coherence Index]
\begin{equation}
\boxed{
K^{utility} = \frac{1}{3} \left( \text{sign}(\gamma_{INU,KNU}) + \text{sign}(\gamma_{INU,IDN}) + \text{sign}(\gamma_{KNU,IDN}) \right)
}
\label{eq:app-c-utility-coherence}
\end{equation}
\end{tcolorbox}

\begin{table}[htbp]
\centering
\caption{Utility Coherence States}
\label{tab:app-c-utility-coherence-states}
\renewcommand{\arraystretch}{1.3}
\begin{tabular}{@{}cccll@{}}
\toprule
$\gamma_{INU,KNU}$ & $\gamma_{INU,IDN}$ & $\gamma_{KNU,IDN}$ & $K^{utility}$ & \textbf{State} \\
\midrule
+ & + & + & +1 & \textbf{Full coherence} \\
+ & + & $-$ & +1/3 & Partial coherence \\
+ & $-$ & + & +1/3 & Partial coherence \\
$-$ & + & + & +1/3 & Partial coherence \\
+ & $-$ & $-$ & $-$1/3 & Partial conflict \\
$-$ & + & $-$ & $-$1/3 & Partial conflict \\
$-$ & $-$ & + & $-$1/3 & Partial conflict \\
$-$ & $-$ & $-$ & $-$1 & \textbf{Full conflict} \\
\bottomrule
\end{tabular}
\end{table}

\paragraph{Full Coherence (All Positive).}
When all three $\gamma > 0$:
\begin{itemize}
\item Individual success benefits the group
\item Personal achievement confirms identity
\item Group success validates individual identity
\end{itemize}

This is the \emph{ideal state}---all utility categories reinforce each other.
Example: A successful entrepreneur who builds a thriving community business 
that expresses their values.

\paragraph{Full Conflict (All Negative).}
When all three $\gamma < 0$:
\begin{itemize}
\item Individual gain harms the group
\item Personal success violates identity
\item Group demands betray individual values
\end{itemize}

This is a \emph{pathological state}---no choice can satisfy all categories.
Example: A corrupt official who enriches themselves at public expense while 
violating their self-image as honest and serving a system they despise.

%-----------------------------------------------------------------------
\subsection{Complementarities and Behavioral Change}
%-----------------------------------------------------------------------

\paragraph{Why Complementarities Matter for EBF.}
The complementarity structure determines whether behavioral change is 
reinforced or resisted:

\begin{table}[htbp]
\centering
\caption{Complementarities and Change Dynamics}
\label{tab:app-c-complementarities-change}
\renewcommand{\arraystretch}{1.3}
\begin{tabular}{@{}lll@{}}
\toprule
\textbf{Complementarity Pattern} & \textbf{Change Dynamic} & \textbf{Implication} \\
\midrule
All positive & Self-reinforcing & Change in one $\to$ supports others \\
Mixed & Partial resistance & Some dimensions resist \\
All negative & Self-undermining & Change in one $\to$ harms others \\
\bottomrule
\end{tabular}
\end{table}

\paragraph{The Coal Miner Revisited.}
Consider the coal miner's transition to coder through the complementarity lens:

\begin{table}[htbp]
\centering
\caption{Coal Miner Transition: Complementarity Analysis}
\label{tab:app-c-miner-complementarity}
\renewcommand{\arraystretch}{1.3}
\begin{tabular}{@{}lll@{}}
\toprule
\textbf{Complementarity} & \textbf{Status Quo (Mining)} & \textbf{After Transition (Coding)} \\
\midrule
$\gamma_{INU,KNU}$ & + (supports community) & $-$ (leaves community) \\
$\gamma_{INU,IDN}$ & + (mining fits identity) & $-$ (coding violates identity) \\
$\gamma_{KNU,IDN}$ & + (``we are miners'') & $-$ (betraying ``us'') \\
\midrule
$K^{utility}$ & +1 (coherent) & $-$1 (conflicted) \\
\bottomrule
\end{tabular}
\end{table}

The transition moves from full coherence to full conflict---even if $U^{INU}$ 
increases, the complementarity collapse makes total welfare decrease.

\paragraph{Intervention Implication.}
Successful behavioral change requires maintaining or rebuilding complementarities:
\begin{itemize}
\item \textbf{Identity bridge:} Frame new behavior as consistent with identity
\item \textbf{Collective transition:} Move the group together, preserving KNU
\item \textbf{Narrative integration:} Help construct a story where change fits
\end{itemize}

%-----------------------------------------------------------------------
\subsection{Context Effects on Complementarities}
%-----------------------------------------------------------------------

\paragraph{Complementarities Are Context-Dependent.}
The $\gamma_{ij}$ coefficients respond to the $\Psi$-dimensions (Chapter 9):

\begin{tcolorbox}[colback=green!5!white, colframe=green!50!black]
\begin{equation}
\boxed{
\gamma_{ij}(\Psi) = \bar{\gamma}_{ij} + \sum_k \Gamma^{\gamma}_{ij,k} \cdot \Psi_k
}
\label{eq:app-c-gamma-context}
\end{equation}
\end{tcolorbox}

\paragraph{Key Context Effects.}

\begin{table}[htbp]
\centering
\caption{Context Effects on Complementarities}
\label{tab:app-c-context-complementarities}
\renewcommand{\arraystretch}{1.3}
\begin{tabular}{@{}llll@{}}
\toprule
\textbf{$\Psi$-Dimension} & \textbf{Effect on $\gamma_{INU,KNU}$} & \textbf{Effect on $\gamma_{INU,IDN}$} & \textbf{Effect on $\gamma_{KNU,IDN}$} \\
\midrule
$\Psi_S$ (Trust) & $\uparrow$ (alignment increases) & $\sim$ (neutral) & $\uparrow$ (group identity strengthens) \\
$\Psi_I$ (Institutions) & $\uparrow$ (less zero-sum) & $\sim$ (neutral) & $\sim$ (neutral) \\
$\Psi_M$ (Markets) & $\downarrow$ (more competition) & $\downarrow$ (selling out easier) & $\downarrow$ (individualism rises) \\
$\Psi_T$ (Stability) & $\uparrow$ (long-term alignment) & $\uparrow$ (identity more stable) & $\uparrow$ (group bonds strengthen) \\
\bottomrule
\end{tabular}
\end{table}

\paragraph{Example: Globalization Effect.}
Market integration ($\Psi_M \uparrow$) tends to:
\begin{itemize}
\item Decrease $\gamma_{INU,KNU}$: More competition, less collective benefit sharing
\item Decrease $\gamma_{INU,IDN}$: More opportunities to ``sell out''
\item Decrease $\gamma_{KNU,IDN}$: Weaker group identities, more individualism
\end{itemize}

This explains why globalization can increase material welfare ($U^{INU}_F$) 
while decreasing overall wellbeing---the complementarity structure deteriorates.

%-----------------------------------------------------------------------
\subsection{Summary: Inter-Category Complementarities}
%-----------------------------------------------------------------------

\begin{tcolorbox}[colback=gray!5!white, colframe=gray!75!black, title=Section C.9 Summary]
\textbf{The Extended Welfare Function:}
$$Q = \sum_{i} \omega_i \cdot U^i + \sum_{i < j} \gamma_{ij} \cdot U^i \cdot U^j$$

\textbf{The Three Complementarities:}
\begin{itemize}
\item $\gamma_{INU,KNU}$: Individual--Collective alignment ($\pm 0.15$ to $\pm 0.50$)
\item $\gamma_{INU,IDN}$: Individual--Identity alignment ($\pm 0.10$ to $\pm 0.60$)
\item $\gamma_{KNU,IDN}$: Collective--Identity alignment ($\pm 0.10$ to $\pm 0.60$)
\end{itemize}

\textbf{Interpretation:}
\begin{itemize}
\item $\gamma > 0$: Categories reinforce each other (complement)
\item $\gamma < 0$: Categories undermine each other (substitute)
\item $\gamma = 0$: Categories are independent
\end{itemize}

\textbf{Utility Coherence Index:}
$$K^{utility} = \frac{1}{3} \left( \text{sign}(\gamma_{INU,KNU}) + \text{sign}(\gamma_{INU,IDN}) + \text{sign}(\gamma_{KNU,IDN}) \right)$$
\begin{itemize}
\item $K = +1$: Full coherence (ideal state)
\item $K = -1$: Full conflict (pathological state)
\item $|K| < 1$: Partial coherence/conflict
\end{itemize}

\textbf{Context Modulation:}
$$\gamma_{ij}(\Psi) = \bar{\gamma}_{ij} + \sum_k \Gamma^{\gamma}_{ij,k} \cdot \Psi_k$$

\textbf{Behavioral Change Implication:}
\begin{itemize}
\item Successful change requires maintaining complementarity structure
\item Identity bridges, collective transitions, narrative integration
\end{itemize}
\end{tcolorbox}

%%%%%%%%%%%%%%%%%%%%%%%%%%%%%%%%%%%%%%%%%%%%%%%%%%%%%%%%%%%%%%%%%%%%%%%%%%%%%%%
%%%%%%%%%%%%%%%%%%%%%%%%%%%%%%%%%%%%%%%%%%%%%%%%%%%%%%%%%%%%%%%%%%%%%%%%%%%%%%%
\section{Counting Verification}
\label{sec:app-c-counting}
%%%%%%%%%%%%%%%%%%%%%%%%%%%%%%%%%%%%%%%%%%%%%%%%%%%%%%%%%%%%%%%%%%%%%%%%%%%%%%%
%%%%%%%%%%%%%%%%%%%%%%%%%%%%%%%%%%%%%%%%%%%%%%%%%%%%%%%%%%%%%%%%%%%%%%%%%%%%%%%

\begin{table}[htbp]
\centering
\caption{Counting Verification: Four Paths to 144}
\label{tab:app-c-counting}
\renewcommand{\arraystretch}{1.3}
\begin{tabular}{@{}lll@{}}
\toprule
\textbf{Counting Path} & \textbf{Calculation} & \textbf{Result} \\
\midrule
By Category & $(48_{\text{INU}} + 48_{\text{KNU}} + 48_{\text{IDN}})$ & 144 \\
By Dimension & $(24_{\text{per dim}} \times 6_{\text{dims}})$ & 144 \\
By Time & $(36_{\text{per horizon}} \times 4_{\text{horizons}})$ & 144 \\
By Valence & $(72_{\text{Gain}} + 72_{\text{Pain}})$ & 144 \\
\bottomrule
\end{tabular}
\end{table}

%%%%%%%%%%%%%%%%%%%%%%%%%%%%%%%%%%%%%%%%%%%%%%%%%%%%%%%%%%%%%%%%%%%%%%%%%%%%%%%
%%%%%%%%%%%%%%%%%%%%%%%%%%%%%%%%%%%%%%%%%%%%%%%%%%%%%%%%%%%%%%%%%%%%%%%%%%%%%%%
\section{Worked Example: Labor Transition Decision}
\label{sec:app-c-example}
%%%%%%%%%%%%%%%%%%%%%%%%%%%%%%%%%%%%%%%%%%%%%%%%%%%%%%%%%%%%%%%%%%%%%%%%%%%%%%%
%%%%%%%%%%%%%%%%%%%%%%%%%%%%%%%%%%%%%%%%%%%%%%%%%%%%%%%%%%%%%%%%%%%%%%%%%%%%%%%

\begin{cross-reference}
\textbf{Chapter Reference:} See Chapter~\ref{sec:welfare-fepsde}, 
Section~\ref{subsec:10-9-applications}.
\end{cross-reference}

\subsection{Case: Coal Miner $\to$ Coder Retraining}

A coal miner is offered a \$50,000 retraining package to become a software developer.
Standard economics predicts acceptance (NPV positive). Observed: 70\%+ rejection.

\subsection{144-Component Analysis}

\begin{table}[htbp]
\centering
\caption{Coal Miner Decision: Key Utility Components}
\label{tab:app-c-example}
\small
\renewcommand{\arraystretch}{1.3}
\begin{tabular}{@{}llccc@{}}
\toprule
\textbf{Category} & \textbf{Component} & \textbf{Gain} & \textbf{Pain} & \textbf{Net} \\
\midrule
\multirow{3}{*}{INU} 
& $U^{\text{INU}}_{F,2}$ (Medium-term financial) & +\$20k/yr & & +0.8 \\
& $U^{\text{INU}}_{P,2}$ (Medium-term physical) & Safer work & & +0.4 \\
& $U^{\text{INU}}_{E,1}$ (Short-term emotional) & & Anxiety & --0.3 \\
\midrule
\multirow{2}{*}{KNU}
& $U^{\text{KNU}}_{S,2}$ (Community social) & & Leaving community & --0.5 \\
& $U^{\text{KNU}}_{F,3}$ (Long-term collective) & Family benefits & Community loss & $\pm$0.0 \\
\midrule
\multirow{3}{*}{IDN}
& $U^{\text{IDN}}_{S,0}$ (Instant identity) & & ``Not a miner'' & \textbf{--1.2} \\
& $U^{\text{IDN}}_{S,2}$ (Medium identity) & & ``Betraying roots'' & \textbf{--0.8} \\
& $U^{\text{IDN}}_{E,3}$ (Long-term meaning) & & ``Who am I?'' & --0.5 \\
\midrule
\multicolumn{4}{l}{\textbf{Aggregate by Category:}} & \\
& $U^{\text{INU}}$ & & & $\approx +0.9$ \\
& $U^{\text{KNU}}$ & & & $\approx -0.5$ \\
& $U^{\text{IDN}}$ & & & $\approx \mathbf{-2.5}$ \\
\midrule
\multicolumn{4}{l}{\textbf{Total Effective Utility:}} & $\mathbf{-2.1}$ \\
\bottomrule
\end{tabular}
\end{table}

\subsection{Key Insight}

The decision is dominated by \textbf{Identity Utility} ($U^{\text{IDN}} \approx -2.5$), 
which overwhelms the positive Individual Utility ($U^{\text{INU}} \approx +0.9$).

%%%%%%%%%%%%%%%%%%%%%%%%%%%%%%%%%%%%%%%%%%%%%%%%%%%%%%%%%%%%%%%%%%%%%%%%%%%%%%%
%%%%%%%%%%%%%%%%%%%%%%%%%%%%%%%%%%%%%%%%%%%%%%%%%%%%%%%%%%%%%%%%%%%%%%%%%%%%%%%
\section{Validation}
\label{sec:app-c-validation}
%%%%%%%%%%%%%%%%%%%%%%%%%%%%%%%%%%%%%%%%%%%%%%%%%%%%%%%%%%%%%%%%%%%%%%%%%%%%%%%
%%%%%%%%%%%%%%%%%%%%%%%%%%%%%%%%%%%%%%%%%%%%%%%%%%%%%%%%%%%%%%%%%%%%%%%%%%%%%%%

\begin{table}[htbp]
\centering
\caption{$\Gamma$-Matrix Validation Against Empirical Benchmarks}
\label{tab:app-c-validation}
\renewcommand{\arraystretch}{1.3}
\begin{tabular}{@{}llccc@{}}
\toprule
\textbf{Coefficient} & \textbf{Empirical Source} & \textbf{Empirical} & \textbf{LLM-MC} & \textbf{$|\Delta|$} \\
\midrule
$\gamma_{F,I}$ & Acemoglu et al. (2001) & 0.28 & 0.31 & 0.03 \\
$\gamma_{E,S}$ & Helliwell \& Putnam (2004) & 0.44 & 0.47 & 0.03 \\
$\gamma_{P,C}$ & Cutler \& Lleras-Muney (2010) & 0.35 & 0.31 & 0.04 \\
$\gamma_{S,M}$ & Autor et al. (2013) & --0.21 & --0.18 & 0.03 \\
$\gamma_{Eco,E}$ & EKC meta-analyses & --0.15 & --0.19 & 0.04 \\
\midrule
\multicolumn{4}{r}{\textbf{Mean Absolute Error:}} & \textbf{0.034} \\
\multicolumn{4}{r}{\textbf{Sign Agreement:}} & \textbf{5/5 (100\%)} \\
\bottomrule
\end{tabular}
\end{table}

%%%%%%%%%%%%%%%%%%%%%%%%%%%%%%%%%%%%%%%%%%%%%%%%%%%%%%%%%%%%%%%%%%%%%%%%%%%%%%%
%%%%%%%%%%%%%%%%%%%%%%%%%%%%%%%%%%%%%%%%%%%%%%%%%%%%%%%%%%%%%%%%%%%%%%%%%%%%%%%
%==============================================================================
% CRITICAL FOUNDATIONS / OBJECTIONS (Required for CORE)
%==============================================================================

\section{Critical Foundations and Objections}
\label{sec:C-foundations}

\subsection{Objection 1: Why Exactly Six Dimensions?}

\textbf{Objection:} The FEPSDE framework seems arbitrary. Why not five dimensions, 
or seven? What makes this particular decomposition privileged?

\textbf{Response:}
\begin{enumerate}[nosep]
    \item \textbf{Empirical grounding:} The six dimensions emerge from factor 
    analysis of wellbeing surveys (Stiglitz-Sen-Fitoussi 2009; Diener et al. 1999).
    \item \textbf{Completeness:} Any welfare-relevant outcome can be mapped to 
    at least one dimension.
    \item \textbf{Orthogonality:} Dimensions are approximately independent in 
    variation---high F does not imply high P.
    \item \textbf{Policy relevance:} Each dimension corresponds to distinct 
    policy levers (financial: taxes; health: healthcare; etc.).
    \item \textbf{Extensibility:} The framework is modular---dimensions can be 
    added if empirically warranted.
\end{enumerate}

\subsection{Objection 2: 144 Components is Overwhelming}

\textbf{Objection:} No practical application can handle 144 separate utility 
components. The decomposition is too fine-grained.

\textbf{Response:}
\begin{enumerate}[nosep]
    \item \textbf{Principle of selection:} For any specific application, only 
    a subset of components is relevant.
    \item \textbf{Aggregation:} The framework supports aggregation---sum over 
    time, or over valence, as needed.
    \item \textbf{Sparsity:} Most cross-effects are negligible; the effective 
    dimensionality is much lower.
    \item \textbf{Completeness value:} The full specification ensures nothing 
    is overlooked \textit{a priori}.
\end{enumerate}

\subsection{Objection 3: Ecological Dimension is Novel}

\textbf{Objection:} Traditional welfare economics includes F, E, P, S. Adding 
"Eco" seems like a political statement rather than scientific necessity.

\textbf{Response:}
\begin{enumerate}[nosep]
    \item \textbf{Empirical evidence:} Environmental quality affects reported 
    wellbeing independently of other dimensions (Levinson 2012).
    \item \textbf{Intergenerational:} Ecological utility captures 
    sustainability---welfare of future selves and generations.
    \item \textbf{Complementarity:} Eco interacts non-trivially with other 
    dimensions (health-environment, wealth-sustainability).
    \item \textbf{Policy relevance:} Climate and environment are major policy 
    domains that require explicit representation.
\end{enumerate}

\subsection{Objection 4: Digital Dimension is Too New}

\textbf{Objection:} "Digital" as a welfare dimension is a recent phenomenon. 
How can it be fundamental?

\textbf{Response:}
\begin{enumerate}[nosep]
    \item \textbf{Ubiquity:} For much of the world's population, digital access 
    is as fundamental as other dimensions.
    \item \textbf{Distinct effects:} Digital wellbeing (connectivity, privacy, 
    data autonomy) varies independently of traditional dimensions.
    \item \textbf{Future-proofing:} As AI and digital technologies expand, 
    this dimension becomes more, not less, relevant.
    \item \textbf{Empirical:} Digital exclusion has measurable welfare effects 
    distinct from financial poverty.
\end{enumerate}


%==============================================================================
% OPEN ISSUES / RESEARCH AGENDA (Required)
%==============================================================================

\section{Open Issues and Research Agenda}
\label{sec:C-open-issues}

\subsection{Measurement Priorities}

\begin{enumerate}
    \item \textbf{Ecological utility measurement:} Developing reliable, 
    cross-culturally valid measures of Eco-dimension utility.
    
    \item \textbf{Digital wellbeing scales:} Standardized instruments for 
    D-dimension that distinguish access, quality, and autonomy.
    
    \item \textbf{Time horizon calibration:} Empirical determination of the 
    four time horizons (instant/short/medium/long) by dimension and context.
    
    \item \textbf{Loss aversion heterogeneity:} Does $\lambda$ vary 
    systematically across FEPSDE dimensions and populations?
\end{enumerate}

\subsection{Theoretical Extensions}

\begin{enumerate}
    \item \textbf{Dimension emergence:} Under what conditions might new 
    dimensions become empirically necessary? What would trigger adding a 
    seventh dimension?
    
    \item \textbf{Hierarchy vs. orthogonality:} Is there a natural ordering 
    among dimensions (à la Maslow), or are they truly orthogonal?
    
    \item \textbf{Identity utility foundations:} Deeper microfoundations for 
    IDN---when does identity-based utility dominate INU?
    
    \item \textbf{Cross-level interaction:} How do INU, KNU, and IDN interact 
    when they conflict?
\end{enumerate}

\subsection{Empirical Validation}

\begin{enumerate}
    \item \textbf{Factor structure:} Confirmatory factor analysis of FEPSDE 
    across diverse populations and cultures.
    
    \item \textbf{Predictive validity:} Does the 144-component decomposition 
    predict behavior better than simpler models?
    
    \item \textbf{Axiom testing:} Experimental tests of specific INU/KNU/IDN 
    axioms, especially controversial ones.
    
    \item \textbf{Temporal dynamics:} How stable are dimension weights over 
    the life course?
\end{enumerate}


\section{Summary}
\label{sec:app-c-summary}
%%%%%%%%%%%%%%%%%%%%%%%%%%%%%%%%%%%%%%%%%%%%%%%%%%%%%%%%%%%%%%%%%%%%%%%%%%%%%%%
%%%%%%%%%%%%%%%%%%%%%%%%%%%%%%%%%%%%%%%%%%%%%%%%%%%%%%%%%%%%%%%%%%%%%%%%%%%%%%%

\begin{tcolorbox}[colback=gray!5!white, colframe=gray!75!black, title=Appendix C Summary]

\textbf{Meta-Axiom U-M0 (universal):}
$$\boxed{U^{i} = \sum_{d} \sum_{t} \delta^{i}_{d}(t) \cdot [G^{i}_{d,t} - \lambda^{i}_{d} \cdot P^{i}_{d,t}]}$$

\textbf{Structure:} $3 \times 2 \times 6 \times 4 = 144$ components

\textbf{Axiom Systems (220 Total):}
\begin{itemize}
\item INU: 59 axioms defining $G^{INU}$, $P^{INU}$, $\lambda^{INU}$, $\delta^{INU}$
\item KNU: 65 axioms defining $G^{KNU}$, $P^{KNU}$, $\lambda^{KNU}$, $\delta^{KNU}$
\item IDN: 95 axioms defining $G^{IDN}$, $P^{IDN}$, $\lambda^{IDN}$, $\delta^{IDN}$
\end{itemize}

\textbf{Parameter Structure (all categories):}
\begin{itemize}
\item $G^i$: Gain with modulators $\Omega^{G,i}$
\item $P^i$: Pain with modulators $\Omega^{P,i}$
\item $\lambda^i$: Loss aversion with context modulators $\Lambda^i$
\item $\delta^i$: Discounting with time modulators $\Delta^i$
\end{itemize}

\textbf{Inter-dimensional Complementarities:}
\begin{itemize}
\item INU: 15 coefficients (A43-A57)
\item KNU: 15 coefficients (A51-A65)
\item IDN: 15 coefficients (A80-A94)
\end{itemize}

\textbf{$\Gamma$-Matrix:} 48 coefficients ($6 \times 8$) for context-welfare interaction

\textbf{Key Findings:}
\begin{itemize}
\item Strongest positive: $\gamma_{S,S} = 0.51$ (trust $\rightarrow$ social welfare)
\item Structural trade-offs: $\gamma_{S,M} = -0.18$, $\gamma_{Eco,M} = -0.21$
\item Identity loss aversion: $\lambda^{IDN} > \lambda^{INU}$, $\lambda^{IDN} > \lambda^{KNU}$ (persistence effect)
\end{itemize}

\textbf{Validation:} MAE = 0.034 against empirical benchmarks; 100\% sign agreement.

\end{tcolorbox}

% =============================================================================
% APPENDIX C GLOSSARY
% =============================================================================
% Complete reference of symbols, variables, and terms used in Appendix C
% =============================================================================

\section{Glossary of Symbols and Terms}
\label{sec:app-c-glossary}

%%%%%%%%%%%%%%%%%%%%%%%%%%%%%%%%%%%%%%%%%%%%%%%%%%%%%%%%%%%%%%%%%%%%%%%%%%%%%%%
\subsection{Core Utility Symbols}
%%%%%%%%%%%%%%%%%%%%%%%%%%%%%%%%%%%%%%%%%%%%%%%%%%%%%%%%%%%%%%%%%%%%%%%%%%%%%%%

\begin{longtable}{@{}lp{4cm}p{7cm}@{}}
\caption{Core Utility Symbols} \\
\toprule
\textbf{Symbol} & \textbf{Name} & \textbf{Definition} \\
\midrule
\endfirsthead
\midrule
\endhead
\bottomrule
\endlastfoot

$U^i$ & Utility (Category $i$) & Total utility in category $i \in \{INU, KNU, IDN\}$ \\
$U^{INU}$ & Individual Utility & Self-regarding utility: ``What do I want for MYSELF?'' \\
$U^{KNU}$ & Collective Utility & Other-regarding utility: ``What do I want for US?'' \\
$U^{IDN}$ & Identity Utility & Self-definitional utility: ``Who AM I?'' \\
$U^{total}$ & Aggregate Utility & Weighted sum of all three categories with complementarities \\
\addlinespace
$G^i_{d,t}$ & Gain & Positive utility in category $i$, dimension $d$, time $t$ \\
$P^i_{d,t}$ & Pain & Negative utility in category $i$, dimension $d$, time $t$ \\
$\lambda^i_d$ & Loss Aversion & Coefficient for asymmetric weighting of pains vs. gains \\
$\delta^i_d(t)$ & Discount Factor & Time preference factor for future outcomes \\
\addlinespace
$R^i_{d,t}$ & Reference Point & Benchmark against which outcomes are evaluated \\
$x_{d,t}$ & Outcome & Actual outcome in dimension $d$ at time $t$ \\
$[x]^+$ & Positive Part & $\max(0, x)$ -- only positive values count \\

\end{longtable}

%%%%%%%%%%%%%%%%%%%%%%%%%%%%%%%%%%%%%%%%%%%%%%%%%%%%%%%%%%%%%%%%%%%%%%%%%%%%%%%
\subsection{FEPSDE Dimensions}
%%%%%%%%%%%%%%%%%%%%%%%%%%%%%%%%%%%%%%%%%%%%%%%%%%%%%%%%%%%%%%%%%%%%%%%%%%%%%%%

\begin{longtable}{@{}lp{3cm}p{8cm}@{}}
\caption{FEPSDE Welfare Dimensions} \\
\toprule
\textbf{Symbol} & \textbf{Dimension} & \textbf{Examples} \\
\midrule
\endfirsthead
\midrule
\endhead
\bottomrule
\endlastfoot

$F$ & Financial & Income, wealth, economic security, purchasing power \\
$E$ & Emotional & Joy, satisfaction, fulfillment, meaning, psychological wellbeing \\
$P$ & Physical & Health, energy, fitness, bodily integrity, longevity \\
$S$ & Social & Status, belonging, recognition, relationships, reputation \\
$D$ & Digital & Connectivity, privacy, digital competence, online presence \\
$Eco$ & Ecological & Environmental quality, sustainability, nature connection \\
\addlinespace
$\mathcal{D}$ & Dimension Set & $\mathcal{D} = \{F, E, P, S, D, Eco\}$ (6 dimensions) \\
$d$ & Dimension Index & $d \in \mathcal{D}$ \\

\end{longtable}

%%%%%%%%%%%%%%%%%%%%%%%%%%%%%%%%%%%%%%%%%%%%%%%%%%%%%%%%%%%%%%%%%%%%%%%%%%%%%%%
\subsection{Time Horizons}
%%%%%%%%%%%%%%%%%%%%%%%%%%%%%%%%%%%%%%%%%%%%%%%%%%%%%%%%%%%%%%%%%%%%%%%%%%%%%%%

\begin{longtable}{@{}lp{3cm}p{8cm}@{}}
\caption{Time Horizons} \\
\toprule
\textbf{Symbol} & \textbf{Horizon} & \textbf{Definition} \\
\midrule
\endfirsthead
\midrule
\endhead
\bottomrule
\endlastfoot

$t = 0$ & Instant & Immediate, visceral, now (seconds to minutes) \\
$t = 1$ & Short-term & Days to weeks \\
$t = 2$ & Medium-term & Months to years \\
$t = 3$ & Long-term & Years to decades, legacy effects \\
\addlinespace
$\beta$ & Present Bias & Quasi-hyperbolic discounting parameter ($\beta < 1$) \\
$\bar{\delta}$ & Base Discount & Standard exponential discount rate \\

\end{longtable}

%%%%%%%%%%%%%%%%%%%%%%%%%%%%%%%%%%%%%%%%%%%%%%%%%%%%%%%%%%%%%%%%%%%%%%%%%%%%%%%
\subsection{INU-Specific Symbols}
%%%%%%%%%%%%%%%%%%%%%%%%%%%%%%%%%%%%%%%%%%%%%%%%%%%%%%%%%%%%%%%%%%%%%%%%%%%%%%%

\begin{longtable}{@{}lp{4cm}p{7cm}@{}}
\caption{INU (Individual Utility) Symbols} \\
\toprule
\textbf{Symbol} & \textbf{Name} & \textbf{Definition / Axiom} \\
\midrule
\endfirsthead
\midrule
\endhead
\bottomrule
\endlastfoot

$V_{base}$ & Baseline Value & Intrinsic value of ``doing'' (A3) \\
$V_{intr}$ & Intrinsic Value & Task-inherent motivation (A4) \\
$V_{extr}$ & Extrinsic Value & External reward value (A5) \\
$\kappa$ & Intr-Extr Complementarity & Interaction between intrinsic and extrinsic (A6) \\
\addlinespace
$\Omega^G_d$ & Gain Modulator & Combined cognitive/perceptual filters on gains \\
$\Omega^P_d$ & Pain Modulator & Combined stress/coping filters on pains \\
$\Lambda_d$ & Lambda Modulator & Context effects on loss aversion \\
$\Delta_d(t)$ & Delta Modulator & Context effects on discounting \\
\addlinespace
$\phi_{risk}$ & Risk Aversion & Reduction due to risk (A9) \\
$\phi_{amb}$ & Ambiguity Aversion & Reduction due to ambiguity (A9) \\
$\phi_{motiv}$ & Motivated Beliefs & Belief distortion for pain avoidance (A10) \\
$\phi_{blind}$ & Blind Spots & Unperceived utility components (A11) \\
$S_d$ & Salience & Attention-weighted prominence (A15) \\
$H_d$ & Habit & Habituation dampening factor (A20) \\
$\alpha_d$ & Attention Share & Bounded attention allocation (A32) \\
\addlinespace
$Cope_d$ & Coping & Stress coping capacity (A22) \\
$Res_d$ & Resilience & Stress resilience capacity (A22) \\
$\mathcal{E}^-$ & Negative Emotions & Shame, guilt, embarrassment (A30) \\
\addlinespace
$\gamma_{gender}$ & Gender Effect & Gender-based parameter modulation (A17) \\
$\gamma_{culture}$ & Culture Effect & Culture-based parameter modulation (A18) \\
$\gamma_{lifecycle}$ & Lifecycle Effect & Age/phase-based parameter modulation (A16) \\

\end{longtable}

%%%%%%%%%%%%%%%%%%%%%%%%%%%%%%%%%%%%%%%%%%%%%%%%%%%%%%%%%%%%%%%%%%%%%%%%%%%%%%%
\subsection{KNU-Specific Symbols}
%%%%%%%%%%%%%%%%%%%%%%%%%%%%%%%%%%%%%%%%%%%%%%%%%%%%%%%%%%%%%%%%%%%%%%%%%%%%%%%

\begin{longtable}{@{}lp{4cm}p{7cm}@{}}
\caption{KNU (Collective Utility) Symbols} \\
\toprule
\textbf{Symbol} & \textbf{Name} & \textbf{Definition / Axiom} \\
\midrule
\endfirsthead
\midrule
\endhead
\bottomrule
\endlastfoot

$N_{kol}$ & Collective Utility & Aggregate group welfare (A1) \\
$\omega_d$ & Dimension Weight & Relative importance of dimension $d$ (A2) \\
$\gamma_d$ & Benefit Coefficient & Gain scaling per dimension (A14-A19) \\
$\delta_d$ & Cost Coefficient & Pain scaling per dimension (A14-A19) \\
\addlinespace
$\rho_K$ & Cooperation Type & Actor type: $\{0, f(B), 1\}$ = \{free-rider, conditional, unconditional\} (A9) \\
$B^{eff}$ & Effective Beliefs & Beliefs about others' cooperation (A8) \\
$U_{env}$ & Environmental Uncertainty & Uncertainty discount factor (A7) \\
\addlinespace
$P_{sanction}$ & Sanction Pain & Total punishment costs (A31-A32) \\
$P_{dec}$ & Decentralized Punishment & Peer punishment cost $c$ (A31) \\
$P_{cen}$ & Centralized Punishment & Institutional punishment $P_{inst} \cdot I_{acc}$ (A32) \\
$P_{DTX}$ & Dark Triad Destruction & Antisocial punishment amplified by Dark Triad (A46) \\
\addlinespace
$\Theta_{DT}$ & Dark Triad Score & $\alpha M + \beta N + \gamma P$ (Machiavellianism, Narcissism, Psychopathy) (A44) \\
$\Theta_{Beh}$ & Behavioral Indicator & Observable manipulation behaviors (A45) \\
$P_{anti}$ & Antisocial Punishment & Punishment of cooperators by defectors (A43) \\
\addlinespace
$Legit$ & Legitimacy & Perceived legitimacy of norms/institutions (A59) \\
$Coh$ & Norm Coherence & Consistency of normative expectations (A49) \\
$Stab_{KNU}$ & System Stability & Equilibrium metric (A50) \\
\addlinespace
$\varepsilon$ & Resilience Loss & System vulnerability parameter (A24) \\
$\xi$ & Tipping Point Shock & Discontinuous change trigger (A25) \\
$\delta^g$ & Intergenerational Discount & Discount across generations (A21) \\

\end{longtable}

%%%%%%%%%%%%%%%%%%%%%%%%%%%%%%%%%%%%%%%%%%%%%%%%%%%%%%%%%%%%%%%%%%%%%%%%%%%%%%%
\subsection{IDN-Specific Symbols}
%%%%%%%%%%%%%%%%%%%%%%%%%%%%%%%%%%%%%%%%%%%%%%%%%%%%%%%%%%%%%%%%%%%%%%%%%%%%%%%

\begin{longtable}{@{}lp{4cm}p{7cm}@{}}
\caption{IDN (Identity Utility) Symbols} \\
\toprule
\textbf{Symbol} & \textbf{Name} & \textbf{Definition / Axiom} \\
\midrule
\endfirsthead
\midrule
\endhead
\bottomrule
\endlastfoot

$j$ & Identity Index & Index for multiple identities (A2) \\
$\pi_j$ & Identity Weight & Salience/importance of identity $j$ (A2, A6) \\
$A_j$ & Awareness & Conscious activation of identity $j$ (A9) \\
$U^j$ & Identity-Specific Utility & Utility from identity $j$ \\
\addlinespace
$Recog_j$ & Recognition & Social acknowledgment of identity $j$ (A21) \\
$Belong_j$ & Belonging & Group membership feeling (A21) \\
$Meaning_j$ & Meaning & Purpose/significance from identity $j$ (A21) \\
$Exclus_j$ & Exclusion & Social rejection pain (A22) \\
$Stigma_j$ & Stigma & Identity-based discrimination (A22) \\
\addlinespace
$\kappa^+_{jk}$ & Positive Complementarity & Synergy between identities $j$ and $k$ (A24) \\
$\kappa^-_{jk}$ & Negative Complementarity & Conflict between identities $j$ and $k$ (A25) \\
$C_{conflict}$ & Conflict Cost & Total inter-identity conflict pain \\
$C_{frag}$ & Fragmentation Cost & Cost of too many contradictory identities (A47, A70) \\
\addlinespace
$\Omega$ & Fragmentation Penalty & Coherence loss weight (A70) \\
$K^{IDN}$ & Identity Coherence Index & Proportion of positive complementarities \\
$\psi$ & Emotional Intensity & Emotion amplification factor (A57) \\
\addlinespace
$Profile^j$ & Temporal Profile & $\{\omega_{inst}, \omega_{short}, \omega_{mid}, \omega_{long}\}$ per identity (A17) \\
$Persistence$ & Pain Persistence & Duration of identity pain effects (A38) \\
\addlinespace
$I_t$ & Self-Image & Current self-perception (A70b) \\
$S_t$ & Social Feedback & External perception/mirror (A70b) \\
$Action(t)$ & Synthesized Action & ArgMax of weighted identity utilities (A78) \\
$State(t)$ & System State & Stable vs. Collapse condition (A79) \\

\end{longtable}

%%%%%%%%%%%%%%%%%%%%%%%%%%%%%%%%%%%%%%%%%%%%%%%%%%%%%%%%%%%%%%%%%%%%%%%%%%%%%%%
\subsection{Context ($\Psi$) Dimensions}
%%%%%%%%%%%%%%%%%%%%%%%%%%%%%%%%%%%%%%%%%%%%%%%%%%%%%%%%%%%%%%%%%%%%%%%%%%%%%%%

\begin{longtable}{@{}lp{4cm}p{7cm}@{}}
\caption{8$\Psi$-Context Dimensions} \\
\toprule
\textbf{Symbol} & \textbf{Dimension} & \textbf{Definition} \\
\midrule
\endfirsthead
\midrule
\endhead
\bottomrule
\endlastfoot

$\Psi_I$ & Institutional Quality & Rule of law, property rights, contract enforcement \\
$\Psi_S$ & Social Capital & Trust, networks, norms of reciprocity \\
$\Psi_C$ & Cognitive Load & Information complexity, decision difficulty \\
$\Psi_K$ & Information Ecology & Media quality, misinformation prevalence \\
$\Psi_E$ & Economic Development & GDP, infrastructure, market sophistication \\
$\Psi_T$ & Temporal Stability & Predictability, planning horizon reliability \\
$\Psi_M$ & Market Integration & Globalization, competition intensity \\
$\Psi_F$ & Flexibility/Agency & Degrees of freedom, choice availability \\

\end{longtable}

%%%%%%%%%%%%%%%%%%%%%%%%%%%%%%%%%%%%%%%%%%%%%%%%%%%%%%%%%%%%%%%%%%%%%%%%%%%%%%%
\subsection{$\Gamma$-Matrix Symbols}
%%%%%%%%%%%%%%%%%%%%%%%%%%%%%%%%%%%%%%%%%%%%%%%%%%%%%%%%%%%%%%%%%%%%%%%%%%%%%%%

\begin{longtable}{@{}lp{4cm}p{7cm}@{}}
\caption{Context-Welfare Interaction Symbols} \\
\toprule
\textbf{Symbol} & \textbf{Name} & \textbf{Definition} \\
\midrule
\endfirsthead
\midrule
\endhead
\bottomrule
\endlastfoot

$\Gamma$ & Gamma Matrix & $6 \times 8$ matrix of context-welfare coefficients \\
$\gamma_{dj}$ & Interaction Coefficient & Effect of context dimension $j$ on welfare dimension $d$ \\
$\bar{U}_d$ & Baseline Welfare & Welfare in dimension $d$ at neutral context \\
\addlinespace
\multicolumn{3}{@{}l}{\textit{Key Coefficients:}} \\
$\gamma_{S,S}$ & Trust $\to$ Social & Strongest positive effect (0.51) \\
$\gamma_{F,E}$ & Development $\to$ Financial & Economic development effect (0.42) \\
$\gamma_{S,M}$ & Markets $\to$ Social & Negative trade-off ($-0.18$) \\
$\gamma_{Eco,M}$ & Markets $\to$ Ecological & Negative trade-off ($-0.21$) \\

\end{longtable}

%%%%%%%%%%%%%%%%%%%%%%%%%%%%%%%%%%%%%%%%%%%%%%%%%%%%%%%%%%%%%%%%%%%%%%%%%%%%%%%
\subsection{Complementarity Symbols}
%%%%%%%%%%%%%%%%%%%%%%%%%%%%%%%%%%%%%%%%%%%%%%%%%%%%%%%%%%%%%%%%%%%%%%%%%%%%%%%

\begin{longtable}{@{}lp{4cm}p{7cm}@{}}
\caption{Complementarity Symbols} \\
\toprule
\textbf{Symbol} & \textbf{Name} & \textbf{Definition} \\
\midrule
\endfirsthead
\midrule
\endhead
\bottomrule
\endlastfoot

$\sigma_{d_1,d_2}$ & FEPSDE Complementarity & Inter-dimensional interaction coefficient \\
$\sigma^{INU}_{d_1,d_2}$ & INU Complementarity & Individual utility dimension interaction (A43-A57) \\
$\sigma^{KNU}_{d_1,d_2}$ & KNU Complementarity & Collective utility dimension interaction (A51-A65) \\
$\sigma^{IDN}_{d_1,d_2}$ & IDN Complementarity & Identity utility dimension interaction (A80-A94) \\
\addlinespace
$\gamma_{ij}$ & Inter-Category Compl. & Interaction between utility categories $i$ and $j$ \\
$\gamma_{INU,KNU}$ & INU-KNU Interaction & Individual-Collective alignment/conflict \\
$\gamma_{INU,IDN}$ & INU-IDN Interaction & Individual-Identity alignment/conflict \\
$\gamma_{KNU,IDN}$ & KNU-IDN Interaction & Collective-Identity alignment/conflict \\

\end{longtable}

%%%%%%%%%%%%%%%%%%%%%%%%%%%%%%%%%%%%%%%%%%%%%%%%%%%%%%%%%%%%%%%%%%%%%%%%%%%%%%%
\subsection{Mathematical Operators and Notation}
%%%%%%%%%%%%%%%%%%%%%%%%%%%%%%%%%%%%%%%%%%%%%%%%%%%%%%%%%%%%%%%%%%%%%%%%%%%%%%%

\begin{longtable}{@{}lp{4cm}p{7cm}@{}}
\caption{Mathematical Notation} \\
\toprule
\textbf{Symbol} & \textbf{Name} & \textbf{Definition} \\
\midrule
\endfirsthead
\midrule
\endhead
\bottomrule
\endlastfoot

$\sum_d$ & Sum over dimensions & $\sum_{d \in \mathcal{D}} = \sum_{d \in \{F,E,P,S,D,Eco\}}$ \\
$\sum_t$ & Sum over time & $\sum_{t=0}^{3}$ (four horizons) \\
$\sum_j$ & Sum over identities & Sum over all active identities \\
$[x]^+$ & Positive part & $\max(0, x)$ \\
$|x|$ & Absolute value & Magnitude regardless of sign \\
\addlinespace
$f(\cdot)$ & Function & General functional relationship \\
$g(\cdot)$ & Function & Alternative functional relationship \\
$\partial / \partial x$ & Partial derivative & Rate of change with respect to $x$ \\
$\int_t^\infty$ & Integral & Cumulative effect from $t$ to infinity \\
\addlinespace
$\mathbb{1}[\cdot]$ & Indicator function & 1 if condition true, 0 otherwise \\
$\text{ArgMax}(\cdot)$ & Argument of maximum & Value that maximizes the expression \\
$\propto$ & Proportional to & Scales with (not exact equality) \\
$\gg$ & Much greater than & Significantly larger \\

\end{longtable}

%%%%%%%%%%%%%%%%%%%%%%%%%%%%%%%%%%%%%%%%%%%%%%%%%%%%%%%%%%%%%%%%%%%%%%%%%%%%%%%
\subsection{Acronyms and Abbreviations}
%%%%%%%%%%%%%%%%%%%%%%%%%%%%%%%%%%%%%%%%%%%%%%%%%%%%%%%%%%%%%%%%%%%%%%%%%%%%%%%

\begin{longtable}{@{}lp{4cm}p{7cm}@{}}
\caption{Acronyms and Abbreviations} \\
\toprule
\textbf{Acronym} & \textbf{Full Form} & \textbf{Definition} \\
\midrule
\endfirsthead
\midrule
\endhead
\bottomrule
\endlastfoot

BCM & Behavioral Change Model & The overarching framework (Version 2.0) \\
BEATRIX & -- & BCM software architecture \\
FEPSDE & Financial, Emotional, Physical, Social, Digital, Ecological & Six welfare dimensions \\
\addlinespace
INU & Individual Utility & Self-regarding utility module (6040) \\
KNU & Collective Utility & Other-regarding utility module (6041) \\
IDN & Identity Utility & Self-definitional utility module (6042) \\
AWX & Awareness Module & Attention and salience module \\
KON & Context Module & 8$\Psi$-dimensions context module \\
\addlinespace
LLM-MC & LLM Monte Carlo & Estimation methodology using language models \\
MAE & Mean Absolute Error & Validation metric \\
RT & Reaction Time & Neuroeconomic validation measure \\
\addlinespace
DTX & Dark Triad Index & Machiavellianism + Narcissism + Psychopathy \\
ESS & Evolutionarily Stable Strategy & Game-theoretic equilibrium concept \\
ROI & Return on Investment & Benefit-cost ratio \\

\end{longtable}

%%%%%%%%%%%%%%%%%%%%%%%%%%%%%%%%%%%%%%%%%%%%%%%%%%%%%%%%%%%%%%%%%%%%%%%%%%%%%%%
\subsection{Key Thresholds and Benchmarks}
%%%%%%%%%%%%%%%%%%%%%%%%%%%%%%%%%%%%%%%%%%%%%%%%%%%%%%%%%%%%%%%%%%%%%%%%%%%%%%%

\begin{longtable}{@{}lp{3cm}p{8cm}@{}}
\caption{Key Thresholds and Benchmarks} \\
\toprule
\textbf{Parameter} & \textbf{Value} & \textbf{Interpretation} \\
\midrule
\endfirsthead
\midrule
\endhead
\bottomrule
\endlastfoot

$\lambda$ (baseline) & $\approx 2.0$ & Kahneman-Tversky loss aversion coefficient \\
$\lambda^{IDN}$ & $> 2.0$ & Higher due to persistence effects \\
$\beta$ (present bias) & $\approx 0.7$ & Quasi-hyperbolic discounting parameter \\
\addlinespace
$Stab_{KNU} \geq 0.95$ & CE\_OK\_FULL & System considered stable \\
$Coherence < Min$ & Collapse & Identity system breakdown threshold \\
\addlinespace
$\gamma_{S,S} = 0.51$ & Strongest + & Trust $\to$ Social welfare (validated) \\
$\gamma_{S,M} = -0.18$ & Trade-off & Markets $\to$ Social welfare \\
$\gamma_{Eco,M} = -0.21$ & Trade-off & Markets $\to$ Ecological welfare \\
\addlinespace
MAE & 0.034 & Validation error against empirical benchmarks \\
Sign Agreement & 100\% & All coefficient signs match literature \\

\end{longtable}
% =============================================================================
% APPENDIX C BIBLIOGRAPHY
% =============================================================================
% Complete reference list for all citations in Appendix C
% Organized alphabetically by author
% =============================================================================


%------------------------------------------------------------------------------
% BBB PARAMETER CALIBRATION REFERENCE
%------------------------------------------------------------------------------
\begin{tcolorbox}[colback=blue!5!white, colframe=blue!75!black, 
    title=\textbf{Parameter Calibration: See Appendix BBB}]
    
The parameters introduced in this appendix require empirical calibration. 
\textbf{Appendix BBB} (\textit{The Epistemics of Numbers}) provides:

\begin{itemize}[nosep]
    \item The three-tier hybrid estimation strategy (empirical, LLM-MC, axiomatic)
    \item Epistemic status tags $E(\theta)$ for all parameters
    \item The complete $\Gamma$-matrix, $\Lambda$-matrix, and $\delta$-tensor
    \item Worked examples demonstrating the calibration workflow
\end{itemize}

\textbf{Important:} Parameter values marked as \texttt{ILL} (illustrative) in 
Appendix BBB are placeholders requiring local calibration before application.

\end{tcolorbox}




%------------------------------------------------------------------------------
% CENTRAL REFERENCES (Required per Template)
%------------------------------------------------------------------------------

\begin{tcolorbox}[colback=blue!5!white, colframe=blue!75!black, 
    title=\textbf{Central References}]

\textbf{Glossary:} For complete symbol definitions across the EBF framework, 
see \textbf{Appendix G} (\textit{Glossary and Notation}).

\textbf{Bibliography:} For full reference details, see the 
\textbf{Master Bibliography} (\texttt{bcm2\_master\_references.bib}) 
in the repository.

All symbols used in this appendix are consistent with Appendix G definitions.

\end{tcolorbox}

\section{References for Appendix C}
\label{sec:app-c-references}

\begin{thebibliography}{99}

% === A ===
\bibitem{acemoglu2005}
Acemoglu, D., Johnson, S., \& Robinson, J. A. (2005).
Institutions as a fundamental cause of long-run growth.
\textit{Handbook of Economic Growth}, 1, 385--472.

\bibitem{acemoglu2006}
Acemoglu, D., \& Robinson, J. A. (2006).
De facto political power and institutional persistence.
\textit{American Economic Review}, 96(2), 325--330.

\bibitem{acemoglu2012}
Acemoglu, D., \& Robinson, J. A. (2012).
\textit{Why Nations Fail: The Origins of Power, Prosperity, and Poverty}.
Crown Business.

\bibitem{akerlof1980}
Akerlof, G. A. (1980).
A theory of social custom, of which unemployment may be one consequence.
\textit{Quarterly Journal of Economics}, 94(4), 749--775.

\bibitem{akerlof2000}
Akerlof, G. A., \& Kranton, R. E. (2000).
Economics and identity.
\textit{Quarterly Journal of Economics}, 115(3), 715--753.

\bibitem{akerlof2010}
Akerlof, G. A., \& Kranton, R. E. (2010).
\textit{Identity Economics: How Our Identities Shape Our Work, Wages, and Well-Being}.
Princeton University Press.

\bibitem{anderson1983}
Anderson, B. (1983).
\textit{Imagined Communities: Reflections on the Origin and Spread of Nationalism}.
Verso.

\bibitem{antonovsky1979}
Antonovsky, A. (1979).
\textit{Health, Stress, and Coping}.
Jossey-Bass.

% === B ===
\bibitem{banerjee2011}
Banerjee, A. V., \& Duflo, E. (2011).
\textit{Poor Economics: A Radical Rethinking of the Way to Fight Global Poverty}.
PublicAffairs.

\bibitem{bargh2006}
Bargh, J. A. (2006).
What have we been priming all these years? On the development, mechanisms, and ecology of nonconscious social behavior.
\textit{European Journal of Social Psychology}, 36(2), 147--168.

\bibitem{barth1969}
Barth, F. (1969).
\textit{Ethnic Groups and Boundaries: The Social Organization of Culture Difference}.
Little, Brown.

\bibitem{baumeister1995}
Baumeister, R. F., \& Leary, M. R. (1995).
The need to belong: Desire for interpersonal attachments as a fundamental human motivation.
\textit{Psychological Bulletin}, 117(3), 497--529.

\bibitem{baumeister2001}
Baumeister, R. F., Bratslavsky, E., Finkenauer, C., \& Vohs, K. D. (2001).
Bad is stronger than good.
\textit{Review of General Psychology}, 5(4), 323--370.

\bibitem{becker1964}
Becker, G. S. (1964).
\textit{Human Capital: A Theoretical and Empirical Analysis}.
University of Chicago Press.

\bibitem{becker1968}
Becker, G. S. (1968).
Crime and punishment: An economic approach.
\textit{Journal of Political Economy}, 76(2), 169--217.

\bibitem{becker1988}
Becker, G. S., \& Murphy, K. M. (1988).
A theory of rational addiction.
\textit{Journal of Political Economy}, 96(4), 675--700.

\bibitem{belk1988}
Belk, R. W. (1988).
Possessions and the extended self.
\textit{Journal of Consumer Research}, 15(2), 139--168.

\bibitem{belk2013}
Belk, R. W. (2013).
Extended self in a digital world.
\textit{Journal of Consumer Research}, 40(3), 477--500.

\bibitem{benabou2011}
Bénabou, R., \& Tirole, J. (2011).
Laws and norms.
\textit{NBER Working Paper} No. 17579.

\bibitem{benabou2016}
Bénabou, R., \& Tirole, J. (2016).
Mindful economics: The production, consumption, and value of beliefs.
\textit{Journal of Economic Perspectives}, 30(3), 141--164.

\bibitem{benkler2006}
Benkler, Y. (2006).
\textit{The Wealth of Networks: How Social Production Transforms Markets and Freedom}.
Yale University Press.

\bibitem{berkman2000}
Berkman, L. F., \& Glass, T. (2000).
Social integration, social networks, social support, and health.
\textit{Social Epidemiology}, 1, 137--173.

\bibitem{bicchieri2006}
Bicchieri, C. (2006).
\textit{The Grammar of Society: The Nature and Dynamics of Social Norms}.
Cambridge University Press.

\bibitem{bonanno2004}
Bonanno, G. A. (2004).
Loss, trauma, and human resilience: Have we underestimated the human capacity to thrive after extremely aversive events?
\textit{American Psychologist}, 59(1), 20--28.

\bibitem{bourdieu1984}
Bourdieu, P. (1984).
\textit{Distinction: A Social Critique of the Judgement of Taste}.
Harvard University Press.

\bibitem{bourdieu1986}
Bourdieu, P. (1986).
The forms of capital.
In J. Richardson (Ed.), \textit{Handbook of Theory and Research for the Sociology of Education} (pp. 241--258).
Greenwood.

\bibitem{boyd1992}
Boyd, R., \& Richerson, P. J. (1992).
Punishment allows the evolution of cooperation (or anything else) in sizable groups.
\textit{Ethology and Sociobiology}, 13(3), 171--195.

\bibitem{boyd2014}
boyd, d. (2014).
\textit{It's Complicated: The Social Lives of Networked Teens}.
Yale University Press.

\bibitem{braudel1958}
Braudel, F. (1958).
Histoire et sciences sociales: La longue durée.
\textit{Annales. Histoire, Sciences Sociales}, 13(4), 725--753.

\bibitem{bulow1985}
Bulow, J. I., Geanakoplos, J. D., \& Klemperer, P. D. (1985).
Multimarket oligopoly: Strategic substitutes and complements.
\textit{Journal of Political Economy}, 93(3), 488--511.

\bibitem{butler1990}
Butler, J. (1990).
\textit{Gender Trouble: Feminism and the Subversion of Identity}.
Routledge.

% === C ===
\bibitem{carstensen2006}
Carstensen, L. L. (2006).
The influence of a sense of time on human development.
\textit{Science}, 312(5782), 1913--1915.

\bibitem{cialdini1984}
Cialdini, R. B. (1984).
\textit{Influence: The Psychology of Persuasion}.
Harper Business.

\bibitem{clayton2003}
Clayton, S. (2003).
Environmental identity: A conceptual and an operational definition.
\textit{Identity and the Natural Environment}, 45--65.

\bibitem{cole2009}
Cole, E. R. (2009).
Intersectionality and research in psychology.
\textit{American Psychologist}, 64(3), 170--180.

\bibitem{coleman1990}
Coleman, J. S. (1990).
\textit{Foundations of Social Theory}.
Harvard University Press.

\bibitem{collins2004}
Collins, R. (2004).
\textit{Interaction Ritual Chains}.
Princeton University Press.

\bibitem{cooley1902}
Cooley, C. H. (1902).
\textit{Human Nature and the Social Order}.
Scribner's.

\bibitem{crenshaw1989}
Crenshaw, K. (1989).
Demarginalizing the intersection of race and sex.
\textit{University of Chicago Legal Forum}, 1989(1), 139--167.

% === D ===
\bibitem{daly1990}
Daly, H. E. (1990).
Toward some operational principles of sustainable development.
\textit{Ecological Economics}, 2(1), 1--6.

\bibitem{damasio1994}
Damasio, A. R. (1994).
\textit{Descartes' Error: Emotion, Reason, and the Human Brain}.
Putnam.

\bibitem{deaton2008}
Deaton, A. (2008).
Income, health, and well-being around the world: Evidence from the Gallup World Poll.
\textit{Journal of Economic Perspectives}, 22(2), 53--72.

\bibitem{dietz2005}
Dietz, T., Fitzgerald, A., \& Shwom, R. (2005).
Environmental values.
\textit{Annual Review of Environment and Resources}, 30, 335--372.

\bibitem{donahue1993}
Donahue, E. M., Robins, R. W., Roberts, B. W., \& John, O. P. (1993).
The divided self: Concurrent and longitudinal effects of psychological adjustment and social roles on self-concept differentiation.
\textit{Journal of Personality and Social Psychology}, 64(5), 834--846.

\bibitem{durkheim1893}
Durkheim, E. (1893).
\textit{The Division of Labor in Society}.
Free Press (1997 reprint).

\bibitem{durkheim1897}
Durkheim, E. (1897).
\textit{Suicide: A Study in Sociology}.
Free Press (1951 reprint).

\bibitem{dyck2010}
Dyck, A., Morse, A., \& Zingales, L. (2010).
Who blows the whistle on corporate fraud?
\textit{Journal of Finance}, 65(6), 2213--2253.

% === E ===
\bibitem{easterlin1974}
Easterlin, R. A. (1974).
Does economic growth improve the human lot? Some empirical evidence.
\textit{Nations and Households in Economic Growth}, 89, 89--125.

\bibitem{ebaugh1988}
Ebaugh, H. R. F. (1988).
\textit{Becoming an Ex: The Process of Role Exit}.
University of Chicago Press.

\bibitem{edmondson1999}
Edmondson, A. (1999).
Psychological safety and learning behavior in work teams.
\textit{Administrative Science Quarterly}, 44(2), 350--383.

\bibitem{eisenberger2003}
Eisenberger, N. I., Lieberman, M. D., \& Williams, K. D. (2003).
Does rejection hurt? An fMRI study of social exclusion.
\textit{Science}, 302(5643), 290--292.

\bibitem{elster1986}
Elster, J. (Ed.). (1986).
\textit{The Multiple Self}.
Cambridge University Press.

\bibitem{elster1996}
Elster, J. (1996).
Rationality and the emotions.
\textit{Economic Journal}, 106(438), 1386--1397.

\bibitem{erikson1959}
Erikson, E. H. (1959).
\textit{Identity and the Life Cycle}.
International Universities Press.

\bibitem{erikson1968}
Erikson, E. H. (1968).
\textit{Identity: Youth and Crisis}.
Norton.

% === F ===
\bibitem{fehr1999}
Fehr, E., \& Schmidt, K. M. (1999).
A theory of fairness, competition, and cooperation.
\textit{Quarterly Journal of Economics}, 114(3), 817--868.

\bibitem{fehr2002}
Fehr, E., \& Gächter, S. (2002).
Altruistic punishment in humans.
\textit{Nature}, 415(6868), 137--140.

\bibitem{festinger1954}
Festinger, L. (1954).
A theory of social comparison processes.
\textit{Human Relations}, 7(2), 117--140.

\bibitem{festinger1957}
Festinger, L. (1957).
\textit{A Theory of Cognitive Dissonance}.
Stanford University Press.

\bibitem{fielding2008}
Fielding, K. S., McDonald, R., \& Louis, W. R. (2008).
Theory of planned behaviour, identity and intentions to engage in environmental activism.
\textit{Journal of Environmental Psychology}, 28(4), 318--326.

\bibitem{fischbacher2001}
Fischbacher, U., Gächter, S., \& Fehr, E. (2001).
Are people conditionally cooperative? Evidence from a public goods experiment.
\textit{Economics Letters}, 71(3), 397--404.

\bibitem{frederick2002}
Frederick, S., Loewenstein, G., \& O'Donoghue, T. (2002).
Time discounting and time preference: A critical review.
\textit{Journal of Economic Literature}, 40(2), 351--401.

\bibitem{freud1915}
Freud, S. (1915).
The unconscious.
\textit{Standard Edition}, 14, 159--215.

% === G ===
\bibitem{gesi2015}
GeSI (2015).
\textit{\#SMARTer2030: ICT Solutions for 21st Century Challenges}.
Global e-Sustainability Initiative.

\bibitem{glaeser2003}
Glaeser, E., Scheinkman, J., \& Shleifer, A. (2003).
The injustice of inequality.
\textit{Journal of Monetary Economics}, 50(1), 199--222.

\bibitem{gneezy2011}
Gneezy, U., Meier, S., \& Rey-Biel, P. (2011).
When and why incentives (don't) work to modify behavior.
\textit{Journal of Economic Perspectives}, 25(4), 191--210.

\bibitem{goffman1959}
Goffman, E. (1959).
\textit{The Presentation of Self in Everyday Life}.
Doubleday.

\bibitem{goffman1963}
Goffman, E. (1963).
\textit{Stigma: Notes on the Management of Spoiled Identity}.
Prentice-Hall.

\bibitem{granovetter1973}
Granovetter, M. S. (1973).
The strength of weak ties.
\textit{American Journal of Sociology}, 78(6), 1360--1380.

\bibitem{greenwald1995}
Greenwald, A. G., \& Banaji, M. R. (1995).
Implicit social cognition: Attitudes, self-esteem, and stereotypes.
\textit{Psychological Review}, 102(1), 4--27.

\bibitem{greif1993}
Greif, A. (1993).
Contract enforceability and economic institutions in early trade: The Maghribi traders' coalition.
\textit{American Economic Review}, 83(3), 525--548.

% === H ===
\bibitem{haines2006}
Haines, A., Kovats, R. S., Campbell-Lendrum, D., \& Corvalán, C. (2006).
Climate change and human health: Impacts, vulnerability and public health.
\textit{Public Health}, 120(7), 585--596.

\bibitem{hardin1968}
Hardin, G. (1968).
The tragedy of the commons.
\textit{Science}, 162(3859), 1243--1248.

\bibitem{hargittai2002}
Hargittai, E. (2002).
Second-level digital divide: Differences in people's online skills.
\textit{First Monday}, 7(4).

\bibitem{heiner1983}
Heiner, R. A. (1983).
The origin of predictable behavior.
\textit{American Economic Review}, 73(4), 560--595.

\bibitem{helliwell2004}
Helliwell, J. F., \& Putnam, R. D. (2004).
The social context of well-being.
\textit{Philosophical Transactions of the Royal Society B}, 359(1449), 1435--1446.

\bibitem{henrich2006}
Henrich, J., McElreath, R., Barr, A., Ensminger, J., Barrett, C., Bolyanatz, A., ... \& Ziker, J. (2006).
Costly punishment across human societies.
\textit{Science}, 312(5781), 1767--1770.

\bibitem{herrmann2008}
Herrmann, B., Thöni, C., \& Gächter, S. (2008).
Antisocial punishment across societies.
\textit{Science}, 319(5868), 1362--1367.

\bibitem{hilty2015}
Hilty, L. M., \& Aebischer, B. (Eds.). (2015).
\textit{ICT Innovations for Sustainability}.
Springer.

\bibitem{hobbes1651}
Hobbes, T. (1651).
\textit{Leviathan}.
Andrew Crooke.

\bibitem{hofstede2001}
Hofstede, G. (2001).
\textit{Culture's Consequences: Comparing Values, Behaviors, Institutions and Organizations Across Nations} (2nd ed.).
Sage.

\bibitem{hogg2007}
Hogg, M. A. (2007).
Uncertainty–identity theory.
\textit{Advances in Experimental Social Psychology}, 39, 69--126.

\bibitem{holling1973}
Holling, C. S. (1973).
Resilience and stability of ecological systems.
\textit{Annual Review of Ecology and Systematics}, 4(1), 1--23.

\bibitem{honneth1992}
Honneth, A. (1992).
\textit{Kampf um Anerkennung} [The Struggle for Recognition].
Suhrkamp.

% === K ===
\bibitem{kahneman1979}
Kahneman, D., \& Tversky, A. (1979).
Prospect theory: An analysis of decision under risk.
\textit{Econometrica}, 47(2), 263--291.

\bibitem{kahneman2011}
Kahneman, D. (2011).
\textit{Thinking, Fast and Slow}.
Farrar, Straus and Giroux.

\bibitem{kals1999}
Kals, E., Schumacher, D., \& Montada, L. (1999).
Emotional affinity toward nature as a motivational basis to protect nature.
\textit{Environment and Behavior}, 31(2), 178--202.

\bibitem{kawachi1997}
Kawachi, I., Kennedy, B. P., Lochner, K., \& Prothrow-Stith, D. (1997).
Social capital, income inequality, and mortality.
\textit{American Journal of Public Health}, 87(9), 1491--1498.

\bibitem{kegan1982}
Kegan, R. (1982).
\textit{The Evolving Self: Problem and Process in Human Development}.
Harvard University Press.

\bibitem{korsgaard2009}
Korsgaard, C. M. (2009).
\textit{Self-Constitution: Agency, Identity, and Integrity}.
Oxford University Press.

\bibitem{kurzban2005}
Kurzban, R., \& Houser, D. (2005).
Experiments investigating cooperative types in humans: A complement to evolutionary theory and simulations.
\textit{Proceedings of the National Academy of Sciences}, 102(5), 1803--1807.

% === L ===
\bibitem{laibson1997}
Laibson, D. (1997).
Golden eggs and hyperbolic discounting.
\textit{Quarterly Journal of Economics}, 112(2), 443--478.

\bibitem{lea2006}
Lea, S. E., \& Webley, P. (2006).
Money as tool, money as drug: The biological psychology of a strong incentive.
\textit{Behavioral and Brain Sciences}, 29(2), 161--209.

\bibitem{lewin1947}
Lewin, K. (1947).
Frontiers in group dynamics: Concept, method and reality in social science.
\textit{Human Relations}, 1(1), 5--41.

\bibitem{loewenstein1992}
Loewenstein, G., \& Prelec, D. (1992).
Anomalies in intertemporal choice: Evidence and an interpretation.
\textit{Quarterly Journal of Economics}, 107(2), 573--597.

\bibitem{loewenstein1996}
Loewenstein, G. (1996).
Out of control: Visceral influences on behavior.
\textit{Organizational Behavior and Human Decision Processes}, 65(3), 272--292.

\bibitem{longo2019}
Longo, S. B., Clark, B., Shriver, T. E., \& Clausen, R. (2019).
Sustainability and environmental sociology: Putting the economy in its place and moving toward an integrative socio-ecology.
\textit{Sustainability}, 11(18), 4983.

\bibitem{luhmann1979}
Luhmann, N. (1979).
\textit{Trust and Power}.
Wiley.

\bibitem{luhmann1984}
Luhmann, N. (1984).
\textit{Soziale Systeme} [Social Systems].
Suhrkamp.

% === M ===
\bibitem{marcia1966}
Marcia, J. E. (1966).
Development and validation of ego-identity status.
\textit{Journal of Personality and Social Psychology}, 3(5), 551--558.

\bibitem{marmot2004}
Marmot, M. (2004).
\textit{Status Syndrome: How Your Social Standing Directly Affects Your Health}.
Bloomsbury.

\bibitem{maslow1943}
Maslow, A. H. (1943).
A theory of human motivation.
\textit{Psychological Review}, 50(4), 370--396.

\bibitem{maynardsmith1982}
Maynard Smith, J. (1982).
\textit{Evolution and the Theory of Games}.
Cambridge University Press.

\bibitem{mead1934}
Mead, G. H. (1934).
\textit{Mind, Self, and Society}.
University of Chicago Press.

\bibitem{milgrom1990}
Milgrom, P., \& Roberts, J. (1990).
The economics of modern manufacturing: Technology, strategy, and organization.
\textit{American Economic Review}, 80(3), 511--528.

\bibitem{milgrom1995}
Milgrom, P., \& Roberts, J. (1995).
Complementarities and fit: Strategy, structure, and organizational change in manufacturing.
\textit{Journal of Accounting and Economics}, 19(2-3), 179--208.

\bibitem{mirrlees1971}
Mirrlees, J. A. (1971).
An exploration in the theory of optimum income taxation.
\textit{Review of Economic Studies}, 38(2), 175--208.

% === N ===
\bibitem{nash1950}
Nash, J. F. (1950).
Equilibrium points in n-person games.
\textit{Proceedings of the National Academy of Sciences}, 36(1), 48--49.

\bibitem{nissenbaum2009}
Nissenbaum, H. (2009).
\textit{Privacy in Context: Technology, Policy, and the Integrity of Social Life}.
Stanford University Press.

\bibitem{nordhaus1994}
Nordhaus, W. D. (1994).
\textit{Managing the Global Commons: The Economics of Climate Change}.
MIT Press.

\bibitem{north1990}
North, D. C. (1990).
\textit{Institutions, Institutional Change and Economic Performance}.
Cambridge University Press.

% === O ===
\bibitem{odonoghue1999}
O'Donoghue, T., \& Rabin, M. (1999).
Doing it now or later.
\textit{American Economic Review}, 89(1), 103--124.

\bibitem{olson1965}
Olson, M. (1965).
\textit{The Logic of Collective Action: Public Goods and the Theory of Groups}.
Harvard University Press.

\bibitem{ostrom1990}
Ostrom, E. (1990).
\textit{Governing the Commons: The Evolution of Institutions for Collective Action}.
Cambridge University Press.

\bibitem{ostrom2005}
Ostrom, E. (2005).
\textit{Understanding Institutional Diversity}.
Princeton University Press.

% === P ===
\bibitem{parsons1951}
Parsons, T. (1951).
\textit{The Social System}.
Free Press.

\bibitem{paulhus2002}
Paulhus, D. L., \& Williams, K. M. (2002).
The dark triad of personality: Narcissism, Machiavellianism, and psychopathy.
\textit{Journal of Research in Personality}, 36(6), 556--563.

\bibitem{pigou1920}
Pigou, A. C. (1920).
\textit{The Economics of Welfare}.
Macmillan.

\bibitem{piketty2014}
Piketty, T. (2014).
\textit{Capital in the Twenty-First Century}.
Harvard University Press.

\bibitem{pindyck2007}
Pindyck, R. S. (2007).
Uncertainty in environmental economics.
\textit{Review of Environmental Economics and Policy}, 1(1), 45--65.

\bibitem{polinsky2000}
Polinsky, A. M., \& Shavell, S. (2000).
The economic theory of public enforcement of law.
\textit{Journal of Economic Literature}, 38(1), 45--76.

\bibitem{putnam2000}
Putnam, R. D. (2000).
\textit{Bowling Alone: The Collapse and Revival of American Community}.
Simon \& Schuster.

% === R ===
\bibitem{rawls1971}
Rawls, J. (1971).
\textit{A Theory of Justice}.
Harvard University Press.

\bibitem{rheingold2000}
Rheingold, H. (2000).
\textit{The Virtual Community: Homesteading on the Electronic Frontier} (2nd ed.).
MIT Press.

\bibitem{ricoeur1990}
Ricoeur, P. (1990).
\textit{Oneself as Another}.
University of Chicago Press.

\bibitem{robertson1995}
Robertson, R. (1995).
Glocalization: Time-space and homogeneity-heterogeneity.
\textit{Global Modernities}, 2(1), 25--44.

\bibitem{roccas2002}
Roccas, S., \& Brewer, M. B. (2002).
Social identity complexity.
\textit{Personality and Social Psychology Review}, 6(2), 88--106.

% === S ===
\bibitem{samuelson1937}
Samuelson, P. A. (1937).
A note on measurement of utility.
\textit{Review of Economic Studies}, 4(2), 155--161.

\bibitem{samuelson1954}
Samuelson, P. A. (1954).
The pure theory of public expenditure.
\textit{Review of Economics and Statistics}, 36(4), 387--389.

\bibitem{sandel2012}
Sandel, M. J. (2012).
\textit{What Money Can't Buy: The Moral Limits of Markets}.
Farrar, Straus and Giroux.

\bibitem{scheffer2009}
Scheffer, M., Bascompte, J., Brock, W. A., Brovkin, V., Carpenter, S. R., Dakos, V., ... \& Sugihara, G. (2009).
Early-warning signals for critical transitions.
\textit{Nature}, 461(7260), 53--59.

\bibitem{schelling1978}
Schelling, T. C. (1978).
\textit{Micromotives and Macrobehavior}.
Norton.

\bibitem{schultz2002}
Schultz, P. W. (2002).
Inclusion with nature: The psychology of human-nature relations.
In P. Schmuck \& W. P. Schultz (Eds.), \textit{Psychology of Sustainable Development} (pp. 61--78).
Springer.

\bibitem{sen1999}
Sen, A. (1999).
\textit{Development as Freedom}.
Knopf.

\bibitem{sen2006}
Sen, A. (2006).
\textit{Identity and Violence: The Illusion of Destiny}.
Norton.

\bibitem{simon1957}
Simon, H. A. (1957).
\textit{Models of Man: Social and Rational}.
Wiley.

\bibitem{simon1996}
Simon, G. K. (1996).
\textit{In Sheep's Clothing: Understanding and Dealing with Manipulative People}.
A. J. Christopher.

\bibitem{spence1973}
Spence, M. (1973).
Job market signaling.
\textit{Quarterly Journal of Economics}, 87(3), 355--374.

\bibitem{stern1999}
Stern, P. C., Dietz, T., Abel, T., Guagnano, G. A., \& Kalof, L. (1999).
A value-belief-norm theory of support for social movements: The case of environmentalism.
\textit{Human Ecology Review}, 6(2), 81--97.

\bibitem{stern2007}
Stern, N. (2007).
\textit{The Economics of Climate Change: The Stern Review}.
Cambridge University Press.

\bibitem{stiglitz1999}
Stiglitz, J. E. (1999).
Knowledge as a global public good.
\textit{Global Public Goods}, 1(9), 308--325.

\bibitem{stiglitz2002}
Stiglitz, J. E. (2002).
Transparency in government.
In \textit{The Right to Tell: The Role of Mass Media in Economic Development} (pp. 27--44).
World Bank.

\bibitem{stryker1980}
Stryker, S. (1980).
\textit{Symbolic Interactionism: A Social Structural Version}.
Benjamin/Cummings.

\bibitem{sunstein1996}
Sunstein, C. R. (1996).
Social norms and social roles.
\textit{Columbia Law Review}, 96(4), 903--968.

\bibitem{sunstein2002}
Sunstein, C. R. (2002).
The law of group polarization.
\textit{Journal of Political Philosophy}, 10(2), 175--195.

\bibitem{sydow2009}
Sydow, J., Schreyögg, G., \& Koch, J. (2009).
Organizational path dependence: Opening the black box.
\textit{Academy of Management Review}, 34(4), 689--709.

% === T ===
\bibitem{tajfel1974}
Tajfel, H. (1974).
Social identity and intergroup behaviour.
\textit{Social Science Information}, 13(2), 65--93.

\bibitem{tajfel1979}
Tajfel, H., \& Turner, J. C. (1979).
An integrative theory of intergroup conflict.
In W. G. Austin \& S. Worchel (Eds.), \textit{The Social Psychology of Intergroup Relations} (pp. 33--47).
Brooks/Cole.

\bibitem{taleb2007}
Taleb, N. N. (2007).
\textit{The Black Swan: The Impact of the Highly Improbable}.
Random House.

\bibitem{taleb2012}
Taleb, N. N. (2012).
\textit{Antifragile: Things That Gain from Disorder}.
Random House.

\bibitem{thaler1999}
Thaler, R. H. (1999).
Mental accounting matters.
\textit{Journal of Behavioral Decision Making}, 12(3), 183--206.

\bibitem{PAP-thaler2008}
Thaler, R. H., \& Sunstein, C. R. (2008).
\textit{Nudge: Improving Decisions About Health, Wealth, and Happiness}.
Yale University Press.

\bibitem{topkis1998}
Topkis, D. M. (1998).
\textit{Supermodularity and Complementarity}.
Princeton University Press.

\bibitem{topol2019}
Topol, E. J. (2019).
\textit{Deep Medicine: How Artificial Intelligence Can Make Healthcare Human Again}.
Basic Books.

\bibitem{turkle2011}
Turkle, S. (2011).
\textit{Alone Together: Why We Expect More from Technology and Less from Each Other}.
Basic Books.

\bibitem{turner1987}
Turner, J. C., Hogg, M. A., Oakes, P. J., Reicher, S. D., \& Wetherell, M. S. (1987).
\textit{Rediscovering the Social Group: A Self-Categorization Theory}.
Blackwell.

\bibitem{tyler1990}
Tyler, T. R. (1990).
\textit{Why People Obey the Law}.
Yale University Press.

\bibitem{tyler2006}
Tyler, T. R. (2006).
Psychological perspectives on legitimacy and legitimation.
\textit{Annual Review of Psychology}, 57, 375--400.

% === U ===
\bibitem{uslaner2002}
Uslaner, E. M. (2002).
\textit{The Moral Foundations of Trust}.
Cambridge University Press.

% === V ===
\bibitem{vandijk2020}
Van Dijk, J. (2020).
\textit{The Digital Divide}.
Polity Press.

\bibitem{veblen1899}
Veblen, T. (1899).
\textit{The Theory of the Leisure Class}.
Macmillan.

% === W ===
\bibitem{walther1996}
Walther, J. B. (1996).
Computer-mediated communication: Impersonal, interpersonal, and hyperpersonal interaction.
\textit{Communication Research}, 23(1), 3--43.

\bibitem{walzer1983}
Walzer, M. (1983).
\textit{Spheres of Justice: A Defense of Pluralism and Equality}.
Basic Books.

\bibitem{weitzman2009}
Weitzman, M. L. (2009).
On modeling and interpreting the economics of catastrophic climate change.
\textit{Review of Economics and Statistics}, 91(1), 1--19.

\bibitem{wellman2001}
Wellman, B. (2001).
Physical place and cyberplace: The rise of personalized networking.
\textit{International Journal of Urban and Regional Research}, 25(2), 227--252.

\bibitem{wilkinson2009}
Wilkinson, R. G., \& Pickett, K. (2009).
\textit{The Spirit Level: Why Greater Equality Makes Societies Stronger}.
Bloomsbury Press.

\bibitem{williams2007}
Williams, K. D. (2007).
Ostracism.
\textit{Annual Review of Psychology}, 58, 425--452.

\bibitem{wilson1984}
Wilson, E. O. (1984).
\textit{Biophilia}.
Harvard University Press.

% === Y ===
\bibitem{yee2007}
Yee, N., \& Bailenson, J. (2007).
The Proteus effect: The effect of transformed self-representation on behavior.
\textit{Human Communication Research}, 33(3), 271--290.

% === Z ===
\bibitem{zuboff2019}
Zuboff, S. (2019).
\textit{The Age of Surveillance Capitalism: The Fight for a Human Future at the New Frontier of Power}.
PublicAffairs.

\end{thebibliography}

%%%%%%%%%%%%%%%%%%%%%%%%%%%%%%%%%%%%%%%%%%%%%%%%%%%%%%%%%%%%%%%%%%%%%%%%%%%%%%%
% END OF APPENDIX C
%%%%%%%%%%%%%%%%%%%%%%%%%%%%%%%%%%%%%%%%%%%%%%%%%%%%%%%%%%%%%%%%%%%%%%%%%%%%%%%


%==============================================================================
% CROSS-REFERENCES TO OTHER CORE APPENDICES (for C)
%==============================================================================

\subsection{Integration with CORE Framework}
\label{sec:C-core-integration}

This appendix answers the CORE question ``WHAT is utility?'' and integrates 
with the other three CORE appendices:

\subsubsection{Connection to Appendix AAA (WHO - Aggregation)}

The six FEPSDE dimensions apply at all aggregation levels $L$, but with 
level-specific interpretations:

\begin{table}[h]
\centering
\small
\begin{tabular}{l|ll}
\toprule
\textbf{Dimension} & \textbf{Individual ($L=1$)} & \textbf{National ($L=5$)} \\
\midrule
Financial (F) & Personal income/wealth & GDP, fiscal balance \\
Emotional (E) & Subjective wellbeing & National happiness index \\
Physical (P) & Personal health & Public health outcomes \\
Social (S) & Personal relationships & Social cohesion, trust \\
Digital (D) & Tech access/literacy & Digital infrastructure \\
Ecological (E) & Personal footprint & Environmental quality \\
\bottomrule
\end{tabular}
\caption{FEPSDE interpretation by aggregation level}
\end{table}

See Appendix AAA for the complete aggregation framework.

\subsubsection{Connection to Appendix B (HOW - Complementarity)}

The FEPSDE dimensions define the rows and columns of the complementarity 
matrix $\Gamma$:

\begin{equation}
\Gamma = [\gamma_{dd'}]_{d, d' \in \{F, E, P, S, D, E\}}
\end{equation}

The sign and magnitude of $\gamma_{dd'}$ determines whether dimensions 
reinforce (complements) or offset (substitutes) each other.

See Appendix B for the complete complementarity theory including:
\begin{itemize}
    \item Teil II: Utility $\times$ Utility interactions ($\Gamma$-matrix)
    \item Teil III: Context modulation $\gamma_{dd'}(\Psi)$
\end{itemize}

\subsubsection{Connection to Appendix V (WHEN - Context)}

Context $\Psi$ affects both the weights and the values of FEPSDE dimensions:

\begin{equation}
U = \sum_{d \in \text{FEPSDE}} w_d(\Psi) \cdot u_d(a; \Psi)
\end{equation}

Key insight: The same action $a$ produces different utility profiles 
$(u_F, u_E, u_P, u_S, u_D, u_E)$ depending on context $\Psi$.

See Appendix V for the 8$\Psi$-Dimensions and their empirical validation.

\vspace{0.5em}
\hrule
\vspace{0.5em}
\noindent\textit{Cross-references: Appendix AAA (Aggregation), 
Appendix B (Complementarity), Appendix V (8$\Psi$-Dimensions), 
Appendix BBB (Parameter Calibration)}

